\documentclass[11pt,a4paper]{article}
\usepackage[T1]{fontenc}
\usepackage{lmodern}
\usepackage[margin=27mm]{geometry}
\usepackage{microtype}
\usepackage{amsmath,amssymb,amsthm,mathtools}
\usepackage{booktabs}
\usepackage[hidelinks]{hyperref}
\usepackage[nameinlink,noabbrev]{cleveref}

\hypersetup{pdftitle={Local reductions and certified norm bounds for the Li--Liu spectral hypothesis S}}
\allowdisplaybreaks[2]
\setlength{\emergencystretch}{2em}
\numberwithin{equation}{section}
\newtheorem{theorem}{Theorem}[section]
\newcommand{\SSS}{\mathfrak S}
\newcommand{\grh}{\mathrm{GRH}}
\newcommand{\ud}{\mathrm d}
\newcommand{\E}{\mathcal E_N}
\newcommand{\B}{\mathcal B_N}
\newcommand{\NN}{\mathcal N_N}
\newcommand{\NNzero}{\mathcal N_{N,0}}
\newcommand{\Tbal}{T_N^{\mathrm{bal}}}
\newcommand{\Trect}{T_N^{X,Y}}
\newcommand{\dd}{\operatorname{dist}}
\newcommand{\tr}{\operatorname{tr}}

\title{\bfseries Local reductions and certified norm bounds\\
for the Li--Liu spectral hypothesis S}
\author{}
\date{26 November 2026}

\begin{document}
\maketitle

\begin{abstract}
Li and Liu's spectral hypothesis S asks for an asymptotic bound on a
Hardy-space operator associated with the binary Goldbach problem.
Under the modulus-specific generalized Riemann hypothesis, we prove
that the norm of its middle-third restriction vanishes on the
singular-series scale precisely when the normalized weighted
Goldbach measure converges to interval length uniformly over
subintervals of every fixed interior window.
For the full norm, energy-orthogonal interpolation gives an explicit
finite matrix. The normalized operator has trace norm
$O((\log N)^2)$, which permits a further spectral compression to
dimension $O((\log N)^2\log\log N)$ with uniform vanishing error.
We evaluate the deterministic contribution to its signed trace
and deduce a necessary local mean-square law for $L'/L$.
Hypothesis S also forces sine-kernel pair correlation for the
odd-character zeros at a fixed modulus. Under GRH alone the
pair law holds for tests with Fourier support in $(-1,1)$.
An equivalent full-norm criterion uses zero ordinates of
polylogarithmic height and prime-power data of polylogarithmic size.
A sieve controls pairs far below $xy=N^2$ and in thin bands
around sharp product cutoffs. Through $xy=CN^2$, norm decay
is equivalent to decay of signed rectangle discrepancies.
Under modulus-specific GRH a smooth upper product cutoff at
$N^{2\log\log N+\kappa}$ costs $O(e^{-\kappa})$.
A two-sided norm comparison lowers the product threshold to
$N^{(3/2)\log\log N+\kappa}$ when the larger input exceeds
$N^{\log\log N+\kappa}$, with smooth transitions.
For $x,y\ge N^{1+\delta}$, the signed form equals a covariance
of Dirichlet zeros minus the mean archimedean density, up to a
power-saving operator error.
At a fixed modulus, sufficiently high pairs of supports have
vanishing mean norm over a common logarithmic translation.
For each fixed $A$, the restriction to $x,y\le N(\log N)^A$
has rapidly vanishing mean-square norm over the moduli. Global
Dirichlet GRH extends this to $x,y\le N^{1+\eta}$ for every
fixed $0<\eta<1/3$, with a power saving. For inputs above
$N^\delta$, the power saving holds through the exact product
cutoff $xy\le N^{2+\gamma}$ for $0<\gamma<\delta/2$ and
every fixed $0<\delta<1$.
For moduli with two suitably sized prime factors, a
product cutoff at $N^2(\log N)^B$ gives mean-square decay
with arbitrary test supports, for every fixed $B\ge0$.
At each modulus, all but $O(\log N\log\log N)$ proper prime
powers can be removed from the arithmetic form while preserving
the full-norm smallness criterion.
The remaining progression terms vanish through $N^2(\log N)^B$
for fixed inputs at almost all moduli. They also vanish for
polylogarithmic inputs on a family of positive lower density.
We certify two-sided norm bounds at five balanced moduli and for
three separated rectangular restrictions. A rational test gives
$\|T_{32762}\|>0.2485744893$ in the original normalization.
The asymptotic estimate for the full operator remains open.
\end{abstract}

\section{Introduction}

Li and Liu~\cite{LiLiu} express a weighted Goldbach correlation using odd
Dirichlet characters modulo the even integer being represented. After they
subtract the diagonal and a continuous model, the remainder defines a compact
self-adjoint operator $T_N$ under their modulus-specific generalized Riemann
hypothesis, denoted here by $\grh(N)$. Their spectral hypothesis S controls
the ensuing operator norm. We study the eventual smallness condition
\begin{equation}
  \frac{\|T_N\|}{\SSS(N)}\longrightarrow 0
  \qquad (N\to\infty,\;N\text{ even}),
  \label{eq:goal}
\end{equation}
which is the norm formulation of S when the admissible auxiliary function
may be chosen to grow arbitrarily slowly. Here $\SSS(N)$ is the ordinary
two-prime singular series.

The norm in \cref{eq:goal} is a supremum over test functions whose supports
may move with $N$. When both supports lie in $(N/3,2N/3)$, the only
off-diagonal congruence contribution is the reflection $m+n=N$. This
restriction makes it possible to compare operator norm with subinterval
counting discrepancy. It also reduces the computation to an interval
problem on a line. The full norm can sample other regions, including a
pair of windows with respective scales $\sqrt N$ and $N^{3/2}$.

\Cref{thm:disc} gives the equivalence for the balanced window.
\Cref{thm:cutoff} proves an asymptotic truncation of the full norm, including
its mixed terms between retained and discarded supports.
\Cref{thm:mesh} reduces the truncated problem to a matrix on a logarithmic
grid with an explicit error bound. \Cref{thm:tracecompression}
compresses this matrix by selecting spectral directions of the separate
arithmetic and model forms. \Cref{thm:signedtrace} expresses the
signed trace in terms of a weighted logarithmic-derivative mean.
\Cref{thm:localderivative,thm:necessarypair} derive a local
mean-square law and a pair-correlation consequence of hypothesis S.
\Cref{thm:restrictedzeropairs} proves a mixed zero-covariance law
through the critical product from $\grh(N)$.
\Cref{thm:fullzerodata} gives an equivalent finite matrix test
in terms of low zero ordinates and short prime-power sums.
\Cref{thm:subcritical} excludes the
region below a product threshold in its entries, and \cref{thm:zeros}
identifies a Dirichlet-zero covariance when both variables are large.
\Cref{thm:criticaldiscrepancy} characterizes the sharply
truncated norm through $xy=CN^2$ by signed rectangle counts.
\Cref{thm:criticalgrid} reduces those counts to prefix rectangles
on $O(\log N)$ dyadic pairs at fixed precision.
\Cref{thm:shiftmean} bounds the mean norm of translated pairs of
high-support windows at a single modulus.
\Cref{thm:upperproduct} gives a uniform upper cutoff on the product
and isolates the character term at large products.
\Cref{thm:asymupperproduct} tightens the cutoff for separated
supports by using positivity on the smaller input.
\Cref{thm:almostall} bounds the entire supported norm on a growing
window for almost all even moduli.
\Cref{thm:mixedpower} bounds mixed windows with upper scales
$N^a,N^b$ for $0<a\le b$ and $b+a/2<2$.
\Cref{thm:productwindow,thm:powerproduct,thm:productfactored} extend the estimates
to product ranges, including the smallest inputs
on a family of positive lower density.
\Cref{thm:ppcut} reduces the full problem to a prime-only form
and a finite set of opposite-residue progression sums.
\Cref{thm:criticalprofiles} bounds those sums near the square
of the modulus using equidistribution in progressions with
factored moduli.
We also give two-sided numerical certificates for the balanced restriction
and three rectangular restrictions, and a lower bound on the unrestricted
norm at a single modulus.

\section{The form in physical coordinates}
\label{sec:form}

Let $N>e^{10}$ be even, put $\alpha=1/\log N$, and write $\Lambda$ for the
von Mangoldt function. For a positive even integer $h$, define
\begin{equation}
  \SSS(h)=2\prod_{p>2}\left(1-\frac{1}{(p-1)^2}\right)
     \prod_{\substack{p\mid h\\p>2}}\frac{p-1}{p-2}.
  \label{eq:series}
\end{equation}
For $f,g\in C_c^\infty(0,\infty)$, the variables of $f$ and $g$ are scaled
by $N$. The Hardy vectors of Li and Liu~\cite[Proposition~7.4]{LiLiu} have
squared norm
\begin{equation}
  \E(f)=\|v_N(f)\|^2
    =\frac{1}{2\pi}\int_0^\infty
      \frac{|f(u)/2-uf'(u)|^2}{(Nu)^{-2\alpha}+(Nu)^{-1/4}}\,\ud u.
  \label{eq:energy}
\end{equation}
In particular, on any fixed compact interval of $(0,\infty)$,
$\E(f)^{1/2}$ is uniformly equivalent to the $H^1_0$ norm of a supported
test. To see this, write
\begin{equation}
\begin{split}
  a(u)&=\bigl[2\pi((Nu)^{-2\alpha}+(Nu)^{-1/4})\bigr]^{-1},
  \qquad q(u)=\frac{u a'(u)}{a(u)},\\
  p(u)&=u^2a(u),\qquad d(u)=a(u)\left(\frac34+\frac{q(u)}2\right).
\end{split}
  \label{eq:coeff}
\end{equation}
Integration by parts for zero boundary values gives
\begin{equation}
  \E(f)=\int_0^\infty \bigl[p(u)|f'(u)|^2+d(u)|f(u)|^2\bigr]\,\ud u.
  \label{eq:sturm}
\end{equation}
The function $q$ is a weighted mean of $2\alpha$ and $1/4$. Its derivative
with respect to $\log u$ is the negative weighted variance of these two
numbers. Thus $a$ and $p$ are increasing, and $q$ is decreasing.

We use the convention $\B(f,g)=\langle T_N v_N(f),v_N(g)\rangle$,
with the inner product linear in its first argument. In the following
identity a test is extended by zero outside $(0,\infty)$:
\begin{align}
 \B(f,g)
  &=\frac1N
    \sum_{\substack{m,n\ge1,\ (mn,N)=1\\m\ne n}}
      \Lambda(m)\Lambda(n)f(m/N)\overline{g(n/N)}
      \bigl(\mathbf 1_{m\equiv n\ (N)}-\mathbf 1_{m\equiv -n\ (N)}\bigr)
    \notag\\
  &\quad-\sum_{k\ge1}\SSS(kN)J_{f,g}(k),
  \label{eq:form}\\
 J_{f,g}(y)
  &=\int_0^\infty \bigl[
      f(u)\overline{g(u+y)}
      +f(u+y)\overline{g(u)}
      -f(u)\overline{g(y-u)}
      \bigr]\,\ud u.
  \label{eq:model}
\end{align}
For completeness, the character part follows by shifting the contours for
the two logarithmic derivatives to absolute convergence and applying
Mellin inversion. The odd-character average distinguishes equal from
opposite unit residues, and the exact diagonal is removed. For the model,
the three beta integrals of Li and Liu~\cite[Lemma~7.2]{LiLiu} give the
Mellin transform of $J_{f,g}$. The first two terms of \cref{eq:model}
partition the double integral of $f(u)\overline{g(v)}$ when integrated
over $y>0$. The third term has the same double integral. Consequently
$\int_0^\infty J_{f,g}(y)\,\ud y=0$, cancelling the pole of
$\sum_{k\ge1}\SSS(kN)k^{-s}$ at $s=1$ before shifting the model contour.
This gives the last sum in \cref{eq:form}, which is finite for compact
supports. Compare the two-shift evaluation in~\cite[Lemma~2.1]{YLiu}.
Here $\grh(N)$ is the hypothesis used in the contour shift for the
character part.

On bounded support intervals in $(0,\infty)$, \cref{eq:form,eq:energy}
extend to $H^1_0$ by density and continuity of point evaluation. A
restriction norm is the supremum of $|\B(f,g)|/(\E(f)\E(g))^{1/2}$
over nonzero tests in its two support spaces and is bounded above by
the full norm $\|T_N\|$.

\section{The balanced restriction and subinterval discrepancy}
\label{sec:balanced}

Set $I=(1/3,2/3)$. The balanced restriction $\Tbal$ uses $f,g\in H^1_0(I)$.
Here $|u-v|<1$ and $u+v\in(2/3,4/3)$, so the only term of
\cref{eq:form} is the signed $m+n=N$ term and its continuous model. Define the
positive measure $\mu_N$ on $I$ and its signed discrepancy measure by
\begin{equation}
  \mu_N=\frac1{N\SSS(N)}
    \sum_{\substack{n/N\in I\\(n,N)=1}}
       \Lambda(n)\Lambda(N-n)\,\delta_{n/N},
  \qquad
  \nu_N=\ud u-\mu_N.
  \label{eq:nu}
\end{equation}
Then
\begin{equation}
  \frac{\B(f,g)}{\SSS(N)}
    =\int_I f(u)\overline{g(1-u)}\,\ud\nu_N(u).
  \label{eq:balance}
\end{equation}
For $0<\eta<1/6$, put $I_\eta=[1/3+\eta,\,2/3-\eta]$ and define
\begin{equation}
  E_N=\frac{\|\Tbal\|}{\SSS(N)},\qquad
  D_N(\eta)=\sup_{J\subset I_\eta}
     \bigl|\mu_N(J)-|J|\bigr|,
  \label{eq:discrepancy}
\end{equation}
where the supremum allows any interval $J$, with arbitrary choices
of endpoint inclusion.

\begin{theorem}[Balanced norm and counting discrepancy]
\label{thm:disc}
For even $N>e^{10}$ satisfying $\grh(N)$,
\begin{equation}
  D_N(\eta)\ll_\eta E_N^{2/3},\qquad
  E_N\ll D_N(\eta)+\eta .
  \label{eq:equiv}
\end{equation}
The implied constant in the second bound is absolute. Along any sequence
of such moduli tending to infinity, $E_N\to0$ if and only if
$D_N(\eta)\to0$ for every fixed $\eta\in(0,1/6)$.
\end{theorem}

\begin{proof}
First, $\mu_N(I)\ll1$ uniformly. The standard sieve upper bound for the
number of Goldbach representations is $\ll N\SSS(N)/(\log N)^2$
\cite[Theorem~1.2]{Lichtman}; each ordered prime pair contributes at most
$(\log N)^2$ before normalization. Pairs with at least one proper prime
power contribute $O(\sqrt N(\log N)^3)$ before normalization, using the
elementary bound $O(\sqrt N\log N)$ for the number of proper prime
powers up to $N$. Since $\SSS(N)\gg1$, this contribution is bounded as
well. The unit-residue restriction only decreases the count.

Fix $\eta$ and a smooth $g$ with $g(1-u)=1$ on $I_{\eta/2}$ and support
in $I$. Choose smooth upper and lower approximations $f_+,f_-$ to
$\mathbf 1_J$ with transition width $\delta<\eta/4$. They may be taken
between zero and one, supported in $I_{\eta/2}$, and satisfying
$\int_I |f_\pm-\mathbf 1_J|\,\ud u\ll\delta$ and
$\E(f_\pm)^{1/2}\ll_\eta\delta^{-1/2}$. When $J$ is shorter than
$2\delta$, one can take $f_-=0$. Positivity of $\mu_N$ and
\cref{eq:balance} now give
\[
  \bigl|\mu_N(J)-|J|\bigr|\ll_\eta \delta+E_N\delta^{-1/2}.
\]
Take $\delta=E_N^{2/3}$ when this is smaller than $\eta/4$. The
total-mass bound handles larger $E_N$, and $\delta\downarrow0$ handles
$E_N=0$.

For the converse, normalize $\E(f)=\E(g)=1$ and set
$P(u)=f(u)\overline{g(1-u)}$. By \cref{eq:sturm},
$\|P\|_\infty+\int_I|P'|\ll1$. Also
$|P(u)|\ll\dd(u,\partial I)$ by Cauchy--Schwarz from the common
zero boundary values. The integral of $P$ against $\nu_N$ on
$I\setminus I_\eta$ is therefore $O(\eta)$. On $I_\eta=[c,d]$,
integrate by parts with the cumulative signed mass
$F(t)=\nu_N([c,t])$. The endpoint term and the integral of $P'F$
are $O(D_N(\eta))$. Taking the supremum of \cref{eq:balance} proves
the second inequality. For the equivalence of limits, take $N\to\infty$
at fixed $\eta$ and then send $\eta\downarrow0$.
\end{proof}

For example, on any fixed nondegenerate $J\subset I_\eta$, the
condition $E_N\to0$ implies the ordinary ordered prime-pair
asymptotic
\[
  \#\{(p,q): p,q\ \text{prime},\ p+q=N,\ p/N\in J\}
    =\bigl(|J|+o(1)\bigr)\frac{N\SSS(N)}{(\log N)^2}.
\]
Indeed, on this support $\log p,\log q=\log N+O(1)$, and proper
prime powers contribute $o(N\SSS(N))$ to the weighted sum. The only
prime pairs removed by the unit condition have $p=q=N/2$.

\section{Certification on the balanced window}
\label{sec:certificate}

\subsection{A finite comparison for the full Sobolev norm}

Put $h(u)=g(1-u)$. The bound $E_N\le\lambda$ is equivalent to
nonnegativity, for all real $f,h\in H^1_0(I)$, of
\begin{equation}
 Q_\lambda(f,h)=\lambda\bigl(\E(f)+\E(h(1-\cdot))\bigr)
       +2\int_I fh\,\ud\nu_N .
 \label{eq:Q}
\end{equation}
This follows by testing both signs of $h$ and rescaling $f$ and $h$
independently. For complex functions the corresponding Hermitian
form is the sum of its real and imaginary component forms.

Partition $I$ using every atom of $\nu_N$ and a uniform auxiliary
mesh. On a cell of length $\ell$, bound the coefficient pairs for
$f(u)$ and $h(u)$ in \cref{eq:sturm} using endpoint values of
$p(u),a(u),q(u)$ and their reflections. For instance, over a cell
$[b,c]$, the lower and upper mass coefficients for $f$ are
$a(b)(3/4+q(c)/2)$ and $a(c)(3/4+q(b)/2)$. Let $P_-$ be the diagonal
matrix of lower derivative coefficients and let
\[
 M_-=\begin{pmatrix}\lambda d_{f,-}&1\\1&\lambda d_{h,-}\end{pmatrix},
 \qquad
 p_*=\min\operatorname{diag}P_-,
 \qquad
 d_*=\max(d_{f,-},d_{h,-}).
\]
Decompose the vector of cell functions into its linear interpolant
$v$ and a remainder $r$ that vanishes at the endpoints. The derivative
cross term for $P_-$ vanishes, because $v'$ is constant. Also
$M_-\succeq-\mathrm{Id}$ and $\|M_-\|\le1+\lambda d_*$. Poincar\'e's
inequality with $\pi^2\ge9$ and completion of the square give
\begin{align}
 &\int_{\text{cell}}
       \bigl[\lambda (v'+r')^\top P_-(v'+r')
             +(v+r)^\top M_-(v+r)\bigr]\,\ud u
     \notag\\
 &\hspace{2em}\ge
     \int_{\text{cell}}\bigl[\lambda v'^\top P_-v'
                 +v^\top(M_--\Delta_\ell\mathrm{Id})v\bigr]\,\ud u,
     \label{eq:cell}\\
 &\hspace{2em}
    \Delta_\ell=
     \frac{\ell^2(1+\lambda d_*)^2}{9\lambda p_*-\ell^2}
     \quad\text{if }9\lambda p_*>\ell^2 .
     \notag
\end{align}
The atoms are already evaluated exactly at the interpolant nodes.
After assembly, positive definiteness of the lower block
tridiagonal matrix proves $E_N\le\lambda$. Conversely, on a fixed
piecewise linear test, using the upper coefficient bounds gives a
finite quadratic form that bounds $Q_\lambda$ from above. A certified
strictly negative value proves $E_N>\lambda$.

\subsection{Interval implementation and results}

The implementation represents the atom positions and auxiliary nodes
as rational numbers. An integer sieve identifies prime powers, and
each weight uses the logarithm of its prime base. The product in
\cref{eq:series} is enclosed by rounding the finite product outward
and bounding the omitted product between $1-(P-1)^{-1}$ and $1$
at prime cutoff $P$. This follows from
$\sum_{j\ge P}j^{-2}\le (P-1)^{-1}$. The balanced computation uses
$P=10^6$, 2,048 uniform auxiliary cells and 45 decimal digits of
interval working precision. Floating point selects proposed
thresholds and trial vectors. The upper bound is then checked by
interval block $LDL^\top$ pivots from \cref{eq:cell}, and the lower
bound by direct interval evaluation on the fixed trial vector.

\begin{table}[htbp]
\centering
\caption{Two-sided enclosures for the normalized balanced norm $E_N$,
conditional on $\grh(N)$. Endpoints are rounded outward. Atom counts
refer to the ordered, von Mangoldt-weighted sum in \cref{eq:nu}.}
\label{tab:balanced}
\small
\begin{tabular}{@{}r r r l@{}}
\toprule
$N$ & Atoms & Mesh cells & Interval for $E_N$\\
\midrule
$32\,762$     & $172$    & $2\,220$  & $[2.96102,\,2.96945]\times10^{-3}$\\
$1\,048\,576$ & $2\,668$ & $4\,716$  & $[3.57933,\,3.59681]\times10^{-4}$\\
$4\,194\,286$ & $8\,788$ & $10\,836$ & $[1.30227,\,1.31193]\times10^{-4}$\\
$4\,194\,304$ & $8\,706$ & $10\,754$ & $[5.22515,\,5.24432]\times10^{-4}$\\
$16\,777\,216$& $29\,212$& $31\,260$ & $[2.10371,\,2.11076]\times10^{-4}$\\
\bottomrule
\end{tabular}
\end{table}

The two moduli near $2^{22}$ have disjoint enclosures and balanced
norms differing by about a factor of four. These five moduli provide
finite-scale information about the restriction in \cref{eq:balance}.

\section{Separated windows}
\label{sec:rect}

To study a different part of the same operator, take supports in
$[X,2X]$ and $[Y,2Y]$ in the unscaled variable $x$, where
$2\le X<N/4$ and $Y>4X$. Denote the norm of this rectangular
restriction by $\|\Trect\|$. For a physical weight $w(x)$ write
$F(r)=e^{r/2}w(e^r)$, where $r=\log x$, and set
\[
  W_N(r)=\bigl(e^{-2\alpha r}+e^{-r/4}\bigr)^{-1}.
\]
Substituting $f(u)=\sqrt N\,w(Nu)$ in \cref{eq:energy} yields
\begin{equation}
  \NN(F)^2=\frac1{2\pi}\int W_N(r)|F(r)-F'(r)|^2\,\ud r.
  \label{eq:logenergy}
\end{equation}
The derivative energy $\NNzero(F)^2=(2\pi)^{-1}\int W_N|F'|^2$
is a lower bound on \cref{eq:logenergy}, since integration by parts
adds the nonnegative term $(2\pi)^{-1}\int(W_N+W_N')|F|^2$.
On a log-window $[a,b]=[\log Z,\log(2Z)]$, Poincar\'e's inequality
also gives
\begin{equation}
  \NN(F)^2\le R_Z\,\NNzero(F)^2,\qquad
  R_Z=1+\left(1+2\alpha+\tfrac14\right)
        \frac{W_N(b)}{W_N(a)}\left(\frac{\log 2}{\pi}\right)^2.
  \label{eq:ratio}
\end{equation}

For the derivative energy, the Dirichlet Green kernel on $[a,b]$ is
\begin{equation}
 K(r,s)=\frac{2\pi}{L}\,
          \xi(\min(r,s))\bigl[L-\xi(\max(r,s))\bigr],
 \quad \xi(r)=\int_a^r W_N(t)^{-1}\,\ud t,\quad L=\xi(b).
 \label{eq:green}
\end{equation}
This is the usual Dirichlet kernel for the energy
$(2\pi)^{-1}\int|\partial_\xi F|^2\,\ud\xi$ in the coordinate $\xi$.
As $2X<N/2$, the arithmetic part of the form is a finite sum over
low-window prime powers $n$. For each $n$ the left functional
is the weighted evaluation at $n$ and the right functional is
the difference of the weighted sums over $m\equiv n$ and
$m\equiv -n\bmod N$ in the high window. The weight at an integer
$t$ in either log-window is $\Lambda(t)/\sqrt t$.

Let $G_X,G_Y$ be the Gram matrices of these two lists of
functionals, evaluated by \cref{eq:green} in the derivative
energies. The eigenvalues $\lambda_j$ of $G_XG_Y$ are the squared
singular values of the arithmetic rectangle, with zeros included
as needed, since the nonzero spectrum agrees with that of the
corresponding positive Gram compression. With $P_j=\tr((G_XG_Y)^j)$
and $P_2>0$, its norm $A_0$ satisfies
\begin{equation}
   \left(\frac{P_4}{P_2}\right)^{1/4}
          \le A_0\le P_4^{1/8}.
   \label{eq:traces}
\end{equation}
The left inequality is $P_4\le(\max\lambda_j)^2P_2$ and the right
inequality is $P_4\ge(\max\lambda_j)^4$. The case $P_2=0$ gives zero
arithmetic norm. Both traces can be evaluated with enclosing intervals
using integer fixed-point products of Gram entries.

The continuous rectangle pairs the two shifts $kN+x$ and $kN-x$. Write
$\rho_k=\SSS(kN)/\SSS(N)$ and let $\mathcal K$ be the set of integers
$k\ge1$ for which $[kN-2X,kN+2X]$ meets $[Y,2Y]$. Its normalized norm
is bounded above by
\begin{equation}
  C_{X,Y}=2\sqrt3\,X\log 2
    \left(W_N(\log X)^{-1}W_N(\log Y)^{-1}
                 \sum_{k\in\mathcal K}\frac{\rho_k^2}{(kN-2X)^2}\right)^{1/2}.
  \label{eq:contrect}
\end{equation}
In fact, for the high-window physical test $z$ extended by zero,
\[
 |z(kN+x)-z(kN-x)|
  \le\frac{\sqrt{2x}}{kN-2X}
       \left(\int_{kN-2X}^{kN+2X}t^2|z'(t)|^2\,\ud t\right)^{1/2}.
\]
The integration intervals are disjoint because $4X<N$. Sum by
Cauchy--Schwarz over $k$. Integrating against
$|F(\log x)|/\sqrt{x}$ uses
$\int_X^{2X}|F(\log x)|\,\ud x\le\sqrt{3/2}\,X\|F\|_{L^2(\ud r)}$,
followed by Poincar\'e's inequality. For the high-window transform
$G(r)=e^{r/2}z(e^r)$,
$\int t^2|z'(t)|^2\,\ud t=\int |G'-G/2|^2\,\ud r
\le\int |G'-G|^2\,\ud r$ by the zero boundary values.
These estimates with \cref{eq:logenergy} give \cref{eq:contrect}.

If $[L_0,U_0]$ encloses $A_0$, \cref{eq:ratio,eq:contrect} give
\begin{equation}
  \max\left(0,\frac{L_0}{\SSS(N)\sqrt{R_XR_Y}}-C_{X,Y}\right)
  \le\frac{\|\Trect\|}{\SSS(N)}
  \le\frac{U_0}{\SSS(N)}+C_{X,Y}.
  \label{eq:rectbounds}
\end{equation}
For \cref{tab:rect}, $X=\lfloor\sqrt N\rfloor$ and $Y$ is the
nearest integer to $N^2/X$. The arithmetic terms are enumerated
by sieving the required progressions. The Green calculation uses
45 decimal digits of interval precision and integer enclosures at
scale $2^{70}$. The singular-series cutoff is $P=10^5$.

\begin{table}[htbp]
\centering
\caption{Certified normalized norms of the rectangular restrictions
from $[X,2X]$ to $[Y,2Y]$, conditional on $\grh(N)$.
Endpoints are rounded outward.}
\label{tab:rect}
\small
\begin{tabular}{@{}r r r l@{}}
\toprule
$N$ & $X$ & $Y$ & Interval for $\|\Trect\|/\SSS(N)$\\
\midrule
$4\,194\,286$  & $2\,047$ & $8\,594\,057\,181$  & $[4.8569,\,5.3947]\times10^{-4}$\\
$4\,194\,304$  & $2\,048$ & $8\,589\,934\,592$  & $[4.7177,\,5.2008]\times10^{-4}$\\
$16\,777\,216$ & $4\,096$ & $68\,719\,476\,736$ & $[2.3338,\,2.5592]\times10^{-4}$\\
\bottomrule
\end{tabular}
\end{table}

At $N=4\,194\,286$ and $N=16\,777\,216$, the lower bounds in
\cref{tab:rect} exceed the corresponding balanced upper bounds in
\cref{tab:balanced}. The full norm is therefore strictly larger
than the balanced norm at both moduli.

\section{A finite-modulus lower bound for the full operator}
\label{sec:witness}

For $N=32\,762$ we verified a piecewise linear physical weight $w$
on 257 rational knots. A floating-point generalized eigenvector
selected the trial, which was then saved with rational node values.
Let $f(u)=\sqrt N\,w(Nu)$. On each shift line the product
of the two linear pieces is quadratic, so splitting at their
breakpoints permits exact Simpson integration. The prime-power sum
and diagonal are enclosed using scaled integers with outward
rounding of the logarithms. Monotonicity of the coefficient in
\cref{eq:energy} gives a cellwise upper bound on the energy.
Under $\grh(32762)$ the resulting bounds are
\begin{align*}
  \B(f,f)&\in[0.2501768904360,\,0.2502592036273],\\
  \E(f)&\le 1.0064463618199,
  \qquad \|T_{32762}\|>0.2485744893 .
\end{align*}
The comparison uses the original norm. For this modulus the hypothesis
R threshold
$\delta_N=(\log N)^{3/2}\sqrt{\SSS(N)/N}$ is below $0.212829603$,
and the norm threshold for the S window of fixed width parameter
$1/4$, namely $2\delta_N/\sqrt3$, is below $0.245754457$.
The certified Rayleigh quotient exceeds both finite-modulus thresholds
under $\grh(32762)$.

\section{A truncation criterion for the full norm}
\label{sec:cutoff}

Write $L=\log N$ in this section. Given $\kappa\ge1$, set
\[
  Z_N=N^{\frac32\log\log N+\kappa},
  \qquad
  E_N^{\mathrm{cut}}
   =\frac{1}{\SSS(N)}
      \sup_{F,G}
      \frac{|\B(f,g)|}{\NN(F)\NN(G)}.
\]
The supremum is over nonzero tests supported in the physical interval
$[N^{-1},Z_N]$ in both variables. Here $F(r)=e^{r/2}w(e^r)$,
$f(u)=\sqrt N\,w(Nu)$ and the same convention relates $G$ and $g$.

\begin{theorem}[Truncation on the singular-series scale]
\label{thm:cutoff}
Under $\grh(N)$ and for sufficiently large even $N$, uniformly
for $\kappa\ge1$,
\begin{equation}
 E_N^{\mathrm{cut}}\le\frac{\|T_N\|}{\SSS(N)}
 \le (1+O(L^{-1})) E_N^{\mathrm{cut}}+O(e^{-\kappa}+N^{-1/2}).
 \label{eq:cutoff}
\end{equation}
Consequently, for every choice $\kappa=\kappa(N)\to\infty$,
\cref{eq:goal} is equivalent to $E_N^{\mathrm{cut}}\to0$.
For instance, $\kappa=\varepsilon\log\log N$ yields an error term
$O_\varepsilon((\log N)^{-\varepsilon})$ in \cref{eq:cutoff}.
\end{theorem}

\subsection{Density and the character form}

The images of $F\in C_c^\infty(\mathbb R)$ are dense in the Hardy
space of $T_N$. Indeed, the strip Paley--Wiener representation
\cite[Theorem~2.1]{MaiQian} writes a Hardy function as
$h(s)=\int_{\mathbb R} A(r)e^{-sr}\,\ud r$, with
$\int W_N^{-1}|A|^2<\infty$ by Plancherel on the boundary lines.
Under this representation the functional for $F$ acts by
$\int (F-F')A\,\ud r$. This follows from the multiplier $1+s$
in the defining Mellin functional of~\cite[Proposition~7.4]{LiLiu}.
An element annihilated by all the test functionals would satisfy
$A'+A=0$ distributionally. The solution $A(r)=c e^{-r}$ belongs
to the required weighted space only for $c=0$.

In log coordinates the form in \cref{eq:form} is
\begin{equation}
  \mathcal A(F,G)-\mathcal D(F,G)-\mathcal M(F,G),
  \label{eq:split}
\end{equation}
where $\mathcal M$ is the continuous model, $\mathcal D$ is the
arithmetic diagonal and $\mathcal A$ includes that diagonal.
Set
\[
  \ell_\chi(F)
    =\sum_{n\ge2}\frac{\Lambda(n)\chi(n)}{\sqrt n}F(\log n),
  \qquad
  \|U(F)\|^2=\frac{2}{\varphi(N)}\sum_{\chi\ {\rm odd}}|\ell_\chi(F)|^2.
\]
Odd-character orthogonality gives
$\mathcal A(F,G)=\langle U(F),U(G)\rangle$ with the indicated normalized
inner product. We use $\|v\|_\chi^2=(2/\varphi(N))\sum_{\chi\ {\rm odd}}
|v_\chi|^2$ for this character norm.
For each odd $\chi$ let $\chi^*$ be its inducing primitive character,
of conductor $q_\chi$, and let $\mathcal Z_\chi$ be the multiset of
nontrivial zero ordinates under $\grh(N)$. Write
$\widehat F(t)=\int F(r)e^{itr}\,\ud r$ and
$Z_\chi(F)=\sum_{\gamma\in\mathcal Z_\chi}\widehat F(\gamma)$.
Integration by parts in \cref{eq:logenergy} and
$0<W_N'/W_N\le1/4$ show that $\NN(F)^2$ is uniformly equivalent to
$\int W_N(|F|^2+|F'|^2)$. In particular, multiplication by a
smooth $\theta$ with $0\le\theta\le1$ has norm at most
$1+\|\theta'\|_\infty$ in $\NN$. The product rule gives this bound
since $(2\pi)^{-1}\int W_N|F|^2\le\NN(F)^2$.

We claim that for $H\ge2$ and $F$ supported in $[H,\infty)$,
\begin{equation}
       \|U(F)\|\ll L e^{-H/L}\NN(F).
       \label{eq:Utail}
\end{equation}
We first bound the short-interval zero counts in character mean square:
\begin{equation}
 \left\|\#\{\gamma\in\mathcal Z_\chi:|\gamma-t|\le 1/L\}\right\|_\chi
    \ll 1+\frac{\log(2+|t|)}{L}
    \qquad(t\in\mathbb R).
 \label{eq:shortzerocount}
\end{equation}
Choose a fixed nonnegative Schwartz function $\vartheta$ with
$\vartheta\ge1$ on $[-1,1]$ and Fourier transform supported in
$[-1/4,1/4]$ in the convention above. For example, take a positive
multiple of the square of the transform of a smooth even nonnegative
bump supported in $[-1/8,1/8]$. The transform is real on the real
axis and bounded away from zero on $[-1,1]$.
Apply the Weil explicit formula to the even majorant
$\vartheta(L(\gamma-t))+\vartheta(L(\gamma+t))$
for $\chi^*$ \cite[Section~2]{HughesRudnick}.
The conductor and gamma terms are $O(1+\log(2+|t|)/L)$, using
$\log q_\chi\le L$ and the integral of the Schwartz weight
against $O(1+\log(2+|u|))$ for the gamma term.
The prime terms have the form of two character polynomials, in
$\chi^*$ and $\overline{\chi^*}$, with coefficients
$O(\Lambda(n)/(L\sqrt n))$ supported on $n\le N^{1/4}$.
Replacing $\chi^*$ by $\chi$ in them has pointwise error
$O(L^{-1}\sum_{p\mid N}\log p\sum_{a\ge1}p^{-a/2})=O(1)$.
Odd-character orthogonality leaves only the diagonal in the squared
norm of each polynomial, since $n,m\le N^{1/4}$ implies $n+m<N$.
The squared norm is at most
$O(L^{-2}\sum_{n\le N^{1/4}}\Lambda(n)^2/n)=O(1)$ by Chebyshev's
bound and partial summation. Nonnegativity on the real zero
ordinates then proves \cref{eq:shortzerocount}.

For $F$ supported in any interval of length $O(L)$, translate it
to $b$ supported in $[-CL,CL]$. This leaves $|\widehat F|$ unchanged.
On the bins $I_m=[m/L,(m+1)/L)$ put $t_m=m/L$ and
$M_m=\sup_{I_m}|\widehat b|$. The local Sobolev bound
$\sup_I|v|^2\ll L\int_I|v|^2+L^{-1}\int_I|v'|^2$ and
Plancherel give
\begin{equation}
 \sum_m M_m^2\ll L\|F\|_2^2,\qquad
 \sum_m(1+t_m^2) M_m^2\ll L\|F\|_{H^1}^2.
 \label{eq:shortzerosampling}
\end{equation}
For the second bound, multiplication of $b(r)$ by $r$ costs
$O(L)$ in $H^1$ on this support. Apply the triangle inequality in
the character norm to the zero sums in the bins, and apply
Cauchy--Schwarz over the bins using \cref{eq:shortzerocount,eq:shortzerosampling}.
For $T\ge2$ this gives
\begin{align}
 \left\|\sum_{|\gamma|\le T}\widehat F(\gamma)\right\|_\chi
     &\ll L\sqrt T\left(1+\frac{\log(2+T)}{L}\right)\|F\|_2,
       \label{eq:zeroslow}\\
 \left\|\sum_{|\gamma|>T}\widehat F(\gamma)\right\|_\chi
     &\ll \frac{L}{\sqrt T}\left(1+\frac{\log(2+T)}{L}\right)\|F\|_{H^1}.
       \label{eq:zeroshigh}
\end{align}
For the latter estimate, the sum of
$(1+\log(2+|t_m|)/L)^2/(1+t_m^2)$ over bins meeting
$\{|t|>T\}$ is $O((L/T)(1+\log(2+T)/L)^2)$.

For a test supported in $r\ge H\ge2$, the smoothed explicit
formula \cite[Section~16.1]{Helfgott} is
\begin{equation}
 \ell_{\chi^*}(F)=-Z_\chi(F)
       -\sum_{a\ge0}\int F(r)e^{-(3/2+2a)r}\,\ud r .
 \label{eq:poslogexplicit}
\end{equation}
Here the reflected prime sum and the terms at $x=1$ vanish.
The trivial-zero term is $O(e^{-3H/2}\|F\|_2)$.
To pass from $\chi^*$ to $\chi$, for each $p\mid N$ the
geometric sum of $p^{-a/2}$ over $a\log p\ge H$ is
$O(e^{-H/2})$. The deleted prime terms therefore cost
$O(L e^{-H/2}\|F\|_\infty)$.
Using \cref{eq:zeroslow,eq:zeroshigh} with $T=2$ proves
$\|U(F)\|\ll L\|F\|_{H^1}$ on an interval of length $O(L)$.
For general support in $[H,\infty)$, use a smooth partition into
intervals of length $O(L)$ with bounded overlap and derivatives
$O(L^{-1})$, so the $j$th piece, $j\ge0$, is supported at
$r\ge H+jL-O(L)$. Their squared $\NN$ norms sum to $O(\NN(F)^2)$.
Since $W_N(r)\gg e^{2r/L}$ for $r\ge0$, the $j$th piece has
$H^1$ norm at most $O(e^{-H/L-j})$ times its $\NN$ norm.
The triangle inequality and Cauchy--Schwarz over the pieces now
prove \cref{eq:Utail}. The same argument applies directly to
the zero sums in \cref{eq:zeroshigh}, giving on $r\ge H$
\begin{equation}
 \left\|\sum_{|\gamma|>T}\widehat F(\gamma)\right\|_\chi
     \ll \frac{L e^{-H/L}}{\sqrt T}
         \left(1+\frac{\log(2+T)}{L}\right)\NN(F).
 \label{eq:globalzeroshigh}
\end{equation}

\subsection{Bounds for the diagonal and continuous model}

On log packets supported in $[j-1,j+1]$ and $[k-1,k+1]$ the diagonal
vanishes unless their supports overlap. In an overlapping pair its
norm is $O((1+j)e^{-2j/L})$ for $j\ge1$, by the elementary estimate
$\sum_{e^{j-1}\le n\le e^{j+1}}\Lambda(n)^2/n\ll 1+j$ and Sobolev
point evaluation. The finitely many packets around $r=0$ have
bounded diagonal norm, and the diagonal vanishes whenever one
argument is supported in $r\le0$.

For the model, expand the singular series in \cref{eq:series} as
\[
   \frac{\SSS(kN)}{\SSS(N)}
     =\sum_{d\mid k} a_N(d),\qquad
   a_N(d)=\prod_{p\mid d}\frac1{p-2}
\]
for squarefree $d$ coprime to $2N$, with $a_N(d)=0$ otherwise.
Uniformly in $N$, $\sum_{d\le D}a_N(d)\ll 1+\log D$ for $D\ge1$.
To see this, expand each factor as
$1/p+2/[p(p-2)]$ and use the harmonic sum for the terms in $1/p$.
The remaining coefficients have bounded sum, by convergence of
$\prod_{p>2}(1+2/[p(p-2)])$.

Take two packets with $j\le k$, and let $w,z$ be their physical tests.
For $d$ fixed, write $Q=dN$ and let $z_o$ be the odd extension of $z$.
The model contribution divided by $\SSS(N)$ is the sum over $d$ of
$a_N(d)$ times the pairing of $w$ with
\[
  \sum_{\ell\in\mathbb Z} z_o(x+\ell Q)-z(x).
\]
The $d$-contribution vanishes for $dN>2e^{k+1}$, since there is then
no nonzero translated term. The full lattice sum is bounded by
$\int|z_o'|$, by comparing with its zero integral divided by $Q$
on consecutive intervals of length $Q$. Subtracting $z(x)$ costs
at most $\|z\|_\infty$. If $s_j=W_N(j)^{-1/2}$, the packet bounds are
$\|w\|_{L^1}\ll e^{j/2}s_j\NN(F_j)$ and
$\|z\|_\infty+\int|z_o'|\ll e^{-k/2}s_k\NN(G_k)$.
These follow from Sobolev's inequality on the log packets and
the definition of $w,z$. Thus the model block has normalized norm
at most
\begin{equation}
  O\bigl((1+(k-L)_+)s_js_k e^{-(k-j)/2}\bigr),
  \label{eq:Mpacket}
\end{equation}
and vanishes for $2e^{k+1}<N$. Symmetry gives the bound with $j,k$
interchanged. The logarithmic derivative bound on $W_N$ yields
$s_j/s_k\le e^{(k-j)/8}$ for $j\le k$. For $k\ge0$ also
$s_k^2\le2e^{-2k/L}$. The sum of $e^{-3|k-j|/8}$ is bounded
in either index.

Applying these matrix bounds to the $\ell^2$ sequences of packet
norms proves the following estimates for unit tests in $\NN$:
\begin{align}
 \frac{|\mathcal D(F,G)|+|\mathcal M(F,G)|}{\SSS(N)}
    &\ll L, \label{eq:DMtotal}\\
 \frac{|\mathcal D(F,G)|+|\mathcal M(F,G)|}{\SSS(N)}
    &\ll (1+H)e^{-2H/L}
       &&\text{if either test is supported in $r\ge H$, $H\ge L$}.
       \label{eq:DMtail}
\end{align}
We used $\SSS(N)\gg1$ for the diagonal. If one test is
supported in $r\le0$, the model norm is $O(N^{-1/2}\SSS(N))$.
Indeed, for $j\le1$ the factor $s_j e^{j/2}$ is $O(e^{3j/8})$,
while the other index satisfies $k\ge L-O(1)$. Formula
\cref{eq:Mpacket} is then bounded by a product of two summable
sequences whose product of $\ell^2$ norms is $O(e^{-L/2})$.

\subsection{Comparison of the two norms}

Let $H=\log Z_N-L$. Split any smooth unit test into
$F=F_0+F_-+F_+$ by smooth multipliers valued in $[0,1]$, with
derivative $O(1/L)$. The lower transition lies on $[-L,0]$ and the
upper transition on $[H,H+L]$. Thus $F_0$ is supported in
$[-L,\log Z_N]$, $F_-$ in $r\le0$ and $F_+$ in $r\ge H$.
Each has norm at most $1+O(1/L)$. Use the same decomposition
for $G$. The term with both central pieces has normalized norm
at most $(1+O(1/L))E_N^{\mathrm{cut}}$. For the other terms
\cref{eq:DMtail} and the lower-end bound remove the diagonal and
model, up to $O(N^{-1/2}+(1+H)e^{-2H/L})$.

The character form vanishes on the lower pieces. By
\cref{eq:split,eq:DMtotal} its positivity gives
\[
  \|U(F_0)\|^2\ll \SSS(N)(E_N^{\mathrm{cut}}+L).
\]
With \cref{eq:Utail} this bounds each mixed central--upper character
term divided by $\SSS(N)$ by
$O(L e^{-H/L}\sqrt{E_N^{\mathrm{cut}}+L})$.
The upper--upper term is $O(L^2e^{-2H/L})$ on the same scale.
Since $e^{-H/L}=L^{-3/2}e^{1-\kappa}$, the mixed character bound
is $O(e^{-\kappa}\sqrt{1+E_N^{\mathrm{cut}}/L})$.
Use $\sqrt{1+t}\le1+t/2$ for $t\ge0$ to absorb its
$E_N^{\mathrm{cut}}$ contribution in the factor $1+O(L^{-1})$.
The upper--upper bound and $(1+H)e^{-2H/L}$ are both $O(e^{-\kappa})$,
uniformly in $\kappa\ge1$. This gives the upper inequality
in \cref{eq:cutoff}. Density gives the full
operator norm on the left of this comparison. The reverse
inequality follows from the definition of the restriction.

\section{A finite matrix criterion}
\label{sec:mesh}

We keep $L=\log N$ and take $3\le \tau\le L$. Use the cutoff in
\cref{thm:cutoff} with $\kappa=\log\tau$, and set
\begin{equation}
 r_-=-L,\qquad r_+=L(\tfrac32\log L+\log\tau),
 \qquad h=L^{-3}\tau^{-2}.
 \label{eq:grid}
\end{equation}
Partition $[r_-,0]$ into equal cells of length at most $h$ using
the smallest possible number of cells. On each interval
$[qL,\min\{(q+1)L,r_+\}]$ with $q\ge0$ an integer and $qL<r_+$,
do the same with maximum cell length $he^q$. In particular, the
mesh can be coarser at larger positive $r$. All cell lengths are
less than one, since $he^{r_+/L}=L^{-3/2}\tau^{-1}$.

For each $F\in H^1_0([r_-,r_+])$, write $F(r)=e^r V(\xi(r))$ with
\begin{equation}
 \xi(r)=\int_{r_-}^r W_N(t)^{-1}e^{-2t}\,\ud t.
 \label{eq:exactcoord}
\end{equation}
The integrand is $e^{-(2+2\alpha)t}+e^{-9t/4}$, and
\cref{eq:logenergy} becomes exactly
\begin{equation}
 \NN(F)^2=\frac1{2\pi}\int_0^{\xi(r_+)} |V'(\xi)|^2\,\ud\xi.
 \label{eq:exactdirichlet}
\end{equation}
Let $P_hF$ be $e^r$ times the piecewise linear interpolant of $V$
on the transformed nodes $\xi(r_i)$. This is an orthogonal projection
in $\NN$. On each cell the interpolant has constant derivative with
respect to $\xi$, and the residual vanishes at both endpoints.
Thus $\NN(P_hF)$ and $\NN(F-P_hF)$ are each bounded by $\NN(F)$.
Define
\[
 E_N^{\mathrm{mesh}}(\tau)
 =\sup_{F,G\in\operatorname{ran}P_h\setminus\{0\}}
   \frac{|\B(f,g)|}{\SSS(N)\NN(F)\NN(G)}.
\]

\begin{theorem}[Finite matrix criterion]
\label{thm:mesh}
Assume $\grh(N)$ and let $N$ be sufficiently large and even.
Uniformly for $3\le \tau\le L$,
\begin{equation}
 E_N^{\mathrm{mesh}}(\tau)
 \le \frac{\|T_N\|}{\SSS(N)}
 \le (1+O(L^{-1}))E_N^{\mathrm{mesh}}(\tau)+O(\tau^{-1}).
 \label{eq:mesh}
\end{equation}
The dimension is $O(L^4\tau^2)$. In particular, taking
$\tau=\log L$ makes \cref{eq:goal} equivalent to
$E_N^{\mathrm{mesh}}(\log L)\to0$ on matrices of dimension
$O(L^4(\log L)^2)$.
\end{theorem}

\subsection{Control of the interpolation residual}

Let $R=F-P_hF$. Each smooth localization $R_*$ of $R$ still
vanishes at the mesh nodes and the endpoints of the retained
interval. If the cells meeting its support have length at most
$h_*$, the fundamental theorem of calculus and Poincar\'e's
inequality on each cell give
\begin{equation}
 \|R_*\|_{L^\infty}\ll h_*^{1/2}\|R_*'\|_{L^2},
 \qquad
 \|R_*\|_{L^2}\ll h_*\|R_*'\|_{L^2}.
 \label{eq:meshPoincare}
\end{equation}
The same zero conditions hold when a packet crosses a boundary of
the retained interval, since all functions are extended by zero.

To bound $\|U(R)\|$, first split off the part with $r\le3$
by a smooth cutoff with transition on $[2,3]$. This part
contributes $O(h^{1/2}\NN(R))$ by direct point evaluation at
the finitely many positive prime powers with log at most three.
Partition the rest into packets of length $O(L)$ as in the
proof of \cref{eq:Utail}. For a packet $R_*$ near $r=jL$, $j\ge0$,
take $h_*$ to be the largest length allowed by the mesh on any
cell meeting the packet, so $h\le h_*\ll h e^j$.
Apply \cref{eq:zeroslow,eq:zeroshigh} with $T=1/h_*$.
By \cref{eq:meshPoincare} and $\log(1/h)\ll\log L$,
the whole zero sum is $O(L h_*^{1/2}\|R_*\|_{H^1})$ in
character norm. The trivial-zero and imprimitive terms in
\cref{eq:poslogexplicit} obey the same estimate, since their
bounds use $\|R_*\|_2$ and $L\|R_*\|_\infty$. These arguments
extend from smooth to $H^1$ packets by continuity.
As $\|R_*\|_{H^1}\ll e^{-j}\NN(R_*)$, summing with
Cauchy--Schwarz over the bounded-overlap packets gives
\begin{equation}
 \|U(R)\|\ll L\sqrt h\,\NN(F).
 \label{eq:meshU}
\end{equation}
Here $\sum_{j\ge0}e^{-j}\ll1$ and the squared packet norms
sum to $O(\NN(R)^2)$.

For the diagonal and the model, use the unit log packets of
\cref{sec:cutoff}, with cell length bound $h_j=h\exp((j+1)_+/L)$
on the packet $[j-1,j+1]$. For the diagonal,
\cref{eq:meshPoincare} adds $h_j^{1/2}$ to one
point-evaluation factor in the proof of \cref{eq:DMtotal}. Since
$(1+j)e^{-2j/L}h_j^{1/2}\ll Lh^{1/2}$ for $j\ge0$, we have
$|\mathcal D(R,G)|\ll Lh^{1/2}\NN(F)\NN(G)$. For the model,
take two packets indexed by $j\le k$, at least one of which is a
residual piece. In the estimate leading to \cref{eq:Mpacket},
the physical $L^1$ norm of either residual packet has a factor at most $h_k$
by \cref{eq:meshPoincare}. We can apply the lattice quadrature
bound to either packet by the symmetry of $\mathcal M$. Thus
the packet coefficient in \cref{eq:Mpacket} can use
\[
     \min\{ e^{-(k-j)/2},\ h_k e^{(k-j)/2}\}
\]
in place of $e^{-(k-j)/2}$. If the residual is in the lower packet,
one can use the factor $h_k e^{-(k-j)/2}$ directly.
For nonzero blocks $k\ge L-O(1)$ and
$(1+(k-L)_+)s_k^2\ll L e^{-k/L}$.
If $j\ge0$, then $s_j/s_k\ll e^{(k-j)/L}$, so the
coefficient is at most
\[
 C L\min\{e^{-(1/2-1/L)(k-j)},\ h e^{(1/2+1/L)(k-j)}\},
\]
using $e^{-k/L}h_k\ll h$. Its sum in either index is
$O(Lh^{1/2-1/L})=O(L\sqrt h)$ by splitting at $k-j=\log(1/h)$.
If $j<0$, use $s_j/s_k\le e^{(k-j)/8}$ and $k-j\ge L-O(1)$
to obtain row and column sums $O(L e^{-3L/8})$.
Consequently,
\begin{equation}
 \frac{|\mathcal D(R,G)|+|\mathcal M(R,G)|}{\SSS(N)}
         \ll L\sqrt h\,\NN(F)\NN(G).
 \label{eq:meshDM}
\end{equation}

\subsection{Norm comparison and matrix entries}

For two unit tests in the retained interval, split each into its
projection and residual. Positivity of the character form
and \cref{eq:DMtotal} give
\[
 \|U(P_hF)\|^2\ll\SSS(N)(L+E_N^{\mathrm{mesh}}(\tau)).
\]
For the mesh in \cref{eq:grid}, \cref{eq:meshU} is $O(L^{-1/2}\tau^{-1})$
on unit tests. Each mixed character contribution, normalized
by $\SSS(N)$, therefore has size
$O(\tau^{-1}\sqrt{1+E_N^{\mathrm{mesh}}(\tau)/L})$.
The contribution with two residuals is $O(L^{-1}\tau^{-2})$.
Equation \eqref{eq:meshDM} adds $O(L\sqrt h)=O(L^{-1/2}\tau^{-1})$.
The term with both projected tests is bounded
by $E_N^{\mathrm{mesh}}(\tau)$ because $P_h$ has norm one.
Using $\sqrt{1+t}\le1+t/2$ and applying \cref{thm:cutoff}
gives the upper bound in \cref{eq:mesh}. The lower bound follows
by restriction. The number of cells is
$O(L/h+\sum_{q\ge0}(L/h)e^{-q}+\log L)=O(L/h)$, proving the
dimension bound.

For an explicit matrix, let $\ell_i$ be the linear hat function at
the interior node $\xi_i=\xi(r_i)$ and let
$F_i(r)=e^r\ell_i(\xi(r))$. Define $f_i$ from $F_i$ as in
\cref{sec:cutoff}, and let
$B_{ij}=\B(f_j,f_i)/\SSS(N)$ using \cref{eq:form}.
The Gram matrix $G$ in \cref{eq:exactdirichlet} is tridiagonal with
\[
 G_{ii}=\frac{1}{2\pi}
       \left(\frac{1}{\xi_i-\xi_{i-1}}
                  +\frac{1}{\xi_{i+1}-\xi_i}\right),
 \qquad
 G_{i,i+1}=G_{i+1,i}=-\frac{1}{2\pi(\xi_{i+1}-\xi_i)}.
\]
Then $E_N^{\mathrm{mesh}}(\tau)$ is the largest absolute
generalized eigenvalue of $Bv=\lambda Gv$. The small number of
rows is a statement about the test space. Evaluating its entries
still involves prime-power sums and model integrals at scales up
to the cutoff.

\subsection{Trace norm and spectral compression}

Use $\|\cdot\|_{S_1}$ for the trace norm, the sum of singular
values. We view the forms in \cref{eq:split} as operators on
the Hilbert space with norm $\NN$, identified with the Hardy
space by the dense test vectors in \cref{sec:cutoff}.

\begin{theorem}[Trace norm at a fixed modulus]
\label{thm:tracecompression}
Assume $\grh(N)$ and write $L=\log N$. Then
\begin{equation}
 \|\mathcal A\|_{S_1}+\|\mathcal D\|_{S_1}
       +\|\mathcal M\|_{S_1}/\SSS(N)\ll L^2.
 \label{eq:traceterms}
\end{equation}
In particular, $T_N/\SSS(N)$ is trace class, with trace norm
$O(L^2)$. For $3\le\tau\le L$, the finite matrix criterion
in \cref{thm:mesh} admits a compression of dimension
$O(L^2\tau)$. If $E_N^{\rm comp}(\tau)$ is the normalized
norm on this compressed space, then
\begin{equation}
 E_N^{\rm comp}(\tau)\le \frac{\|T_N\|}{\SSS(N)}
  \le (1+O(L^{-1}))E_N^{\rm comp}(\tau)+O(\tau^{-1}).
 \label{eq:compressedmesh}
\end{equation}
The space is chosen from the separate matrices of
$\mathcal A$ and $\mathcal D+\mathcal M$.
\end{theorem}

For the character term,
$\|\mathcal A\|_{S_1}= (2/\varphi(N))\sum_{\chi\ {\rm odd}}
\|\ell_\chi\|_{\NN^*}^2$ by positivity and the finite sum
of rank-one operators. We estimate the mean of these
squared dual norms. First separate the part of the test
with $r\le3$ by a fixed smooth cutoff, leaving a part
$F_+$ supported in $r\ge2$. The omitted evaluations at
prime powers with $\log n\le3$ have bounded dual norm.
In \cref{eq:poslogexplicit}, the trivial-zero and
imprimitive terms cost $O(L)$ in each dual norm, by
the estimates following that equation.

Put $\rho(v)=1/\cosh v$ and write $F_+(r)=\rho(r/L)b(r)$.
Since $W_N(r)\asymp e^{2r/L}$ on $r\ge0$ and
$|\rho'/\rho|\le1$, the product rule gives
$\|b\|_{H^1}\ll\NN(F)$. The dual norm of the truncated
zero functional is therefore bounded by a constant
times the $H^{-1}(\mathbb R)$ norm of
\[
 g_{\chi,T}(r)=\rho(r/L)
       \sum_{\substack{\gamma\in\mathcal Z_\chi\\|\gamma|\le T}}
             e^{ir\gamma}.
\]
Its Fourier transform is
$L\sum_{|\gamma|\le T}\widehat\rho(L(t+\gamma))$ in the
convention above. On bins
$\gamma+t\in[m/L,(m+1)/L)$, the Fourier weight is
$O((1+|m|)^{-3})$. The short zero count
\cref{eq:shortzerocount} and the triangle inequality in
character mean square give
\[
 \|\widehat g_{\chi,T}(t)\|_\chi\ll L+\log(2+|t|)
\]
uniformly in $T$. Indeed the bin-count bound is
$O(1+L^{-1}\log(2+|-t+m/L|))$ and the bin weights
are summable against this bound. Plancherel with the
$H^{-1}$ weight $(1+t^2)^{-1}$ now gives
\[
 \frac2{\varphi(N)}\sum_{\chi\ {\rm odd}}
         \|g_{\chi,T}\|_{H^{-1}}^2
  \ll \int_{\mathbb R}\frac{(L+\log(2+|t|))^2}{1+t^2}\,\ud t
  \ll L^2.
\]
The truncated functionals converge on compact smooth
tests by the ordinary zero count and rapid Fourier
decay. Their dual norms are lower semicontinuous under
this convergence. With the two error terms this proves
the trace estimate for $\mathcal A$.

The diagonal is a positive sum of point evaluations.
At $r\ge0$ the squared dual evaluation norm is
$O(e^{-2r/L})$ by local Sobolev embedding. On an interval
$[j,j+1]$, $j\ge0$, Chebyshev's bound gives
$\sum_{\log n\in[j,j+1]}\Lambda(n)^2/n\ll 1+j$.
Its trace is thus bounded by
$O(\sum_{j\ge0}(1+j)e^{-2j/L})=O(L^2)$.

For the continuous term, use the expansion in
\cref{sec:cutoff}. For $d\ge1$ put $D=dN$ and let
$\mathcal C_D$ be the form of odd periodization on compact
tests, before subtracting the zero translate:
\[
 \mathcal C_D(F,G)=
   \int_0^\infty w(x)\overline{\sum_{\ell\in\mathbb Z}
                                 z_o(x+\ell D)}\,\ud x .
\]
Then $\mathcal M/\SSS(N)=\sum_{d\ge1}a_N(d)
(\mathcal C_{dN}-\mathcal I)$, where
$\mathcal I(F,G)=\int F(r)\overline{G(r)}\,\ud r$.
The periodic bilinear form is half the $L^2(0,D)$ pairing
of the two periodized odd extensions. Its sine expansion
therefore gives
\begin{equation}
 \mathcal C_D(F,G)=\frac4D\sum_{m\ge1} u_{m,D}(F)
                                  \overline{u_{m,D}(G)},
 \qquad
 u_{m,D}(F)=\int_0^\infty w(x)\sin(2\pi mx/D)\,\ud x .
 \label{eq:periodtr}
\end{equation}

Localize both tests to unit log packets with centers $j,k$.
On the $j$th packet, set $X=e^j$ and $s_j=W_N(j)^{-1/2}$.
The $L^1$ bounds for the physical weight and its derivative
are $O(s_j\sqrt X\,\NN(F))$ and
$O(s_jX^{-1/2}\NN(F))$. Integration by parts in
\cref{eq:periodtr} thus gives the packet dual bound
\[
 \|u_{m,D}\|\ll s_j\sqrt X\min(1,D/(mX)).
\]
If $j\le k$ and $V=e^k$, summing the trace norms of the
rank-one terms across the two packets gives at most
\[
 \frac{Cs_js_k\sqrt{XV}}{D}
 \sum_{m\ge1}\min(1,D/(mX))\min(1,D/(mV))
    \ll s_js_k(1+k-j)e^{-(k-j)/2}.
\]
The last bound follows by replacing the decreasing summand
with its integral and splitting at $D/V$ and $D/X$.
The $\mathcal I$ term is zero for disjoint packets.
On either packet the inclusion into unweighted $L^2$
has Hilbert--Schmidt norm $O(s_j)$, by integrating
the squared dual evaluation norm over the packet.
The trace norm of the product of two Hilbert--Schmidt
maps is bounded by the product of their norms.
Thus for overlapping packets the $\mathcal I$ term
obeys the same displayed bound up to a constant.

The difference $\mathcal C_D-\mathcal I$ is zero when
$D>2e^{k+1}$ on these packets, since all nonzero translates
and reflections then miss the two supports.
The localization maps have bounded operator norms, so the
same block trace bounds hold after localization on the full
space. Now $s_j/s_k\le e^{(k-j)/8}$ for $j\le k$.
Summing first over $j\le k$ and then over
$k\ge\log(D/2)-1$ gives $O(LD^{-2/L})$, since
$s_k^2\ll e^{-2k/L}$ there. The opposite ordering is
identical by adjunction. In particular the blocks converge
in trace norm, even though $\mathcal C_D$ and $\mathcal I$
were defined separately on compact tests only.
Finally $\sum_{d\le u}a_N(d)\ll1+\log u$ implies
$\sum_{d\ge1}a_N(d)d^{-2/L}\ll L$ by partial summation.
The series for $\mathcal M/\SSS(N)$ thus has trace norm
$O(L^2)$ and agrees with the original form on compact
tests. This proves \cref{eq:traceterms}.

To get the compressed matrix, use the mesh in
\cref{thm:mesh} and form the two self-adjoint matrices
of $\mathcal A/\SSS(N)$ and
$(\mathcal D+\mathcal M)/\SSS(N)$ separately in an
orthonormal basis for its energy norm. Their trace bounds
persist after compression. Select the span of the
eigenvectors of these two matrices with absolute eigenvalue
greater than $\tau^{-1}$. The span has dimension
$O(L^2\tau)$ by \cref{eq:traceterms}. For its orthogonal
projection $P$, each matrix has norm at most
$\tau^{-1}$ when restricted to the complement in one
argument. Thus their difference changes by at most
$4\tau^{-1}$ in operator norm when compressed to $P$.
Apply \cref{thm:mesh} to get \cref{eq:compressedmesh}.
For example, $\tau=\log L$ gives an additive $o(1)$
error and dimension $O(L^2\log L)$. This construction
starts with the explicit mesh matrices of \cref{sec:mesh}.

The trace bound also gives a scalar criterion. Put
$S_N=T_N/\SSS(N)$ and $m_N=2\lceil\log L\rceil$.
Then \cref{eq:goal} is equivalent to
\begin{equation}
       \bigl(\operatorname{Tr} S_N^{m_N}\bigr)^{1/m_N}
              \longrightarrow 0.
 \label{eq:tracemoment}
\end{equation}
Indeed, for even $m$ the trace is the sum of the $m$th
powers of absolute eigenvalues. Its $m$th root lies between
$\|S_N\|$ and $(C L^2)^{1/m}\|S_N\|^{1-1/m}$ by
\cref{eq:traceterms}, with $(C L^2)^{1/m_N}=O(1)$.

\subsection{The signed trace and a logarithmic-derivative mean}

For $a>0$ define
\[
 \mathcal L_N(a)=\frac2{\varphi(N)}
    \sum_{\chi\ {\rm odd}}\int_{\mathbb R}
       \frac{|(L'/L)(1/2+a+it,\chi)|^2}{(1+a)^2+t^2}\,\ud t .
\]
The $L$-functions are for the characters modulo $N$, including
their imprimitive Euler factors when applicable.

\begin{theorem}[A necessary scalar mean-square law]
\label{thm:signedtrace}
Under $\grh(N)$, uniformly as $N$ tends to infinity through
even moduli, with $L=\log N$,
\begin{equation}
 \operatorname{Tr} T_N
   =\mathcal L_N(1/L)-\frac{\pi}{4}(1-e^{-2})L^2+O(L\log L).
 \label{eq:signedtrace}
\end{equation}
Moreover, with $S_N=T_N/\SSS(N)$,
\begin{equation}
 |\operatorname{Tr} S_N|
    \ll L^{3/2}\bigl(1+(\log L)\|S_N\|\bigr).
 \label{eq:signedtracenorm}
\end{equation}
Consequently, along any sequence on which
\cref{eq:goal} holds, one has
\begin{equation}
 \frac{\mathcal L_N(1/L)}{L^2}\longrightarrow
               \frac{\pi}{4}(1-e^{-2}).
 \label{eq:necessaryderivativelaw}
\end{equation}
\end{theorem}

The evaluation kernel for \cref{eq:logenergy} can be calculated
on the full line. With $F(r)=e^r V(\xi(r))$, use
$\xi(r)=\int_r^\infty W_N(t)^{-1}e^{-2t}\,\ud t$.
This maps the energy to $(2\pi)^{-1}\int_0^\infty|V'|^2$.
The evaluation kernel in $\xi$ is $2\pi\min(\xi,\xi')$,
since evaluation at $\xi$ is integration of the derivative
from zero to $\xi$. Derivatives of compactly supported tests
are dense in $L^2(0,\infty)$: the integral of a compact
derivative candidate can be adjusted by a bump of arbitrarily
small $L^2$ norm spread over a distant interval.
Transforming back gives
\begin{equation}
 K_N(r,s)=
 \sum_{a\in\{\alpha,1/8\}}\frac{\pi}{1+a}
               e^{-a(r+s)} e^{-(1+a)|r-s|}.
 \label{eq:fullenergykernel}
\end{equation}
The individual summand is also the kernel for the energy
with weight $e^{2ar}$ in \cref{eq:logenergy}.

For a finite prime-power sum in $\ell_\chi$, contracting
with \cref{eq:fullenergykernel} and applying the Fourier
transform of $(\lambda^2+t^2)^{-1}$, $\lambda=1+a$,
gives a sum of two weighted integrals of squared Dirichlet
polynomials, one for each $a$. Smooth truncations in
$r$ converge in $\NN^*$ by \cref{eq:Utail}. Since the
kernel is a sum of two positive kernels, the polynomial
limits exist in both weighted $L^2$ spaces. Under GRH,
partial summation using $\psi(x,\chi)\ll \sqrt x\log^2(Nx)$
for odd $\chi$ identifies the pointwise limits with
$-L'/L$ on the respective lines \cite[Section~4]{LeeGRH}.
Thus $\operatorname{Tr}\mathcal A=\mathcal L_N(\alpha)+
\mathcal L_N(1/8)$.

The second term is $O(1)$. To see this, split at
$r=(1/2)L$ with a smooth transition of width one.
The short Dirichlet polynomial has bounded character
mean square for each $t$ by odd orthogonality, since its
indices are below $e\sqrt N$ and no two have sum divisible
by $N$ for large $N$. The diagonal series on this line
is bounded by $\sum_n\Lambda(n)^2 n^{-5/4}<\infty$.
For the long part, the energy with weight $e^{r/4}$
dominates $\NN^2$ by a factor $\gg e^{2(1/8-\alpha)L/2}$
on its support. The character dual bound averaged in
\cref{eq:traceterms} therefore gives a power saving in the
mean squared dual norm in this stronger energy. Its kernel
is the $a=1/8$ summand of \cref{eq:fullenergykernel}, so
this controls the long part of the indicated integral.

For the diagonal, contraction at $r=s=\log n$ gives
\[
 \operatorname{Tr}\mathcal D
    =\sum_{a\in\{\alpha,1/8\}}\frac{\pi}{1+a}
        \sum_{(n,N)=1}\frac{\Lambda(n)^2}{n^{1+2a}}
    =\frac{\pi}{4}L^2+O(L).
\]
Indeed the prime number theorem under RH and partial
summation give $\sum_n\Lambda(n)^2 n^{-1-2a}
=(2a)^{-2}+O(1)$ uniformly for $0<a\le1/8$. Removing
primes dividing $N$ costs $O((\log L)^2)$:
$\sum_{p\mid N}(\log p)^2/(p-1)\ll(\log L)^2$ by
splitting at $p=L$ and using $\sum_{p\mid N}\log p\le L$
on the large primes.

To take the trace of the periodized model
$\mathcal C_D-\mathcal I$ from \cref{eq:periodtr},
write the physical kernel of the $a$ summand as
\[
 k_a(x,y)=\frac{\pi}{1+a}
        \frac{\min(x,y)^{1/2}}{\max(x,y)^{3/2+2a}} .
\]
Its trace contribution is
\begin{equation}
 \int_0^\infty \sum_{\ell\ne0}
        \operatorname{sgn}(x+\ell D)k_a(x,|x+\ell D|)
                   \,\ud x .
 \label{eq:periodtrintegral}
\end{equation}
To justify this integral, multiply both test arguments
by smooth log cutoffs that tend to one and have bounded
derivatives, with compact support at every stage. The
multipliers and their adjoints tend strongly to the identity.
For the adjoints this follows first on the dense span of
evaluation vectors, where they act by the cutoff values.
Convergence of the traces follows from the trace-norm
bound proved for the periodized model. With cutoffs, the
trace is an integral of the kernel by the rank-one
evaluation formula on the finitely many retained lines.
For large $x$, odd lattice summation in $y$ is bounded
by the total variation of the odd extension of $k_a(x,y)$
with the cutoff. This is $O(x^{-1-2\Re a})$, as is the
zero translate being subtracted. The kernel and the
cutoff times the kernel have the same variation bound,
since a derivative of a log cutoff contributes $O(1/y)$.
For $x<D/2$, pair the two signs at $mD\pm x$ for each
$m\ge1$. Differentiation in $y$ gives
$O(x^{3/2}(mD)^{-5/2-2\Re a})$ for the difference.
These estimates justify the removal of the cutoffs and
give a holomorphic integral on $\Re a>0$.

For $\Re a>1/2$ the individual line integrals in
\cref{eq:periodtrintegral} are absolutely summable. Scaling
a positive shift to one, the two difference terms integrate
to $2\mathrm B(3/2,2a)$ after extracting the factor
$\pi/(1+a)$. In the reflection integral the substitution
$u=x/(1-x)$ on $0<x<1/2$ gives
$2\int_0^1 u^{1/2}(1+u)^{2a-1}\,\ud u$.
Thus by analytic continuation \cref{eq:periodtrintegral}
equals $(\pi/(1+a))D^{-2a}\zeta(2a)J(a)$, where
\[
 J(a)=2\mathrm B(3/2,2a)-2\int_0^1 u^{1/2}(1+u)^{2a-1}\,\ud u.
\]
The pole at $a=1/2$ is removable because $J(1/2)=0$.
Summing the traces over $d$ using the absolute
trace-norm convergence gives
\begin{equation}
 \operatorname{Tr}\mathcal M
  =\sum_{a\in\{\alpha,1/8\}}\frac{\pi}{1+a}
      N^{-2a}\zeta(2a)J(a)\SSS(N)\zeta(1+2a)H_N(2a),
 \label{eq:modeltraceexact}
\end{equation}
with the Euler factor $H_N$ of \cref{sec:zeros}.
By \cref{eq:condconstant}, $\SSS(N)H_N(0)=1$. On
$0\le s\le2/L$, logarithmic differentiation of the Euler
product gives $H_N'/H_N=O(1+\sum_{p\mid N}\log p/(p-1))
=O(\log L)$, splitting at $p=L$ as above. The primes
not dividing $N$ have summable derivative bounds
$O((\log p)/p^2)$ after their linear terms cancel.
It follows that $\SSS(N)H_N(2/L)=1+O((\log L)/L)$.
Since $J(a)=1/a+O(1)$ as $a\downarrow0$ and $\zeta(0)=-1/2$,
the $a=\alpha$ term in \cref{eq:modeltraceexact} is
$-\pi e^{-2}L^2/4+O(L\log L)$. The fixed $a=1/8$
term is $o(1)$ by the $N^{-1/4}$ factor and the bounded
Euler factors at $1/4$. Combining the three trace evaluations
proves \cref{eq:signedtrace}.

For \cref{eq:signedtracenorm}, let $H=L\log L$ and
$h=L^{-1/2}$. Take a uniform mesh of maximum step $h$
on $[0,H]$ and let $P$ be the orthogonal projection
onto the span of its evaluation vectors. For $Q=1-P$,
each residual $QF$ vanishes at the nodes. Its parts on
compact log packets in $[0,H]$ have weighted $L^2$ norm
$O(h\NN(F))$ and weighted supremum norm
$O(h^{1/2}\NN(F))$ by the cell estimates from
\cref{sec:mesh}. Use smooth cutoffs with transition
width one to split off the regions $r\le1$ and
$r\ge H-1$.

For the character trace of the residual on the core
$2\le r\le H$, use the function $g_{\chi,T}$ from
the proof of \cref{eq:traceterms}. The residual
factored as $\rho(r/L)b(r)$ has
$\|b\|_{H^1}\ll\NN(F)$ and $\|b\|_2\ll h\NN(F)$.
The squared dual bound for its pairing with the zero
sum uses $\|b\|_2^2$ at frequencies $|t|\le1/h$ and
$\|b\|_{H^1}^2$ with weight $(1+t^2)^{-1}$ above
$1/h$. Averaging the bounds
$\|\widehat g_{\chi,T}(t)\|_\chi\ll L+\log(2+|t|)$
gives $O(L^2h)$ independently of $T$.
The prime evaluations at $r\le3$ cost $O(h)$ in the
mean squared dual norm. The explicit-formula remainders
on the core cost $O(L^2 h)$ by the residual $L^2$
and point-evaluation bounds. On $r\ge H-1$, the same
$H^{-1}$ estimate with $\rho((r-H+1)/L)$ instead gives
$O(L^2e^{-2H/L})$: the translated Fourier weights
have the same magnitudes, and the factored test now
has $H^1$ norm $O(e^{-H/L}\NN(F))$.
It follows that
$\|Q\mathcal A Q\|_{S_1}\ll L^2(h+e^{-2H/L})$.
The diagonal point evaluations give
$\|Q\mathcal D Q\|_{S_1}\ll L^2h+L(H+L)e^{-2H/L}$
by the same split.

For the continuous term, the sine functional bound on
the core residual packet now contains
$\min(h,D/(mX))$ instead of $\min(1,D/(mX))$.
In the trace-norm estimate of \cref{eq:periodtr} both
arguments are residuals. The product integral for the
two sine bounds gains a factor $h$ by rescaling its
integration variable by $h$. The zero translate also
gains $h$ from the two residual point evaluations.
This gives $O(h L^2)$ on the normalized trace scale
after summing over packets and periods. If either test
is supported in $r\ge H-1$, the previous packet sum
starts at larger index $\max(H,\log D)-O(1)$, so its
sum over $d$ is $O(L(H+L)e^{-2H/L})$. Here one splits
the sum at $d=e^H/N$ and uses
$\sum_{d\le u}a_N(d)\ll1+\log u$.
If either test is supported in $r\le1$, the packet
separation for $D=dN$ gives $O(D^{-1/4})$ in trace
norm for that period, by summing
$(1+k-j)e^{-3(k-j)/8}$ for $j\le2$ and
$k\ge\log D-O(1)$. This gives $O(N^{-1/4})$ on the
normalized scale. All these bounds hold after composing
with $Q$ and with the smooth cutoffs, which are bounded
operators. We have therefore shown
\[
 \|Q S_N Q\|_{S_1}\ll L^2h+L(H+L)e^{-2H/L}+N^{-1/4}.
\]
The trace on $P$ is at most $\operatorname{rank}(P)\|S_N\|$
in absolute value, and $\operatorname{rank}(P)\ll H/h+1$.
The cross terms have zero trace. Substitution gives
\cref{eq:signedtracenorm}. Finally $\SSS(N)\ll\log L$
by \cref{eq:subseries} and Mertens' bound applied to the
primes up to $L$, with the contribution of larger
prime divisors bounded by an absolute factor. If
\cref{eq:goal} holds, then \cref{eq:signedtracenorm}
gives $\operatorname{Tr}T_N=o(L^2)$, and
\cref{eq:signedtrace} proves \cref{eq:necessaryderivativelaw}.

\subsection{Local logarithmic-derivative means}

Put $c(b)=(1-e^{-2b})/(4b^2)$ for $b>0$. In this subsection
$a=b/L$, $\sigma=1/2+a$, and $D_\chi(t)=(L'/L)(\sigma+it,\chi)$.
Use $\mathbf E_N$ for the average over odd characters,
with weight $2/\varphi(N)$.

\begin{theorem}[A necessary local mean-square law]
\label{thm:localderivative}
Assume $\grh(N)$ and let $N$ tend to infinity along a sequence
on which $\|S_N\|\ll L^{1-\varepsilon}$ for some fixed
$0<\varepsilon\le1$. This holds in particular under
\cref{eq:goal}. For every fixed $b>0$,
\begin{equation}
 \frac{\mathbf E_N |D_\chi(t)|^2}{L^2}\,\ud t
           \ \longrightarrow\ c(b)\,\ud t
 \label{eq:localderivativelaw}
\end{equation}
against continuous compactly supported tests in $t$ and
on every bounded interval. The corresponding measures
with densities $\mathbf E_N D_\chi(t)/L$ and
$\mathbf E_N D_\chi(t)^2/L^2$ converge to zero already
under $\grh(N)$ against compactly supported continuous
tests. In particular the limit in
\cref{eq:localderivativelaw} is halved if $|D_\chi|^2$
is replaced by $(\Re D_\chi)^2$.
\end{theorem}

Use the pure exponential energy
$(2\pi)^{-1}\int e^{2ar}|F-F'|^2$ and let $T_{a,v}$
be the form of $T_N$ after multiplication of both tests
by $e^{-iv r}$, where $v\in\mathbb R$. Its evaluation
kernel is the $a$ term of \cref{eq:fullenergykernel}.
The arithmetic trace is therefore
\begin{equation}
 \operatorname{Tr}\mathcal A_{a,v}
     =\int_{\mathbb R}\frac{\mathbf E_N|D_\chi(t)|^2}
                             {(1+a)^2+(t-v)^2}\,\ud t .
 \label{eq:shiftedtracearith}
\end{equation}
The argument proving \cref{eq:traceterms} applies with
$\rho(ar)$ in the zero functional and with the phase
$e^{-iv r}$ placed in that functional. Binning the
transformed damping factor in intervals of length $1/L$
still gives the bound $O_b(L+\log(2+|t-v|))$ in the
character norm. The prime-power and trivial-zero terms
use the modulus of the phase only. Hence
\begin{equation}
 \operatorname{Tr}\mathcal A_{a,v}
                  \ll_b (L+\log(2+|v|))^2 .
 \label{eq:localderivativeupper}
\end{equation}
For example, the Fourier proof bounds the arithmetic
dual norm by the $H^{-1}$ norm of the shifted zero sum,
whose mean square integrates
$(L+\log(2+|t-v|))^2/(1+t^2)$.
With either character input restricted to $r\ge H$,
the squared dual bound for that input gains $e^{-2aH}$
by using $\rho(a(r-H))$ instead. The trace norm of the
resulting bilinear form then gains $e^{-aH}$ by the
Hilbert--Schmidt product bound. These estimates also
justify the polynomial limits in \cref{eq:shiftedtracearith},
as in the proof of \cref{eq:signedtrace}.

The diagonal trace is unaffected by the phase. Its
main term in the new energy is $\pi L^2/(4b^2)$ with
error $O_b(L)$. The preceding model trace calculation uses
$k_a(x,y)(x/y)^{-iv}$ in place of $k_a(x,y)$.
The variation and small-$x$ bounds in the proof of
\cref{eq:periodtrintegral} continue to hold with constants
depending on $v$. The resulting factor $J(a)$ in
\cref{eq:modeltraceexact} is replaced by
\[
 \mathrm B(3/2-iv,2a)+\mathrm B(3/2+iv,2a)
  -\int_0^1 (u^{1/2-iv}+u^{1/2+iv})(1+u)^{2a-1}\,\ud u .
\]
It again vanishes at $a=1/2$ and is $1/a+O_v(1)$
at zero. Thus for fixed $b,v$ the deterministic trace is
\begin{equation}
 \operatorname{Tr}(\mathcal D_{a,v}+\mathcal M_{a,v})
        =\pi c(b)L^2+O_{b,v}(L\log L).
 \label{eq:shiftedtracedet}
\end{equation}
Here the estimates for $\SSS(N)H_N(2a)$ from
\cref{eq:modeltraceexact} hold with $b$ fixed, giving
$1+O_b((\log L)/L)$.

For completeness we compare the signed trace in this
energy with the original norm, including when $b<1$.
Take $H=(\varepsilon/8)L\log L$ and $h=L^{-\varepsilon/4}$,
and cut the
test arguments smoothly above $H$ with a transition of
width one. The trace-norm cost for the discarded
character terms is $O_{b,v}(L^2 e^{-aH})$ by the
dual bounds above. For the diagonal it is
$O_b(L(H+L)e^{-2aH})$ by the atom count in
\cref{sec:mesh}. For the model, the packet estimate
following \cref{eq:periodtr} now uses $s_j=e^{-aj}$.
Summing the packets with larger index
$k\ge\max(H,\log(dN))-O(1)$ and then over $d$
costs $O_{b,v}(\SSS(N)L(H+L)e^{-2aH})$.
In particular the discarded trace norms total $o_{b,v}(L^2)$,
since $\SSS(N)\ll\log L$. The same estimates on the full
space with $k\ge\log(dN)-O(1)$ prove trace convergence
of the continuous form in the new energy.

On the retained part, take the orthogonal projection onto
the evaluation vectors at a mesh of maximum step $h$
on $[0,H+1]$ in the pure energy. As in the proof of
\cref{eq:signedtracenorm}, residuals vanish at the mesh
points. In the character trace the low Fourier frequencies
use the residual $L^2$ gain $h$ and the high frequencies
use the $H^{-1}$ weight, split at $1/h$. In the model
trace the two residual sine bounds contain $\min(h,D/(mX))$,
so the product integral gains $h$. The diagonal gains
$h$ by residual point evaluation. These arguments give
$O_{b,v}(\SSS(N)L^2h)$ in trace norm on the complement
of the mesh projection, plus $O_{b,v}(\SSS(N)N^{-1/4})$
for the low model region $r\le1$. The latter is the same
separated-packet estimate as in the proof of
\cref{eq:signedtracenorm}, since $a\le1/8$ for large $N$.

The projected trace has at most $O(H/h+1)$ dimensions.
On the retained supports the old energy norm of each
cutoff and modulated test is at most a constant times
$e^{(\alpha-a)_+H}$ times its pure exponential norm.
Indeed $W_N(r)\le e^{2\alpha r}$ for $r\ge0$ and
$W_N(r)\le e^{r/4}\le e^{2ar}$ for $r\le0$.
The cutoff and phase derivatives are bounded for fixed
$v$. Consequently the projected trace has absolute value
at most
\[
 O_{b,v}\bigl((H/h+1)\SSS(N)\|S_N\|
                       e^{2(\alpha-a)_+H}\bigr)=o(L^2)
\]
along the sequence of the theorem. In fact
$e^{2(\alpha-a)_+H}\le L^{\varepsilon/4}$, so with the
norm bound the power of $L$ before logarithms is at most
$2-\varepsilon/2$.
Together with the residual and discarded traces this
shows $\operatorname{Tr}T_{a,v}=o(L^2)$.
\Cref{eq:shiftedtracearith,eq:shiftedtracedet} give
\[
 \frac1{L^2}\int_{\mathbb R}
     \frac{\mathbf E_N|D_\chi(t)|^2}{(1+a)^2+(t-v)^2}\,\ud t
                        \longrightarrow \pi c(b)
 \qquad(v\in\mathbb R\ \text{fixed}).
\]

The limit of the Poisson integrals determines the local
limit. By \cref{eq:localderivativeupper}, the mass of a
unit interval at height $v$ is
$O_b((1+L^{-1}\log(2+|v|))^2)$ for the measures in
\cref{eq:localderivativelaw}. This makes the weighted
tails uniformly negligible and permits passage to any
vague subsequential limit in the Poisson integrals.
Every such limit is translation bounded and has
convolution with the Poisson kernel at height one
equal to $c(b)$. To see uniqueness, subtract $c(b)\,\ud t$.
The Poisson semigroup makes the convolution vanish at
every height greater than one. Harmonic uniqueness
then gives vanishing at all positive heights, and the
approximate identity recovers the zero measure.

The analogous convolutions for $D_\chi(t)$ and $D_\chi(t)^2$
are particularly simple. With $\lambda=1+a$, the integral
of either power against $((t-v)^2+\lambda^2)^{-1}$ is
$\pi/\lambda$ times that power of
$(L'/L)(\sigma+\lambda+iv,\chi)$. This follows first for
finite Dirichlet polynomials by the Fourier transform
of the Poisson kernel and then by the weighted $L^2$
limit in \cref{eq:shiftedtracearith}. The right side is
$O(1)$ since $\sigma+\lambda>3/2$. For the normalized
measures of $D_\chi$ or $D_\chi^2$, the same unit-interval
bounds hold for total variation, by Cauchy--Schwarz for
the first power and by \cref{eq:localderivativeupper}
for the second. Thus uniqueness gives their zero local limits.
Finally $2(\Re D_\chi)^2=|D_\chi|^2+\Re(D_\chi^2)$.
This finishes the proof of \cref{thm:localderivative}.

\section{Excluding pairs below the product threshold}
\label{sec:subcritical}

Write $x,y$ for the unscaled physical variables of \cref{sec:rect}.
For $Q\ge1$, retain in both the arithmetic and continuous terms of
\cref{eq:form} only pairs with
\[
       x,y\ge1,\qquad xy>N^2/Q.
\]
Denote the resulting form by $\B^{>Q}$ and, after transfer to the
Hardy space, the corresponding operator by $T_N^{>Q}$.

\begin{theorem}[Uniform bound below the product threshold]
\label{thm:subcritical}
For every fixed $\varepsilon>0$, under $\grh(N)$ and for all
sufficiently large even $N$, uniformly in $Q\ge1$,
\begin{equation}
  \frac{\|T_N-T_N^{>Q}\|}{\SSS(N)}
      \ll_\varepsilon Q^{-1/2}+N^{-1/2+\varepsilon}.
  \label{eq:subcritical}
\end{equation}
In particular, for any $Q=Q(N)\to\infty$, the full-norm
condition \cref{eq:goal} is equivalent to
$\|T_N^{>Q}\|/\SSS(N)\to0$.
\end{theorem}

\subsection{Uniform arithmetic bounds for the log packets}

We use the following consequence of the Selberg upper-bound
sieve \cite[Theorem~9.1]{Richert}. On an interval of length $H$,
suppose that at each prime $p$ the excluded set has cardinality
$0\le\nu(p)\le2$ with $\nu(p)<p$. The count after sifting by
primes up to $H^c$, for a sufficiently small fixed $c>0$, is at
most
\begin{equation}
  O\left(H\prod_{p\le H^c}(1-\nu(p)/p)+H^{2/5}\right).
  \label{eq:uniformsieve}
\end{equation}
Indeed, for sieve parameter $z=H^{1/10}$ the Selberg denominator
is $\sum_{d\le z}\mu^2(d)\prod_{p\mid d}\nu(p)/(p-\nu(p))$, with
the sum restricted to primes being sifted. Under the probability
weights obtained by normalizing the product on divisors of
$\prod_{p\le H^c}p$, the expectation of $\log d$ is
$\sum_{p\le H^c}(\nu(p)/p)\log p\ll c\log H$.
Taking $c$ small makes the probability of $d\le z$ at least $1/2$.
The denominator is therefore bounded below by half the inverse
local product. The Chinese remainder theorem gives a remainder
$O(\nu(d))$ at a squarefree modulus $d$, where
$\nu(d)=\prod_{p\mid d}\nu(p)$. In the Selberg upper bound the
total remainder is
$O(\sum_{d\le z^2}3^{\omega(d)}\nu(d))=O(z^4)$.
Here $\omega(d)$ counts distinct prime divisors and
$6^{\omega(d)}\ll d$ by separating the primes below seven.

Use the unit log packets in \cref{sec:cutoff}, with indices
$0\le j\le k$, and write $X=e^j$, $Y=e^k$. From \cref{eq:logenergy}
and Sobolev point evaluation, the physical weights on their supports
satisfy
\begin{equation}
  |w(x)|\ll X^{-1/2}\NN(F_j),\qquad
  |z(y)|\ll Y^{-1/2}\NN(G_k).
  \label{eq:packetpoint}
\end{equation}
For terms with $x,y\ge1$ and $xy\le N^2/Q$ in these packets we
have $XY\le e^2N^2/Q$. Each permissible pair lies on one of the
three lines $x+y=bN$ or $x-y=\pm bN$ with $b\ge1$.
Thus $Y\ge N/(2e)$ and $b\ll Y/N$. Also
\begin{equation}
 \sum_{b\le B}\SSS(bN)\ll B\SSS(N),\qquad
 \SSS(N)\asymp N/\varphi(N).
 \label{eq:subseries}
\end{equation}
The first bound follows from the divisor expansion in
\cref{sec:cutoff}, since
$\sum_d a_N(d)/d\le\prod_{p>2}(1+1/[p(p-2)])$.
The second follows by comparing the local factors in \cref{eq:series}.
Each continuous line has parameter length $O(X)$ in the smaller
coordinate, so its contribution below the threshold can be bounded
in absolute value. Equations \eqref{eq:packetpoint} and
\eqref{eq:subseries} give a normalized coefficient
$O(\sqrt{XY}/N)$ for the sum of the continuous line contributions.

Fix $\delta=\min(\varepsilon/2,1/8)$ for the arithmetic bounds.
If $X\ge N^\delta$, sieve along each line using the two linear
forms in the smaller variable. There is one excluded residue
when $p\mid bN$ and two otherwise. Both prime values exceed
the sifting level in \cref{eq:uniformsieve}. The Euler product
in that bound is $O(\SSS(bN)/(\log X)^2)$ by Mertens' theorem
and the two-prime local factors. The contribution of two primes
to the von Mangoldt-weighted count per line is then
$O_\delta(X\SSS(bN))$ because $\log Y\ll L\ll_\delta\log X$.
Pairs involving a proper prime power contribute
$O(\sqrt X\,L^3)$ per line. To see this, each coordinate ranges
over an interval of length $O(X)$ in its respective packet.
For each exponent $a\ge2$ the number of powers in such a range
at scale at least $X$ is $O(\sqrt X+1)$ by the mean value theorem
applied to $t^{1/a}$. There are $O(L)$ possible exponents and the
weights are $O(L^2)$. This power contribution can be absorbed
in $O_\delta(X\SSS(bN))$. Summing over lines and applying
\cref{eq:packetpoint,eq:subseries} gives a normalized coefficient
$O_\delta(\sqrt{XY}/N)$.

If $X<N^\delta$ and $Y\ge N^{1+\delta}$, fix the smaller integer
and sieve the larger one in either of its two reduced classes
modulo $N$. The progression parameter ranges over a length
$H\asymp Y/N$. One residue is excluded at each prime not dividing
$N$, and none at primes dividing $N$. The prime count in each
class is
$O((Y/N)(N/\varphi(N))/\log(Y/N))$ by \cref{eq:uniformsieve}.
After inserting logarithmic weights it is $O_\delta(Y/\varphi(N))$.
The proper prime powers contribute $O_\delta(Y/N)$ in each class.
For any fixed exponent $a$, the congruence $t^a=u\bmod N$ with
$u$ a unit has at most $2a^{\omega(N)}$ roots modulo $N$.
This follows from the unit groups at prime powers and the Chinese
remainder theorem. For fixed $a$ this is $O_{\delta,a}(N^{\delta/2})$,
seen by separating the primes with $p^{\delta/2}<a$.
As $Y\ll N^2$, the possible bases occupy only boundedly many
complete residue intervals modulo $N$. For exponents $a\ge A$
with $A$ a sufficiently large constant depending on $\delta$,
the total unweighted count even without the residue restriction
is $O(Y^{1/A}L)$. With logarithmic weights this is $O_\delta(Y/N)$
by $Y\ge N^{1+\delta}$. The finitely many exponents below $A$
give the same conclusion from the root bound. Now use
$\sum_{x\le eX}\Lambda(x)\ll X$ and \cref{eq:subseries,eq:packetpoint}
to get the normalized coefficient $O_\delta(\sqrt{XY}/N)$.

In the remaining case $X<N^\delta$, $N/(2e)\le Y<N^{1+\delta}$,
the progression has $O(Y/N)$ possible larger integers for
each smaller one. With the trivial $\Lambda(y)\ll L$ this gives
a normalized coefficient $O(L\sqrt{XY}/N)$.

\subsection{Summation of the deleted contributions}

For the first two cases and the model, the coefficient
$\sqrt{XY}/N=e^{(j+k-2L)/2}$ is summed only over
$j+k\le 2L-\log Q+2$. Each row and each column has geometric
sum $O(Q^{-1/2})$. For the remaining arithmetic case, sum
$L e^{(j+k-2L)/2}$ over $j+k\le(1+2\delta)L$ to get row and
column sums $O(LN^{-1/2+\delta})$. The same bounds hold with
$j,k$ interchanged. The Schur bound and the bounded overlap
of log packets give the desired form norm on the deleted
region with $x,y\ge1$.

For $x<1$ the arithmetic sum is empty. To bound the model there,
we use a stronger estimate valid whenever the smaller variable
is at most $A$ with $1\le A\le N/4$. For $x\le A$, only the
terms $z(bN+x)$ and $z(bN-x)$ occur, with opposite signs. With
$I_b=[bN-A,bN+A]$, Cauchy--Schwarz gives
\[
 |z(bN+x)-z(bN-x)|
   \le\frac{\sqrt{2x}}{bN-A}
       \left(\int_{I_b}y^2|z'(y)|^2\,\ud y\right)^{1/2}.
\]
The intervals $I_b$ are disjoint. Their derivative integrals sum
to $O(\NN(G)^2)$ since $y\ge3N/4$, where $W_N(\log y)\gg1$.
Also $\sum_{b\ge1}(\SSS(bN)/\SSS(N))^2/b^2\ll1$
by the convergent Euler product from \cref{eq:series}.
Summing over $b$ and integrating the low test gives
\[
 \int_0^A |w(x)|\sqrt{x}\,\ud x
  =\int_{-\infty}^{\log A}|F(r)|e^r\,\ud r
  \ll A\NN(F),
\]
using Cauchy--Schwarz and
$\int_{-\infty}^{\log A}e^{2r}W_N(r)^{-1}\,\ud r\ll A^2$ for $A\ge1$.
Applying symmetry bounds the full model restricted to
$\min(x,y)\le A$ in normalized operator norm by $O(A/N)$.
In particular $A=1$ gives $O(N^{-1})$ for the deleted model
terms below one.
Combining these bounds and absorbing $L N^{\delta-\varepsilon}$
proves \cref{eq:subcritical} on smooth tests. Density from
\cref{sec:cutoff} gives the operator statement, including the
bounded extension of the retained form.

In the finite matrix of \cref{thm:mesh}, one may therefore
apply this pairwise cutoff directly in every entry of $B$.
The largest absolute generalized eigenvalue changes by at most
the right-hand side of \cref{eq:subcritical}. For example,
$Q=(\log\log N)^2$ retains the error scale
$O(1/\log\log N)$ of the matrix criterion.
One may also omit the continuous model on $\min(x,y)\le A$
from this matrix at additional normalized cost
$O(A/N+Q^{-1/2})$. In fact the full model on this region costs
$O(A/N)$, and its part with $x,y\ge1$ below the product threshold
has the absolute bound already proved for continuous lines.

\subsection{Sharp product cutoffs and signed rectangle counts}

Fix $C\ge1$. On $x,y\ge1$, write $\mu_N$ for the signed measure
of the physical form in \cref{eq:form}, divided by $\SSS(N)$.
The test weights in its integral are $w(x)\overline{z(y)}$ in
the physical convention of \cref{sec:cutoff}.
Write $\mu_{N,C}=\mathbf1_{xy\le CN^2}\mu_N$,
and let $U_N(C)$ be its form norm for the energy $\NN$.
For $X,Y\ge1$, take intervals $I\subset [X/e,eX]\cap[1,\infty)$
and $J\subset[Y/e,eY]\cap[1,\infty)$, allowing all choices of
endpoint inclusion. Define
\[
 d_N(C)=\sup_{X,Y,I,J}\frac{|\mu_{N,C}(I\times J)|}{\sqrt{XY}},
 \qquad
 \omega(d)=d(1+\log_+(1/d)),\qquad \omega(0)=0 .
\]

\begin{theorem}[A rectangle criterion through the critical product]
\label{thm:criticaldiscrepancy}
For each fixed $C\ge1$ there is $c>0$ such that, for all sufficiently
large even $N$,
\begin{align}
 U_N(C)&\ll_C \omega(d_N(C))+N^{-c}, \notag\\
 d_N(C)&\ll_C U_N(C)^{1/4}+N^{-c}.
 \label{eq:criticaldiscrepancy}
\end{align}
These two estimates hold for the physical form without GRH.
Under $\grh(N)$ one also has
\begin{equation}
 U_N(C)\ll_C \left(\frac{\|T_N\|}{\SSS(N)}\right)^{1/3}
                    +N^{-c}.
 \label{eq:criticalfullrestriction}
\end{equation}
Thus hypothesis S implies decay of the signed rectangle
discrepancy for every $C$. Conversely, discrepancy decay for a
given $C$ suffices for decay of the sharp product restriction
through $CN^2$.
\end{theorem}

We first give two local bounds for the absolute arithmetic and
continuous masses. In these bounds the packet intervals at scales
$1\le X\le Y$ may be enlarged by a fixed multiplicative constant.
We consider only packet pairs meeting the product region and at
least one of the congruence lines, so $XY\ll_C N^2$ and $Y\gg N$.
Write $L=\log N$, fix $\delta=1/10$, and initially assume
$N^{-1/100}\le\eta\le1/10$. All packet mass bounds below are
for the singular-series-normalized measure and include a
further factor $(XY)^{-1/2}$.

First, the upper collar $CN^2<xy\le e^\eta CN^2$ has absolute
mass $O_C(\eta+N^{-c_1})$ on each packet pair. On the two
difference lines the derivative of the product with respect
to $x$ is $x+y\gg Y$. On the reflection line $x+y=bN$ its
magnitude is $|x-y|\gg Y$ if the scales are sufficiently
separated. For comparable scales, define
\[
 \Delta_C=\min_{\substack{b\in\mathbb N\\ b^2/4>C}}(b^2/4-C)>0
\]
and first take $\eta$ small enough that
$C(e^\eta-1)\le\Delta_C/2$. A reflection line with
$b^2/4\le C$ has no points in the upper collar. On any other
reflection line there,
\[
 |x-y|^2=(bN)^2-4xy\ge 2\Delta_C N^2.
\]
In this range $X\asymp_C Y\asymp_C N$. Splitting at $x=y$,
each interval in the inverse image of the collar has length
$O_C(\eta X)$ by the derivative bound. Thus the total length
per line is $O_C(\eta X)$ in each orientation.
This proves the continuous bound using
\cref{eq:subseries}. If $X\ge N^\delta$, the short-interval
sieve \eqref{eq:uniformsieve} gives a weighted two-prime
mass $O_C(\eta X\SSS(bN))$ per line. The intervals
can be padded to length $\asymp_C\eta X$, whose
logarithm is comparable to $L$. The remainder is absorbed
by this bound. The proper powers have weight $O_C(\sqrt X L^3)$
on a whole line as in the proof of \cref{thm:subcritical}.
Since there are $O(Y/N)$ lines, they cost $O_C(N^{-c_1})$.
If $X<N^\delta$, the collar forces $Y\gg_C N^{2-\delta}$.
Fixing an integer in the smaller coordinate, an interval of
length $O_C(\eta Y)$ in the other coordinate contains
weighted prime mass $O_C(\eta Y/\varphi(N))$ in either
reduced residue class. Apply \eqref{eq:uniformsieve} to the
progression parameter on an interval padded to length
$\asymp_C\eta Y/N$. The proper powers in each class have
weighted mass $O_C(N^{-c_1}Y/N)$. This follows from the
root bound in the proof of \cref{thm:subcritical}, which
gives $O_C(N^{\delta/2})$ solutions for each fixed exponent.
For all sufficiently large exponents combined, the unrestricted
power count at $Y\ll_C N^2$ is $O_C(Y^{1/k_0} L)$ with a
fixed $k_0$ large enough that $2/k_0<\delta/2$.
These estimates for a fixed row in fact hold whenever
$X<N^\delta$ and $Y\ge N^{1+\delta}$. Summing the smaller
weights by $\sum_{n\le e^2 X}\Lambda(n)\ll X$ proves the
collar bound in this case as well.
Only $O_C(1)$ unit log packets per row and column meet the
collar. Thus its form has normalized norm
\begin{equation}
 \ll_C \eta+N^{-c_1}
 \label{eq:criticalcollar}
\end{equation}
by point evaluation and the Schur bound. For larger $\eta$
up to $1/10$, the absolute packet bounds from
\cref{sec:subcritical}, with cap $e^{1/10}CN^2$, give a
bounded collar norm, which is $O_C(\eta)$ in that range.
For $0<\eta<N^{-1/100}$ use the same bound with the collar
enlarged to width $N^{-1/100}$, and decrease $c_1$.

The bound is for an upper collar at a fixed $C$. If
$C=b^2/4$, the reflection line $x+y=bN$ meets the lower
collar $e^{-\eta}CN^2<xy\le CN^2$ in an interval of length
$2N\sqrt{C(1-e^{-\eta})}$. Its unsigned continuous
contribution on comparable packets can therefore be of
order $\sqrt\eta$.

Second, in a packet box clipped by $xy\le CN^2$, a strip
of width $O(\eta Y)$ in the larger coordinate has absolute
mass $O_C(\sqrt\eta\log(2/\eta)+N^{-c_1})$.
A strip of width $O(\eta X)$ in the smaller coordinate has
the same bound if $X\ge\eta^{-1}$. To see this for
$X\ge N^\delta$, use the short-interval sieve as above.
A smaller-coordinate strip takes length $O(\eta X)$ on
each line. For a larger-coordinate strip, the length is at
most $O(\min(X,\eta Y))$, and the shift parameters lie in
a bounded number of intervals of length $O((X+\eta Y)/N)$.
On any interval of shift parameters of length $M$ and upper
endpoint $O(Y/N)$, the divisor expansion in
\cref{sec:cutoff} gives
\[
 \sum_b\frac{\SSS(bN)}{\SSS(N)}
   \ll M+\log(2+Y/N).
\]
For each divisor $d$ there are at most $M/d+1$ multiples on
the interval, and $\sum_{d\le u}a_N(d)\ll1+\log u$.
With $r=X/Y$ and $B=Y/N$ the resulting larger-strip bound
for the prime and continuous masses is
\[
 O_C\!\left(\frac{\min(r,\eta)}{\sqrt r}
                 \big((r+\eta)B+\log(2+B)\big)\right).
\]
Here $rB^2\ll_C1$. The first term is $O_C(\eta)$ because
$\min(r,\eta)(r+\eta)\le2r\eta$.
For the second use $\min(r,\eta)/\sqrt r\le
\min(\sqrt\eta,O_C(B^{-1}))$ and split at $B=\eta^{-1/2}$.
This gives $O_C(\sqrt\eta\log(2/\eta))$.
The proper powers on full lines still have a power saving.
The same continuous bounds apply for $X<N^\delta$.

For the arithmetic strips with $X<N^\delta$ and
$Y\ge N^{1+\delta}$, use the row estimates from the collar
argument. For larger-coordinate strips they give the
desired bound on summing the smaller weights. For
smaller-coordinate strips, take an interval $I'$ of
length $O(\eta X)$ in the smaller packet. If
$\eta^{-1}\le X<\eta^{-2}$, it contains $O(\eta X)$ integers,
each with weight $O(\log(2/\eta))$. If $X\ge\eta^{-2}$,
apply the one-prime sieve on an interval padded to length
$\asymp\eta X\ge\sqrt X$.
Proper powers in this case contribute
$O((\eta\sqrt X+1)(\log(2X))^2)$ by counting roots in
the interval for each exponent. Thus
$\sum_{n\in I'}\Lambda(n)\ll\eta X\log(2/\eta)$
in both cases. Use this bound for the smaller weights and
the whole larger packet row bound to obtain
$O_C(\eta\log(2/\eta)+N^{-c_1})$. The remaining regime
$X<N^\delta$, $Y<N^{1+\delta}$ has whole packet arithmetic
mass $O_C(L\sqrt{XY}/N)=O_C(N^{-c_1})$ as in
\cref{sec:subcritical}. This proves the strip estimates.

For the first inequality of \cref{eq:criticaldiscrepancy},
use a partition into unit log packets with bounded overlap.
On a packet at scale $X$, the physical test has supremum
plus total variation $O(X^{-1/2}\NN(F_j))$. On its packet
interval $[a,b]$ the test equals $w(b)-\int_x^b w'(t)\,\ud t$.
Insert this representation in each variable and integrate
against the measure. Every cumulative mass in the expansion
is a rectangle mass, giving a packet bound
$O(d_N(C)\NN(F_j)\NN(G_k))$.
The absolute estimates of \cref{sec:subcritical} also give
a coefficient $O_C(\sqrt{XY}/N+N^{-c_1})$ on these packet
pairs. The same proof applies with the fixed factor $C$
in the cap, since still $Y\ll_C N^2$.
In the minimum of the two coefficients put
$X=e^j$, $Y=e^k$ and sum over $j+k\le 2L+O_C(1)$.
The row and column sums of $\min(d_N(C),e^{(j+k)/2}/N)$
are $O_C(\omega(d_N(C)))$ by a geometric sum below
$e^{(j+k)/2}/N=d_N(C)$ and counting the indices above it.
Only $O_C(L)$ indices can meet the cap with $x,y\ge1$,
so the $N^{-c_1}$ term still gives a power saving.
The Schur bound proves the inequality.

For the reverse inequality take $X\le Y$ by symmetry.
Approximate each interval indicator by a smooth function
between zero and one, with transition of relative width
$\eta$. Use these smooth functions divided by $\sqrt X$
and $\sqrt Y$ as physical test weights.
Their energy norms are $O_C(\eta^{-1/2})$ by
\cref{eq:logenergy}. In fact the log supports have bounded
length and lie below $2L+O_C(1)$, and the two ramps have
total log length $O(\eta)$ and derivative $O(\eta^{-1})$.
If $X\ge\eta^{-1}$, the strip bounds control the absolute
counting error for each endpoint.

If $X<\eta^{-1}$, replace the smaller indicator by a
smooth function with the same values on the integers.
Take a plateau from the first included integer to the last,
with ramps of width $1/4$, or the zero function if none
are included. The energy norm now costs $O_C(\sqrt X)
\ll_C\eta^{-1/2}$, and the larger-coordinate strip still
controls the arithmetic error.
For the continuous signed measure in this case, fix $x$
in the smaller support. Since $x\ll\eta^{-1}\ll N$,
the continuous terms pair the larger weights at $bN+x$
and $bN-x$. With the cap, the larger weight is restricted
to $y\le CN^2/x$. Either the larger indicator or its
smooth approximation has bounded total variation after
that restriction. The intervals $[bN-x,bN+x]$ are disjoint,
so the sum of the absolute differences of these larger
weights before the factor $Y^{-1/2}$ is $O(1)$.
Use $\SSS(bN)/\SSS(N)\ll(bN)^{1/4}$, as follows
from \cref{eq:series}, and $bN\asymp Y$ on the possible
lines. The normalized signed continuous mass for either
pair of tests is therefore $O_C(\sqrt X\,Y^{-1/4})$.
This is $O_C(N^{-c_1})$ since $Y\gg N$ and
$\eta\ge N^{-1/100}$. We conclude in both cases that
\[
 d_N(C)\ll_C
    \eta^{-1}U_N(C)+\sqrt\eta\log(2/\eta)+N^{-c_1}.
\]
Use $\eta=\max(U_N(C)^{2/3},N^{-1/100})$ when it is
at most $1/10$. The absolute packet bounds give
$d_N(C),U_N(C)\ll_C1$, which handles larger values of the
norm. This gives the second inequality of
\cref{eq:criticaldiscrepancy}, after absorbing logarithms
and decreasing $c$.

To prove \cref{eq:criticalfullrestriction}, let
$E=\|T_N\|/\SSS(N)$, $Q\ge2$, and
$v=\log(xy/N^2)$. Use a smooth mask of $v$ equal to one on
$[-\log Q,\log C]$, zero off $[-\log Q-1,\log C+\eta]$,
with the upper transition at scale $\eta$ and the lower
one of fixed width. Also multiply each argument by a
fixed smooth cutoff equal to one for $x\ge1$ and zero
for $x\le1/2$. Write $M$ for the mask. Its first derivative
has $L^1$ norm $O(1)$ and its fourth derivative has $L^1$
norm $O(\eta^{-3})$. Its Fourier transform therefore
satisfies
\[
 |\widehat M(\xi)|\ll_C
      \min(1+\log Q,|\xi|^{-1},\eta^{-3}|\xi|^{-4}).
\]
The two phase multipliers in Fourier inversion incur a joint
energy cost $O((1+|\xi|)^2)$. Integrate the displayed bound
with this weight, splitting at $|\xi|=1$, $\eta^{-1}$ and,
below one, at $(1+\log Q)^{-1}$. The resulting normalized
form norm is $O_C((\eta^{-2}+\log(2+\log Q))E)$.
On $x,y\ge1$ the difference from the sharp form is
bounded by \cref{thm:subcritical} at the lower transition
and by \cref{eq:criticalcollar} at the upper one.
On $x<1$ or $y<1$ only the continuous terms occur.
The estimate below one in \cref{sec:subcritical} continues
to apply with cost $O_C((1+\eta^{-1})/N)$.
For fixed $x<1$, the derivative in the larger variable
of the masked larger weight is bounded by a constant
times $|z'(y)|+(1+\eta^{-1})|z(y)|/y$, and the same
disjoint-interval argument applies. Thus
\[
 U_N(C)\ll_C (\eta^{-2}+\log(2+\log Q))E
       +Q^{-1/2}+\eta+N^{-c_1}.
\]
Choose $\eta=\max(E^{1/3},N^{-1/100})$ and $Q=\eta^{-2}$
when $\eta\le1/10$, using the uniform bound for $U_N(C)$
otherwise. This proves \cref{eq:criticalfullrestriction}.

\subsection{Prefix rectangles at the critical product}
\label{sec:criticalgrid}

Let $B_j=[2^j,2^{j+1})$ for $j\ge0$. Fix $C\ge1$, take
$Q\ge2$ and an integer $10\le M\le N^{1/100}$, and keep
only packet pairs at scales $X=2^j$, $Y=2^k$ with
\begin{equation}
       N^2/Q\le XY\le CN^2 .
 \label{eq:criticalgridband}
\end{equation}
For each packet with $X\ge M$, use the grid points
$X(1+i/M)$ for $0\le i\le M$. For $X<M$, use $X$, $2X$,
and the half-integers strictly between them. For a grid
point $u$ at scale $X$ write $P_X(u)=[X,u)$.
Define $g_N(C;Q,M)$ as the maximum over the kept pairs
and their grid points $u,v$ of
\[
   \frac{|\mu_{N,C}(P_X(u)\times P_Y(v))|}{\sqrt{XY}},
\]
with the maximum of an empty family set to zero.

\begin{theorem}[A finite prefix test for the critical restriction]
\label{thm:criticalgrid}
Uniformly in $Q,M$ in the indicated ranges, for sufficiently
large even $N$,
\begin{equation}
  g_N(C;Q,M)\le d_N(C)
   \ll_C g_N(C;Q,M)+Q^{-1/2}
       + M^{-1/2}\log(2M)+N^{-c}
  \label{eq:criticalgridcomparison}
\end{equation}
for some $c>0$. Thus $U_N(C)\to0$ is equivalent to
$g_N(C;Q,M)\to0$ for each fixed $Q,M$. The test uses
$O_C(M^2\log(2Q)\log(2N))$ prefix rectangles.

For $K\ge8\sqrt C$, one can use only the signed arithmetic
component on the pairs with $\max(X,Y)/\min(X,Y)\ge K^2$.
This changes the maximum by
$O_C(K^{-1}\log(2K))$. For fixed $C,Q,M,K$ the resulting
test has at most $O_{C,Q,M,K}(1)$ entries requiring the
continuous model.
\end{theorem}

First consider $X\le Y$ and $X\le N/8$. For intervals
$I\subset B_j$ and $J\subset B_k$, the continuous part of
the signed rectangle discrepancy satisfies
\begin{equation}
 \frac{|\mu_{N,C}^{\rm cont}(I\times J)|}{\sqrt{XY}}
      \ll \sqrt{X/Y}\,\log(2+Y/N).
 \label{eq:criticalgridmodel}
\end{equation}
For fixed $x\in I$, set $q_x(y)=\mathbf1_J(y)
\mathbf1_{xy\le CN^2}$. It has bounded total variation,
including any specified values at the endpoints.
The model lines contribute the difference
$q_x(bN+x)-q_x(bN-x)$. The corresponding intervals
$[bN-x,bN+x]$ are disjoint since $x<2X\le N/4$.
Hence the sum of the absolute differences before inserting
the singular-series ratios is $O(1)$. Only $b\ll Y/N$
can contribute. The divisor expansion in \cref{sec:cutoff}
gives
\[
  \frac{\SSS(bN)}{\SSS(N)}
       =\sum_{d\mid b}a_N(d)
       \ll 1+\log(2b).
\]
Integrating over $I$ of length at most $X$ and dividing
by $\sqrt{XY}$ proves \cref{eq:criticalgridmodel}. On
a kept pair with $Y/X\ge K^2$, put $s=\sqrt{X/Y}$ and
$B=Y/N$. Then $s\le K^{-1}$, $sB\le\sqrt C$, and
$X\le\sqrt C\,N/K\le N/8$. The bound
$s\log(2+B)\ll_C K^{-1}\log(2K)$ follows by splitting
at $B=K$. Symmetry handles $X>Y$.

Every prefix rectangle is admissible in $d_N(C)$. For
the reverse bound, each interval in its definition splits
among a bounded number of dyadic packets at comparable
scales. The normalized discrepancy is therefore bounded
by a constant times the supremum for rectangles in individual
dyadic pairs. If $XY>CN^2$, their measure is zero.
For $XY<N^2/Q$, the absolute packet estimates from
\cref{sec:subcritical} give
$O_C(\sqrt{XY}/N+N^{-c})=O_C(Q^{-1/2}+N^{-c})$.

On a kept pair with $M\le X\le Y$, round each interval
to a consecutive union of its grid cells, allowing the
empty interval. The change is confined to strips at
the two endpoints, of widths at most $2X/M$ and $2Y/M$
in their respective coordinates. The strip estimates in
\cref{sec:subcritical} apply with $\eta=M^{-1}$ and bound
the total change, on the normalized scale, by
$O_C(M^{-1/2}\log(2M)+N^{-c})$.
If $X<M\le Y$, choose a consecutive union of the
half-integer cells at scale $X$ with exactly the same
integer points as the smaller interval. Each cell contains
at most one such point. Thus the arithmetic change comes
only from the strips at scale $Y$. The two continuous
rectangle masses are individually $O_C(N^{-c})$ by
\cref{eq:criticalgridmodel}, since $X<M\le N^{1/100}$
and any contributing pair has $Y\gg N$. When both scales
are below $M$, no congruence line meets the pair for
sufficiently large $N$.

A rectangle formed by two consecutive unions of cells
has measure equal to a signed sum of four prefix
measures, so its normalized absolute measure is at
most $4g_N(C;Q,M)$. This also applies to empty intervals
and to the first or last cells of each packet. The
estimates above prove \cref{eq:criticalgridcomparison}.
For the asserted equivalence, take the upper limit as
$N\to\infty$ at fixed $Q,M$ and then send both to infinity
in \cref{eq:criticalgridcomparison}, using
\cref{thm:criticaldiscrepancy}.

The interval of possible $j+k$ in
\cref{eq:criticalgridband} has length $O_C(\log(2Q))$,
and each value gives $O_C(\log(2N))$ pairs. Each packet
has at most $M+1$ grid points. In the modified test
the model is needed only for $|j-k|\le 2\log_2 K$.
Together with the bounded interval of possible sums
at fixed $C,Q$, this places each of $j,k$ in an
interval of bounded length at fixed $C,Q,K$.
This proves both counts.

\section{Dirichlet-zero covariance when both variables are large}
\label{sec:zeros}

With the inducing conductors $q_\chi$ of \cref{sec:cutoff}, set
\[
 b_N=\frac{2}{\varphi(N)}\sum_{\chi\ {\rm odd}}\log q_\chi.
\]
Use $\psi_\Gamma=\Gamma'/\Gamma$ for the digamma function.

\begin{theorem}[High-support covariance formula]
\label{thm:zeros}
Fix $\delta>0$. Under $\grh(N)$, there is $c_\delta>0$
such that for compactly supported tests $F,G$ in
$r\ge(1+\delta)\log N$, the following error bound is uniform in
the supports:
\begin{align}
 \B(f,g)
 &=\frac{2}{\varphi(N)}\sum_{\chi\ {\rm odd}}
       Z_\chi(F)\overline{Z_\chi(G)}
    -\mathcal W_N(F,G)
    +O_\delta\bigl(N^{-c_\delta}\NN(F)\NN(G)\bigr),
 \label{eq:zeroform}\\
 \mathcal W_N(F,G)
 &=\frac1{2\pi}\int_{\mathbb R}\widehat F(t)\overline{\widehat G(t)}
     \left(b_N-\log\pi+\Re\psi_\Gamma\!\left(\tfrac34+\tfrac{it}{2}\right)\right)
     \ud t .
 \label{eq:archdensity}
\end{align}
Write $\mathcal O_N(F,G)$ for the average over distinct zero labels
\[
 \frac{2}{\varphi(N)}\sum_{\chi\ {\rm odd}}
     \sum_{\substack{\gamma,\gamma'\in\mathcal Z_\chi\\
                     \text{distinct labels}}}
         \widehat F(\gamma)\overline{\widehat G(\gamma')}.
\]
In a multiple zero with equal ordinates the distinct labels are
included in this sum.
If $|r-s|\le(2-\eta)\log N$ whenever $r,s$ are in the two
log supports, where $0<\eta<2$ is fixed, then
$\B(f,g)=\mathcal O_N(F,G)+O_{\delta,\eta}
(N^{-c}\NN(F)\NN(G))$ for some $c>0$.
For arbitrary support separation the corresponding correction
is a signed progression sum given in \cref{eq:samelabelcorrection}.
\end{theorem}

\subsection{The conductor and gamma terms from the prime model}

Write $L=\log N$ and use log packets with index $j\le k$ as in
\cref{sec:cutoff}. Set $X=e^j$, $R=X/N$ and $D=e^{k-j}$.
Rescale their two physical weights by $x=Xu$, absorbing $\sqrt X$
into each weight. Denote the resulting weights by $w,z$, and write
$I=\int w\overline z\,\ud u=\int F_j\overline{G_k}\,\ud r$.
In these coordinates the continuous model equals
$\sum_{a\ge1}\SSS(aN)J(a/R)$, with $J$ as in \cref{eq:model}
using $w,z$. Thus $J$ is supported in $[0,2eD]$, has
$J(0)=2I$, and $\int_0^\infty J(t)\,\ud t=0$.

The divisor expansion of $\SSS(aN)$ gives the Dirichlet series
\[
 \sum_{a\ge1}\frac{\SSS(aN)}{a^s}
 =\SSS(N)\zeta(s)\zeta(1+s)H_N(s),\qquad
 H_N(s)=\prod_p(1-p^{-1-s})(1+a_N(p)p^{-s}).
\]
The Euler product for $H_N(s)$ converges absolutely for
$\Re s>-1/2$. At primes not dividing $N$ the linear coefficient
after multiplying the two local factors is $a_N(p)-1/p=O(p^{-2})$.
The primes dividing $N$ satisfy the uniform bound
$\SSS(N)|H_N(s)|\ll_\nu N^\nu$ on $\Re s\ge-1/8$ for each $\nu>0$.
Indeed their factors have absolute values at most $1+p^{-7/8}$,
and we may separate the finitely many primes with
$1+p^{-7/8}>p^{\nu/2}$. The remaining factor
$\SSS(N)\ll_\nu N^{\nu/2}$ follows in the same way.
Local multiplication and differentiation give
\begin{equation}
 \SSS(N)H_N(0)=1,\qquad
 c_N:=\frac{H_N'(0)}{H_N(0)}
       =\sum_{p\mid N}\frac{\log p}{p-1}.
 \label{eq:condconstant}
\end{equation}

Write $\widetilde J(s)=\int_0^\infty J(t)t^{s-1}\,\ud t$ in
$\Re s>0$. Its Laurent expansion at zero is $2I/s+J_*+O(s)$, where
\[
 J_*=\int_0^\infty\frac{J(t)-2I\mathbf1_{t<1}}{t}\,\ud t.
\]
Two integrations by parts give
$s(s+1)\widetilde J(s)=\int J''(t)t^{s+1}\,\ud t$.
The latter extends to $\Re s\ge-1/8$.
The $L^\infty$ norm of $J''$ is
$O(\|F_j\|_{H^1}\|G_k\|_{H^1})$ by putting one derivative on
each factor in each convolution. The rescaled $w,z$ have the
required $H^1$ bounds uniformly for $D\ge1$. On the new vertical
line this gives
$|\widetilde J(-1/8+it)|\ll D^2(1+|t|)^{-2}
\|F_j\|_{H^1}\|G_k\|_{H^1}$.
Mellin inversion and a shift to $\Re s=-1/8$ give
\begin{equation}
 \mathcal M(F_j,G_k)
   =-I\bigl(\log R+\gamma_{\mathrm E}+\log(2\pi)+c_N\bigr)
      -\tfrac12 J_*
      +O_\nu\bigl(N^\nu D^2 R^{-1/8}
                   \|F_j\|_{H^1}\|G_k\|_{H^1}\bigr).
 \label{eq:modelresidue}
\end{equation}
Here $\gamma_{\mathrm E}$ is Euler's constant. The pole at one of
$\zeta(s)$ is cancelled by $\widetilde J(1)=0$. At zero use
$\zeta(0)=-1/2$, $\zeta'(0)=-\tfrac12\log(2\pi)$, and the
Laurent series of $\zeta(1+s)$. The shifted integral converges
absolutely, since the functional equation and a truncated
Dirichlet series give
$|\zeta(-1/8+it)\zeta(7/8+it)|\ll (1+|t|)^{3/4}$.
The twice integrated Mellin transform also makes the horizontal
contours tend to zero. A related singular-series residue for a
variance in the prime-modulus family appears in~\cite[Section~8]{TZLiu}.

The diagonal satisfies, on overlapping packets,
\[
 \mathcal D(F_j,G_k)
   =\int r F_j(r)\overline{G_k(r)}\,\ud r
       +O\bigl(X^{-1/2}\log^3(2NX)\|F_j\|_{H^1}\|G_k\|_{H^1}\bigr).
\]
Indeed the Riemann hypothesis for $\zeta$, included in $\grh(N)$,
gives the main term $\int_2^t\log u\,\ud u$ for
$\sum_{n\le t,(n,N)=1}\Lambda(n)^2$ with an
$O(\sqrt t\log^3(2Nt))$ error. This follows by partial summation
of $\sum_{n\le t}\Lambda(n)$, removing the difference between
$\Lambda(n)\log n$ and $\Lambda(n)^2$ at proper powers. The
terms above primes dividing $N$ have weight $O(\log N\log t)$.
Another partial summation against $F_j(\log x)\overline{G_k(\log x)}/x$
gives the displayed packet bound. On disjoint packets both diagonal
terms are zero.

To compute $J_*$, change variables $u=e^v$, $y=e^z$ in the
double integral for the two rescaled physical variables. The
difference term uses the kernel $1/(2\sinh(|v-z|/2))$ and the
reflected term uses $1/(2\cosh((v-z)/2))$.
Cutting out $|u-y|\le a$ corresponds locally to
$|v-z|\le a e^{-v}$, up to $O(a^2)$ on each fixed packet.
If
$K(t)=\operatorname{fp}(1/(2\sinh(|t|/2)))-1/(2\cosh(t/2))$,
with finite part defined by adding $2\log a$ times evaluation at
zero to the integral over $|t|>a$, this gives
\[
 J_*=2\int (r-\log X)F_j(r)\overline{G_k(r)}\,\ud r
      +\iint F_j(r)\overline{G_k(s)}\,K(r-s)\,\ud r\,\ud s .
\]
The local cutoff change is justified by the $1/2$-H\"older
bound for $H^1$ tests. The Fourier symbol of $K$ is
\[
 \widehat K(t)=-2\Re\psi_\Gamma(\tfrac12+it)
                 -2\gamma_{\mathrm E}-\frac{\pi}{\cosh(\pi t)} .
\]
For the first term subtract the integral at $\psi_\Gamma(1)$
from the digamma integral. The transform of the second kernel
is the beta integral $\mathrm B(1/2+it,1/2-it)$.
Digamma reflection and duplication \cite[Sections~5.5 and~5.9]{DLMF}
give
\[
 \Re\psi_\Gamma(\tfrac12+it)+\frac{\pi}{2\cosh(\pi t)}
   =\log 2+\Re\psi_\Gamma(\tfrac34+\tfrac{it}{2}).
\]
Consequently \cref{eq:modelresidue} plus the arithmetic diagonal
has the symbol in \cref{eq:archdensity} with $b_N$ replaced
by $\log N-c_N$.

In the average over all characters modulo $N$, the mean of
$\log q_\chi$ is $\log N-c_N$. For each $p^\nu\Vert N$,
the mean conductor exponent lost at $p$ is the sum of the
probabilities that the conductor divides $N/p^a$ for
$1\le a\le\nu$. These probabilities sum to $1/(p-1)$.
For the average over odd characters the same probabilities
hold whenever $N/p^a\ge3$, since exactly half the characters
of any such modulus are odd. The divisors $1,2$ change the
answer by $O(L/\varphi(N))$. Thus
$b_N=\log N-c_N+O(L/\varphi(N))$.

\subsection{Uniformity of the reduction and the local off-diagonal sum}

Choose $\beta=\min(\delta,1)/100$ and $\nu=\beta$ in
\cref{eq:modelresidue}. For $|k-j|\le\beta L$ the error there is
a power saving in $N$, since $D^2\le N^{2\beta}$ and
$R\ge e^{-1}N^\delta$. On these packets
$\|F_j\|_{H^1}\ll\NN(F_j)$ and similarly for $G_k$.
The diagonal error is a power saving uniformly for all overlapping
packets at these heights. There are $O_\delta(L)$ such nearby
pairs per row or column. For $|k-j|>\beta L$ the arithmetic
diagonal is zero and the model bound in \cref{eq:Mpacket}
has summable tail $O(\SSS(N)L e^{-3\beta L/8})$. The convolution
kernel of \cref{eq:archdensity} away from its diagonal is
$-K/2$, and for $|t|\ge1$ the difference of the two hyperbolic
kernels is $O(e^{-3|t|/2})$. The far tail of the archimedean
form is therefore also a power saving. The Schur bound and
bounded overlap prove that $\mathcal D+\mathcal M$ differs from
$\mathcal W_N$ by $O_\delta(N^{-c_\delta})$ in norm.

For each $\chi$, \cref{eq:poslogexplicit} gives
$\ell_\chi(F)=-Z_\chi(F)$ up to
$O(L N^{-(1+\delta)/2}\NN(F))$ on these supports.
Indeed $\|F\|_{H^1}\ll\NN(F)$ on the positive log axis,
so the trivial-zero and imprimitive estimates following
\cref{eq:poslogexplicit} apply directly. The character norms
of $\ell_\chi(F)$ and $Z_\chi(F)$ are $O(L\NN(F))$
by \cref{eq:Utail} and the pointwise error. Replacing both
character factors in \cref{eq:split} therefore changes the
form by a power saving.
This proves \cref{eq:zeroform}.

To separate the same-label sum, put
\[
 C_{F,G}(v)=\int_{\mathbb R}F(r+v)\overline{G(r)}\,\ud r,
 \qquad S_{F,G}(v)=C_{F,G}(v)+C_{F,G}(-v).
\]
The inverse Fourier transform of the even part of
$\widehat F\overline{\widehat G}$ is $S_{F,G}/2$.
In the Weil explicit formula \cite[Section~2]{HughesRudnick},
we may use this even part since conjugation permutes the odd
characters and reflects their zero ordinates.
The archimedean part is $\mathcal W_N(F,G)$, and the
prime-power terms subtract
\[
 \sum_{n\ge2}\frac{\Lambda(n)}{\sqrt n}S_{F,G}(\log n)
      \frac{2}{\varphi(N)}\sum_{\chi\ {\rm odd}}\chi^*(n).
\]
The mean of $\chi^*(n)$ is real since the average is
closed under conjugation. For $(n,N)=1$ it equals
$\mathbf1_{n\equiv1\ (N)}-\mathbf1_{n\equiv-1\ (N)}$
by odd-character orthogonality.

For $p^\nu\Vert N$ and $n=p^k$, only characters whose conductor
is coprime to $p$ contribute to the mean. They are induced from
characters modulo $M=N/p^\nu$. The mean at $n$ is therefore
\[
 \frac{\mathbf1_{p^k\equiv1\ (M)}-\mathbf1_{p^k\equiv-1\ (M)}}
           {\varphi(p^\nu)}
\]
if $M\ge3$, and is zero otherwise. The numerator can be
nonzero only for $p^k\ge M-1$. Thus the absolute sum of the
coefficients of $S_{F,G}$ from this prime is
$O(\log p/(\varphi(p^\nu)\sqrt M))$ by the geometric series
over $k$. The sum over $p\mid N$ is $O(L/\sqrt N)$, since
$\varphi(p^\nu)\ge p^\nu/2$. Also
$\|S_{F,G}\|_\infty\le2\|F\|_2\|G\|_2\ll\NN(F)\NN(G)$
on the high supports. Combining the same-label formula
with \cref{eq:zeroform} gives, uniformly in these supports,
\begin{align}
 \B(f,g)
  &=\mathcal O_N(F,G)-\mathcal P_N(F,G)
       +O_\delta(N^{-c}\NN(F)\NN(G)), \notag\\
 \mathcal P_N(F,G)
  &=\sum_{n\ge2}\frac{\Lambda(n)}{\sqrt n}S_{F,G}(\log n)
       \bigl(\mathbf1_{n\equiv1\ (N)}-\mathbf1_{n\equiv-1\ (N)}\bigr)
       \label{eq:samelabelcorrection}
\end{align}
for some $c>0$. This invocation of the explicit formula
extends to compact $H^1$ tests by density and the local zero
count. In a common interval of length at most $\log 2$,
the correlation vanishes at every $\pm\log n$, including
the endpoints, so the same-label sum equals $\mathcal W_N$
exactly.

If the two log supports have all differences bounded in
absolute value by $\log X$, only $n\le X$ can occur in
$\mathcal P_N$. The relevant integers are $bN\pm1$ with
$b\ge1$, so the sum of $\Lambda(n)/\sqrt n$ over them
is $O(\sqrt X\,\log(2X)/N)$ for $X\ge N-1$ and is zero for
$X<N-1$. This gives a power saving at $X=N^{2-\eta}$ and
proves the bounded-separation assertion of \cref{thm:zeros}.

For instance, when $N=2p$ with $p$ an odd prime, all the odd
characters in the formula are induced from primitive characters
modulo $p$. For $F(r)=b(r-B\log p)$ and $G=F$, with
$b\in C_c^\infty(0,\log 2)$ and fixed $B>1$, hypothesis S would
give an additive $o(1)$ remainder in the formula that equates
the averaged squared zero sum to \cref{eq:archdensity}.
Here $\SSS(2p)\asymp1$ and $\NN(F)$ stays bounded for fixed $B$.
This specifies the precision needed for an application of
fixed-modulus pair correlation beyond the leading term of
order $\log p$.

\subsection{A finite zero-ordinate cutoff}

Fix $\delta>0$, keep $L=\log N$, and take $3\le\tau\le L$.
On the support region of \cref{thm:zeros}, let $E_\infty$
and $E_T$ be the normalized norms of the covariance form
$\langle Z(F),Z(G)\rangle_\chi-\mathcal W_N(F,G)$ with,
respectively, all ordinates and only $|\gamma|\le T$ in each
character factor. For $T=L^3\tau^2$ they satisfy
\begin{equation}
 \max(E_\infty,E_T)\le (1+O((L\tau)^{-1}))\min(E_\infty,E_T)
                               +O(\tau^{-1}).
 \label{eq:zeronormcut}
\end{equation}
To see this, \cref{eq:globalzeroshigh} bounds the tail character
norm for unit tests by $O(L^{-1/2}\tau^{-1})$. The archimedean
form has norm $O(L)$ by \cref{eq:archdensity} and Plancherel in
$H^1$. For either of the two covariance forms, positivity of
the character square bounds the retained character vector
for a unit test by $O(\sqrt{\SSS(N)(L+E)})$, where $E$ denotes
its normalized form norm. Express the other form by adding
or subtracting the tail vector in each argument. The mixed
terms have normalized norm $O(\tau^{-1}\sqrt{1+E/L})$, and
the term with two tail vectors has norm $O(L^{-1}\tau^{-2})$.
Take the smaller of $E_\infty,E_T$ and use
$\sqrt{1+v}\le1+v/2$ to get \cref{eq:zeronormcut}.
Thus $\tau=\log L$ gives an equivalent high-support smallness
criterion using only ordinates up to $L^3(\log L)^2$, by
\cref{eq:zeroform}.

\subsection{A finite zero-data criterion for the full norm}

The full operator can also be represented using low zero
ordinates, with the correction at primes dividing the
modulus retained. This representation uses the signs and
phases of the two test functions in the zero covariance.

\begin{theorem}[The full norm from finite zero data]
\label{thm:fullzerodata}
Assume $\grh(N)$ and take $3\le\tau\le L=\log N$. There is a
self-adjoint matrix with normalized norm $E_N^{\mathrm z}(\tau)$
and dimension $O(L^2\tau)$ such that
\begin{equation}
 \max\!\left(E_N^{\mathrm z}(\tau),\frac{\|T_N\|}{\SSS(N)}\right)
 \le (1+O(L^{-1}))
    \min\!\left(E_N^{\mathrm z}(\tau),\frac{\|T_N\|}{\SSS(N)}\right)
       +O(\tau^{-1}).
 \label{eq:fullzerodata}
\end{equation}
The data required are the counts of zero labels with
$|\gamma|\le L^3\tau^2$ in bins of width $L^{-5/2}/\tau$
for the inducing primitive odd characters, prime powers up
to $e L^{12}$ and the factorization of $N$. The remaining terms
are explicit integrals involving the zeta function and
a finite Euler product. The matrix is formed by compressing
two initial matrices on the mesh of \cref{thm:mesh}. With
$\tau=\log L$, its norm tending to zero is equivalent to
\cref{eq:goal}.
\end{theorem}

Fix a smooth multiplier $\theta$ equal to zero for
$r\le(\log 2)/2$ and one for $r\ge\log 2$. Then
$\ell_\chi(\theta F)=\ell_\chi(F)$. Write
$P=L^{12}$ and $T=L^3\tau^2$ and define
\begin{align}
 \ell_{\chi}^{\mathrm z}(F)
  ={}&-\sum_{|\gamma|\le T}\widehat{\theta F}(\gamma)
       -\int_0^\infty \frac{\theta(r)F(r)e^{-3r/2}}{1-e^{-2r}}
                    \,\ud r \notag\\
     &-\sum_{\substack{p\mid N,\ k\ge1\\p^k\le P}}
       \frac{\log p\,\chi^*(p^k)}{p^{k/2}} F(\log p^k).
 \label{eq:fullzeroell}
\end{align}
The explicit formula \eqref{eq:poslogexplicit} applies
with this positive-support cutoff: its reflected prime
terms and the terms at zero vanish, and its trivial-zero
series sums to the displayed integral. The last line
removes the Euler factors missing from the induced
character. Its omitted tail is $O(LP^{-1/2}\NN(F))$,
by point evaluation on the positive log axis and summing
each geometric series. The bounds
\eqref{eq:zeroshigh} and \eqref{eq:globalzeroshigh} together
apply to $\theta F$ after a fixed localization between
its initial bounded log interval and $r\ge2$. Thus,
writing $U^{\mathrm z}$ for the vector in
\cref{eq:fullzeroell}, we obtain
\begin{equation}
 \|U^{\mathrm z}(F)-U(F)\|_\chi
       \ll L^{-1/2}\tau^{-1}\NN(F).
 \label{eq:fullzeroU}
\end{equation}
For the binned version, set $\Delta=L^{-5/2}/\tau$. Replace
each retained ordinate $\gamma$ by its bin midpoint
$\widetilde\gamma$, repeating the midpoint according to
the count of retained labels. For any $B$ supported on an
interval of length $O(L)$, we have
\[
 \left\|\sum_{\gamma\in\mathcal Z_\chi}
             \sup_{|u-\gamma|\le\Delta}|\widehat B(u)|\right\|_\chi
 \ll L\|B\|_{H^1}.
\]
Indeed, since $\Delta\le L^{-1}$, a neighborhood of an
ordinate in $I_m$ lies in $I_m$ and its two neighbors.
Apply the triangle inequality in character mean square,
then Cauchy--Schwarz with weights $1+(m/L)^2$ as in
\cref{eq:shortzerosampling}. The other factor in Cauchy--Schwarz
is $O(L^{1/2})$ by \cref{eq:shortzerocount} and the estimate
\[
 \sum_m \frac{(1+L^{-1}\log(2+|m|/L))^2}{1+(m/L)^2}\ll L.
\]
The fundamental theorem of calculus bounds the replacement
error for the zero sum by $\Delta$ times the last
character-norm bound with $B=r\theta F$, after partitioning
it into smooth packets supported on intervals of length
$O(L)$. The packet at $r\asymp jL$, with $j\ge0$
including the initial interval, has $H^1$ norm
$O(L(1+j)e^{-j}\NN(F))$ by the exponential energy weight.
Summation gives an error $O(\Delta L^2\NN(F))$.
Thus \cref{eq:fullzeroU} holds for the binned vector, and
extends from smooth tests to the mesh tests by continuity.

For the diagonal let $\eta(r)$ be a smooth cutoff taking
values in $[0,1]$, zero on $r\le\log P$ and one on
$r\ge\log P+1$, with bounded derivatives. Replace
$\mathcal D$ by
\begin{align}
 \mathcal D^{\mathrm z}(F,G)
  ={}&\sum_{\substack{2\le n\le eP\\(n,N)=1}}
       (1-\eta(\log n))\frac{\Lambda(n)^2}{n}
                        F(\log n)\overline{G(\log n)}
             \notag\\
     &+\int_{\mathbb R}\eta(r)r F(r)\overline{G(r)}\,\ud r .
 \label{eq:fullzeroD}
\end{align}
Under the Riemann hypothesis for $\zeta$, the discrepancy
between $\sum_{n\le x,\ (n,N)=1}\Lambda(n)^2$ and
$\int_1^x\log t\,\ud t$ is
$O(\sqrt x\log^3(2x)+L\log(2x))$. The last term bounds
the contribution of the primes dividing $N$. For
overlapping unit log packets near $j$ and $k$,
the physical weight
$\eta(\log x)F_j(\log x)\overline{G_k(\log x)}/x$
has supremum and total variation
$O(e^{-j}e^{-2j/L}\NN(F_j)\NN(G_k))$ on the positive
log axis. Partial summation on packets meeting
$r\ge\log P$ therefore bounds the diagonal replacement
by a form of norm $O(L^{-4})$, using bounded overlap.
In detail the local coefficient is
$O((1+j)^3e^{-j/2}+L(1+j)e^{-j})$ for
$j\ge\log P-O(1)$, with the weight $e^{-2j/L}$ available.

For the matrix entries of the continuous model, set
\[
 J_{w,z}(u)=\int_0^\infty
       \bigl(w(x)\overline{z(x+u)}
              +w(x+u)\overline{z(x)}
              -w(x)\overline{z(u-x)}\bigr)\,\ud x
\]
with the tests extended by zero, and write
$\widetilde J_{w,z}(s)$ for its Mellin transform.
Mellin inversion and the singular-series expansion in
\cref{sec:zeros} give
\begin{equation}
 \mathcal M(F,G)
   =\frac{\SSS(N)}{2\pi i}\int_{\Re s=\alpha}
       N^{-s}\zeta(s)\zeta(1+s)H_N(s)
                        \widetilde J_{w,z}(s)\,\ud s .
 \label{eq:zeroMellin}
\end{equation}
Initially the line can be taken at $\Re s=2$. The pole
at one is removable since $\widetilde J_{w,z}(1)=0$.
For the mesh tests $w,z$ are compactly supported
$H^1$ functions on the physical axis. Two integrations
by parts, as in \cref{eq:modelresidue}, give
$\widetilde J_{w,z}(s)=O((1+|\Im s|)^{-2})$ in the
closed strip $\alpha\le\Re s\le 2$, with constants
depending on the tests and $N$. The zeta factor on this
strip is $O_\alpha((1+|\Im s|)^{1-\alpha})$ for
$|\Im s|\ge2$, by Euler summation. This justifies
the shift and absolute convergence of \cref{eq:zeroMellin}.
We can also truncate the Euler product in \cref{eq:zeroMellin}.
Let $H_N^{[P]}$ retain the factors at $p\le P$ and at all
$p\mid N$, and let $\mathcal M^{\mathrm z}$ use this
product in the integral. Write $h_N(d)$ for the Dirichlet
coefficients of $H_N$. At $p\nmid N$, the two nonconstant
coefficients are $a_N(p)-1/p$ and $-a_N(p)/p$, both
$O(p^{-2})$. At $p\mid N$ the only such coefficient is
$-1/p$. The absolute product bound
$\prod_{p\mid N}(1+1/p)\ll\log L$ and the sum
$\sum_{p>P}p^{-2}\ll P^{-1}$ show that the coefficients
removed by the truncation have absolute sum
$O(P^{-1}\log L)$.

Since $a_N=h_N*(n\mapsto n^{-1})$, the trace series
for $\mathcal M/\SSS(N)$ in \cref{sec:mesh} can be
written as
\[
 \sum_d h_N(d)\sum_{k\ge1}
          \frac{\mathcal C_{dkN}-\mathcal I}{k}.
\]
It is absolutely convergent in trace norm. Indeed the
bound $O(L D^{-2/L})$ for $\mathcal C_D-\mathcal I$
from the proof of \cref{thm:tracecompression} bounds
the inner series by $O(L^2d^{-2/L})$. Deleting the
indicated $h_N$ terms therefore changes the normalized
trace norm by $O(L^2P^{-1}\log L)$, proving the same
bound for $(\mathcal M-\mathcal M^{\mathrm z})/\SSS(N)$.
The finite product formula agrees with this series by
Mellin inversion on compact tests. The physical mesh
weights are themselves linear combinations of $x^{1/2}$,
$x^{-3/2-2\alpha}$ and $x^{-7/4}$ on each cell, by
\cref{eq:exactcoord}.

On the mesh take the difference of the matrices for
$\langle U^{\mathrm z}(F),U^{\mathrm z}(G)\rangle_\chi$
and $\mathcal D^{\mathrm z}+\mathcal M^{\mathrm z}$, both divided by
$\SSS(N)$. Call its norm $E_*$. For either the exact
or the new pair of forms on this space, positivity
bounds the character vector on unit tests by
$O(\sqrt{\SSS(N)(L+E)})$, where $E$ is its remainder norm.
Here \cref{eq:DMtotal} and the two replacement bounds
control the subtracted form. Starting with the smaller
of $E_*$ and $E_N^{\mathrm{mesh}}$, expansion with
\cref{eq:fullzeroU} then gives
\[
 \max(E_*,E_N^{\mathrm{mesh}})
   \le (1+O((L\tau)^{-1}))\min(E_*,E_N^{\mathrm{mesh}})
                  +O(\tau^{-1}).
\]
Indeed the two mixed terms cost
$O(\tau^{-1}\sqrt{1+E/L})$, and the squared character
error and the two replacements are smaller.

The new positive character matrix has trace norm
$O(L^2)$ on the normalized scale, by the damped
Fourier bound in the proof of \cref{eq:traceterms},
applied to the truncated zero sum in \cref{eq:fullzeroell}.
For the rounded ordinates the same bound follows by using
counts in the neighboring original bins.
The trivial-zero functional is bounded and the finite
Euler-factor correction has dual norm $O(L)$ in each
character. The diagonal of \cref{eq:fullzeroD} has
trace $O(L^2)$ by integrating the squared dual point
evaluation norm in its continuous part and using the
old atom bound in its discrete part. The model inherits
the trace bound of \cref{eq:traceterms}. Compress the two
matrices separately as in \cref{thm:tracecompression},
using the span of their spectral directions above
$\tau^{-1}$ in absolute value. This leaves $O(L^2\tau)$
dimensions and changes the norm of their difference by
at most $4/\tau$. The last comparison and
\cref{thm:mesh} prove \cref{eq:fullzerodata}.

\subsection{A necessary pair-correlation law}

The following consequence concerns pairs from the same odd
character, at a single even modulus. The sine kernel is the
usual pair-correlation prediction for families of Dirichlet
$L$-functions. In the prime-conductor setting, related
form-factor statements are connected to Hooley's variance
by~\cite{TZLiu}.

\begin{theorem}[Pair correlation forced by hypothesis S]
\label{thm:necessarypair}
Assume $\grh(N)$ and the norm bound of \cref{thm:localderivative}
along the sequence of even moduli $N$. For every nonnegative
$f\in C_c^\infty(\mathbb R)$, define the even gap measure
\[
 \nu_{N,f}=\frac1L\mathbf E_N
     \sum_{\gamma,\gamma'\in\mathcal Z_\chi}
          \frac{f(\gamma)+f(\gamma')}{2}\,
          \delta_{L(\gamma'-\gamma)/(2\pi)} .
\]
The two sums are over zero labels and include equal labels.
Then, against every continuous compactly supported test,
\begin{equation}
 \nu_{N,f}\ \longrightarrow\
 \frac{\int f}{2\pi}
 \left[\delta_0+
       \left(1-\left(\frac{\sin(\pi x)}{\pi x}\right)^2\right)
                      \ud x\right].
 \label{eq:necessarypair}
\end{equation}
In every fixed bounded interval of positive length, the
ratio of simple labels to all zero labels tends to one,
counting labels across all the odd characters of a modulus.
\end{theorem}

Fix $b>0$, take $a=b/L$, and put
\[
 K_a(t)=\frac{a}{a^2+t^2},\qquad
 Z_{\chi,a}(t)=\sum_{\gamma\in\mathcal Z_\chi}K_a(t-\gamma).
\]
The real part of the completed logarithmic derivative of
the inducing character at $1/2+a+it$ is $Z_{\chi,a}(t)$.
Indeed the real part of the Hadamard constant cancels at
$1/2+it$ by the functional equation, and the remaining
real parts are the displayed Poisson terms. On a compact
interval in $t$, this gives
\[
 Z_{\chi,a}(t)=\Re D_\chi(t)+\tfrac12\log q_\chi+O(\sqrt L).
\]
The gamma factor is bounded on the interval.
The imprimitive Euler factors contribute $O(\sqrt L)$,
since $\sum_{p\mid N}p^{-1/2}\log p\ll\sqrt L$ by
splitting at $p=L$.
The conductor computation in \cref{sec:zeros} gives
$\mathbf E_N\log q_\chi=L+O(\log L)$. As
$0\le\log q_\chi\le L$, the normalized conductor term
tends to $1/2$ in character mean square. From the
first and second moment limits in \cref{thm:localderivative},
using Cauchy--Schwarz for the mixed error, we obtain
\begin{equation}
 \frac1{L^2}\mathbf E_N\int f(t)Z_{\chi,a}(t)^2\,\ud t
       \longrightarrow \left(\frac{c(b)}2+\frac14\right)\int f .
 \label{eq:poissonzsquare}
\end{equation}

Write $P_y(x)=y/(\pi(y^2+x^2))$ for the Poisson kernel.
We next use
\begin{equation}
 \frac1{L^2}\mathbf E_N\int f(t)Z_{\chi,a}(t)^2\,\ud t
       =\frac{\pi}{2}\int P_{b/\pi}(x)\,\ud\nu_{N,f}(x)+o(1).
 \label{eq:gapstieltjes}
\end{equation}
For an unweighted product of two kernels, the integral is
$\pi K_{2a}(\gamma'-\gamma)$ by the Poisson semigroup.
It is enough to replace $f(t)$ in the product integral by
$f(\gamma)$, since the resulting integral is even in the gap
and hence has the same value with the symmetrized weight.

Here are bounds for this replacement and the tails. On each
gap bin $[k,k+1]$, $k\in\mathbb Z$, the mass of $\nu_{N,f}$
is $O_f(1+L^{-1}\log(2+|k|))$. Indeed the first ordinate
in the unsymmetrized sum lies in $O_f(L)$ bins of width
$1/L$, each of bounded character mean square by
\cref{eq:shortzerocount}. The second ordinate for each such
bin lies in a bounded number of bins at distance $O((1+|k|)/L)$,
and Cauchy--Schwarz gives the bound. Symmetrization preserves it.
For $t$ in a compact set, binning the Poisson kernel using
the same estimate gives $\|Z_{\chi,a}(t)\|_\chi\ll_b L$ there.
The sum over $|\gamma-t|>R/L$ instead has character norm
$O_b(L(1+\log(2+R))/R)$ for $R\ge2$.

On $|t-\gamma|\le R/L$, the replacement costs $O_f(R/L)$
times an integral of the squared kernel sum over a bounded
interval. On the complement, the part with weight $|f(t)|$
costs $O_{b,f}(L^2(1+\log(2+R))/R)$ before normalization,
by the character norm bounds. For the part with weight
$|f(\gamma)|$, gaps of size at least $R/(2L)$ have the same
bound by the gap-bin estimate and the full product integral.
With smaller gaps, $|t-\gamma'|\ge |t-\gamma|/2$ on the
discarded $t$ range, so the product integral there is
$O_b(L/R^3)$. There are $O_f(L(R+1))$ of these weighted
pairs in character mean by the gap-bin estimate, for fixed
$R$ as $N\to\infty$. After dividing by $L^2$, take $N\to\infty$
and then $R\to\infty$. This proves the replacement, and the
normalization of the full product integral gives
\cref{eq:gapstieltjes}.

The two equations identify all the Poisson integrals of any
vague subsequential gap limit, since the gap-bin bound gives
uniformly negligible weighted tails. Their values are
$(c(b)+1/2)(\int f)/\pi$ at height $b/\pi$ for each $b>0$.
The measure on the right of \cref{eq:necessarypair} has
exactly these integrals. In fact the atom contributes $1/b$
before the common factor $(\int f)/(2\pi)$, and convolution
of the squared sinc with $P_{b/\pi}$ at zero gives
\[
 2\int_0^1(1-u)e^{-2bu}\,\ud u
                =\frac1b-\frac{1-e^{-2b}}{2b^2}
\]
by the triangular Fourier transform of that squared sinc.
These integrals determine an even measure with
$\int(1+x^2)^{-1}\,\ud\nu<\infty$. To see uniqueness,
weight the difference of two such measures by $(1+x^2)^{-1}$
and push it forward by $s=(1+x^2)^{-1}$ to $[0,1]$.
The Stieltjes identities become zero integrals of
$(1+(y^2-1)s)^{-1}$ for $y>0$. Expansion near $y=1$
gives zero moments of every order. Polynomial density
and evenness then identify the original measures.
This proves \cref{eq:necessarypair}.
Conversely, under $\grh(N)$, that gap limit for all $f$
implies \cref{eq:localderivativelaw} for every fixed $b>0$.
Indeed the uniform gap-bin bound and \cref{eq:gapstieltjes}
give \cref{eq:poissonzsquare}. The conductor relation and
the vanishing limits for the first and complex second
moments in \cref{thm:localderivative} recover the mean
square of $D_\chi$ from it.

For the last assertion, the first moment from
\cref{thm:localderivative} and the conductor estimate give
$L^{-1}\mathbf E_N\int f(t)Z_{\chi,a}(t)\,\ud t\to \tfrac12\int f$.
The unit-interval zero count, averaged and divided by $L$,
is $O(1+L^{-1}\log(2+|t|))$ by \cref{eq:shortzerocount}.
It allows replacing $P_a*f$ by $f$ in this limit:
the approximation is uniform on compacts and the tails
are $O_f(a/(1+t^2))$ away from the support. Thus the
weighted zero-label count divided by $L$ tends to
$(\int f)/(2\pi)$. The mass $\nu_{N,f}(\{0\})$ counts
the squared multiplicities with the same weights.
By the upper bound for mass on a closed singleton
in \cref{eq:necessarypair}, its upper limit is also
at most $(\int f)/(2\pi)$. Their difference bounds
the labels at multiple zeros. Taking $f$ positive on
the interval in question proves the assertion there,
since the limiting one-level measure has positive constant
density.

\subsection{Zero covariance through the critical product}

For microscopic variables use the Fourier transform
$(\mathcal F h)(u)=\int h(x)e^{-2\pi iux}\,\ud x$. Set
\[
 X_{\chi,h}(t)=\sum_{\gamma\in\mathcal Z_\chi}
                  h\!\left(\frac{L(\gamma-t)}{2\pi}\right).
\]
The expectations below are over all odd characters at the
single modulus $N$. The prime-conductor calculations
of Hughes and Rudnick~\cite{HughesRudnick} give the
corresponding first and restricted second moments for a
common filter centered at zero.

\begin{theorem}[Mixed covariance and pair correlation in restricted Fourier support]
\label{thm:restrictedzeropairs}
Assume $\grh(N)$. Let $h,g$ be real even Schwartz functions
with Fourier supports in $[-a,a]$ and $[-b,b]$, where
$a,b>0$ and $a+b\le2$. For fixed $A,D>0$, uniformly for
$|t|,|s|\le L^A$ and $|t-s|\le D$,
\begin{equation}
 \operatorname{Cov}_{\chi}\bigl(X_{\chi,h}(t),X_{\chi,g}(s)\bigr)
  =\int_{\mathbb R}|u|\,\mathcal Fh(u)\mathcal Fg(u)
                         e^{-i(t-s)Lu}\,\ud u
       +O_{h,g,A,D}\!\left(\frac{\log L}{L}\right).
 \label{eq:restrictedcovariance}
\end{equation}
In particular, the gap measure of \cref{thm:necessarypair}
has the limit \eqref{eq:necessarypair} against real even
Schwartz tests $\Phi$ with $\mathcal F(\Phi)$ compactly
supported in $(-1,1)$. This remains true after replacing
the height weight by $f((t-t_0)/H)$ and dividing both sides
by $H$, where $f\in C_c^\infty(\mathbb R)$ is nonnegative,
$|t_0|\le L^A$, $0<H\le1$, and $LH\to\infty$.
The limiting height factor after division is $\int f$.
For every interval within $[-L^A,L^A]$ of length $H$ in
this range, the ratio of simple labels to all zero labels,
summed over the odd characters, has lower limit at least
\[
 \frac32-\frac1{\sqrt2}\cot\!\left(\frac1{\sqrt2}\right)
     =0.672500703679\ldots.
\]
\end{theorem}

The Weil explicit formula for a translated test
\cite[Lemma~1]{TZLiu} gives
\begin{equation}
 \begin{aligned}
 X_{\chi,h}(t)&=\mathcal Fh(0)-P_{\chi,h}(t)
    -\overline{P_{\chi,h}(t)}+\varepsilon_{\chi,h}(t),\\
 P_{\chi,h}(t)&=\frac1L\sum_{n\ge2}
       \frac{\Lambda(n)\chi(n)}{\sqrt n}\,
          \mathcal Fh(\log n/L)e^{-it\log n}.
 \end{aligned}
 \label{eq:restrictedlinstat}
\end{equation}
We show $\|\varepsilon_{\chi,h}(t)\|_\chi
\ll_{h,A}(\log L)/L$. The gamma term, including $-\log\pi$,
is of this size uniformly for $|t|\le L^A$, since its symbol
is $O(\log(2+|t|))$ and the translated kernel is Schwartz at
scale $1/L$.
For the conductor term, consider first the average over all
characters. If $p^\nu\Vert N$, the lost conductor exponent
$D_p$ satisfies $\Pr(D_p\ge j)\le2p^{-j}$ for
$1\le j\le\nu$, by counting the characters modulo $p^{\nu-j}$.
The $D_p$ are independent in this average by the Chinese
remainder theorem. Thus the mean of $\log(N/q_\chi)$ is
$O(\sum_{p\mid N}(\log p)/p)$, and its variance is
$O(\sum_{p\mid N}(\log p)^2/p)$. The two sums are
$O(\log L)$ and $O((\log L)^2)$, respectively, by
splitting at $p=L$ and using Chebyshev's bound below it.
The odd-character mean square is at most twice the
all-character mean square, so $\|\log(N/q_\chi)\|_\chi\ll\log L$.
In replacing $\chi^*$ by $\chi$, only powers of $p\mid N$
can contribute, and only for characters whose local conductor
at $p$ is trivial. Their odd-character proportion is at most
$2/\varphi(p^\nu)$. The triangle inequality therefore bounds
the unscaled Euler error by
\[
 O_h\!\left(\sum_{p^\nu\Vert N}
      \frac{\log p}{\sqrt{p\varphi(p^\nu)}}\right)=O_h(\log L).
\]
These give the claimed error in \cref{eq:restrictedlinstat}.
Also $\|X_{\chi,h}(t)\|_\chi\ll_{h,A}1$ by
\cref{eq:shortzerocount}, summing the Schwartz majorant in bins.

The character mean of $P_{\chi,h}$ is $O_h(N^{a/2-1})$
for $a<2$. Indeed, only $n\equiv\pm1\pmod N$ with
$2\le n\le N^a$ contribute. Their inverse square roots sum
to $O(N^{a/2-1})$ if the range is nonempty, and
$\Lambda(n)/L\ll_a1$. This proves
$\mathbf E_N X_{\chi,h}(t)=\mathcal Fh(0)+O_{h,A,a}((\log L)/L)$.
The same holds for $g$. In the covariance we may therefore
use the products of $P_{\chi,h}+\overline{P_{\chi,h}}$ and
$P_{\chi,g}+\overline{P_{\chi,g}}$ at cost $O((\log L)/L)$,
by Cauchy--Schwarz and the mean-square bounds above.

In $\mathbf E_N P_{\chi,h}(t)\overline{P_{\chi,g}(s)}$,
odd-character orthogonality produces the equal and opposite
unit residues of \cref{eq:form}. The Fourier restrictions
give $mn\le N^{a+b}\le N^2$. The absolute off-diagonal
estimate in the proof of \cref{thm:subcritical} gives a
normalized bound $O(Q^{-1/2}+N^{-2/5})$ on $mn\le N^2/Q$.
Here it is applied with smooth envelopes for
$\mathcal Fh(r/L)$ and $\mathcal Fg(r/L)$ on the positive log
axis. The envelopes equal one on the respective Fourier
support intervals, have log supports of length $O(L)$,
and each has energy norm $O(\sqrt L)$.
If $N^2/Q<mn\le N^2$ is within the two supports, then
$0\le a-\log n/L<\log Q/L$ and similarly for $b-\log m/L$,
where $n$ is the input to $h$.
The Fourier transforms vanish to every order at the
endpoints. The product of the weights is thus
$O_M((\log Q/L)^M)$ for each $M>0$.
On this part use the same envelope bound with $Q=1$.
Taking $Q$ to be a sufficiently large fixed power of $L$
now gives $O_J(L^{-J})$ for the off-diagonal part of the
averaged product, for every $J>0$, after the factor $L^{-2}$.

For $\mathbf E_N P_{\chi,h}(t)P_{\chi,g}(s)$ the congruence
condition is $nm\equiv\pm1\pmod N$. Write $d=nm$. Its total
von Mangoldt coefficient is at most $(\Lambda*\Lambda)(d)
\le(\log d)^2$, since $\sum_{n\mid d}\Lambda(n)=\log d$.
On $d\le N^2/Q$ the sum of $d^{-1/2}$ over these two
progressions is $O(Q^{-1/2})$ when nonempty. On the upper
part it is $O(1)$, and the endpoint bound for the Fourier
weights applies again. This product consequently contributes
$O_J(L^{-J})$ for every $J>0$.

The remaining two terms are conjugates of the diagonal
in the first product. By partial summation using the
prime-square weight estimate in the proof of
\cref{eq:modelresidue}, that diagonal is
\[
 \frac1{L^2}\int_0^\infty r\,
    \mathcal Fh(r/L)\mathcal Fg(r/L)e^{-i(t-s)r}\,\ud r
       +O_{h,g,D}\!\left(\frac{1+(\log L)^2}{L^2}\right).
\]
Indeed the $O(\sqrt x\log^3(2x))$ discrepancy for the
unrestricted sum has an integrable error against a weight
of derivative $O_{h,g,D}(x^{-2})$. Excluding $p\mid N$
costs $O(\sum_{p\mid N}(\log p)^2/(p-1))=O((\log L)^2)$
before division by $L^2$. The integral and its conjugate
give \cref{eq:restrictedcovariance}.

For the gap measure let $f_H(t)=f((t-t_0)/H)$, with $LH\ge1$.
Fixed compactly supported smooth weights are included by
taking $H=1$. For real even $h,g$ as above,
\begin{equation}
 \int f_H(t)\,\mathbf E_N X_{\chi,h}(t)X_{\chi,g}(t)\,\ud t
      =2\pi\int (h*g)\,\ud\nu_{N,f_H}+O_{f,h,g,A}(L^{-1}).
 \label{eq:restrictedgaptransfer}
\end{equation}
To verify this, replace $f_H(t)$ by $f_H(\gamma)$ inside
the integrated kernel product. Their difference is bounded
by $O_f(|t-\gamma|/H)$ and vanishes unless $t$ or $\gamma$
lies in the support of $f_H$. For the part with $t$ in this
support, \cref{eq:shortzerocount} bounds the character norm
of each absolute kernel sum by $O(1)$, with a gain $L^{-1}$
for the sum containing $|t-\gamma|$. Integrating over a
length $O(H)$ gives $O(L^{-1})$. With $\gamma$ in the
support, the absolute integral for a pair is
$O_M((L^2 H)^{-1}(1+|x|)^{-M})$, where
$x=L(\gamma'-\gamma)/(2\pi)$. This follows by scaling the
product of the two Schwartz kernels. The mean number of
these pairs with $x\in[k,k+1]$ is bounded by
\[
 O_{A,f}\bigl((LH+1)(1+L^{-1}\log(2+|k|))\bigr).
\]
Indeed the first ordinate has $O_f(LH+1)$ bins of width
$1/L$, the gap constrains the second to a bounded number
per first bin, and Cauchy--Schwarz uses
\cref{eq:shortzerocount}. Summation gives the same
$O(L^{-1})$ error. The integrated unweighted kernel
product is $(2\pi/L)(h*g)(x)$. It is even, which allows
symmetrizing the weight in $\nu_{N,f_H}$.

For a real even Schwartz $\Phi$ with $\mathcal F(\Phi)$
compactly supported in $(-1,1)$, take $\mathcal Fh
=\mathcal F(\Phi)$ and a real even $\mathcal Fg\in
C_c^\infty(-1,1)$ equal to one on its support.
Then $h*g=\Phi$, and using the means and covariance at
equal heights gives
\begin{equation}
 \frac1H\int\Phi\,\ud\nu_{N,f_H}
   =\frac{\int f}{2\pi}\left(\mathcal F(\Phi)(0)
                   +\int_{-1}^1 |u|\mathcal F(\Phi)(u)\,\ud u\right)
       +O\!\left(\frac{\log L}{L}+\frac1{LH}\right).
 \label{eq:restrictedpairlaw}
\end{equation}
The expression in parentheses is the sine-kernel functional,
by the triangular transform of the squared sinc. The
conductor and Euler errors and the smoothing transfer are
uniform for $|t_0|\le L^A$.

For the simplicity bound, first use \cref{eq:restrictedlinstat}
in the same transfer with a single $h$ and $\mathcal Fh(0)=1$.
The replacement error is $O(L^{-1})$, by the same argument
using $O_f(LH+1)$ expected labels in the support for the
$\gamma$ term. It follows that
$(LH)^{-1}\mathbf E_N\sum_\gamma f_H(\gamma)\to(\int f)/(2\pi)$.
Now take $\Phi=|\check w|^2$, where $\check w$ is the inverse
Fourier transform of an even real $w\in C_c^\infty(-1/2,1/2)$
with $\int w=1$. Since $\Phi\ge0$ and $\Phi(0)=1$, the right
side of \cref{eq:restrictedpairlaw} bounds the weighted mass
of coincident pairs after division by $H$. Use smooth
approximations to the Montgomery--Taylor test
\cite[Section~3.5]{CarneiroEtAl},
\[
 w(u)=\frac{\cos(\sqrt2 u)}{\sqrt2\sin(1/\sqrt2)}
           \mathbf1_{[-1/2,1/2]}(u).
\]
The expression in parentheses in \cref{eq:restrictedpairlaw}
converges for these approximations to
$\int w^2+\iint |u-v|w(u)w(v)\,\ud u\,\ud v$.
On $[-1/2,1/2]$, the function
$w(u)+\int |u-v|w(v)\,\ud v$ is even and has zero second
derivative. Its value at $1/2$ is
$1/2+(1/\sqrt2)\cot(1/\sqrt2)$. This gives the displayed
quadratic expression since $\int w=1$.
For a sharp height interval, sandwich its indicator between
smooth nonnegative weights whose integrals tend to its
length after scaling. The first moment estimate and the
upper bound on coincident pairs then give the same constant
as an upper limit for the ratio of squared multiplicities to
total labels. A zero of multiplicity $m\ge2$ contributes
$m^2-m\ge m$ to their difference, giving the claimed
lower limit for the simple labels. Related results for
the prime-modulus family without GRH are in~\cite{HuaYang}.

\subsection{Averaging a common logarithmic translation}

Let $E_N(u,v)$ denote the norm of the original form divided
by $\SSS(N)$ when the log supports lie in $[u,u+1]$ and
$[v,v+1]$, respectively. Both tests use the norm $\NN$.

\begin{theorem}[Mean high-support norm at a fixed modulus]
\label{thm:shiftmean}
Fix $\delta>0$ and assume $\grh(N)$. Uniformly for
$H_1,H_2\ge(1+\delta)L$, with $L=\log N$, there is $c_\delta>0$
such that
\begin{equation}
 \frac1L\int_0^L E_N(H_1+s,H_2+s)\,\ud s
     \ll_\delta L e^{-(H_1+H_2)/L}+N^{-c_\delta}.
 \label{eq:shiftmean}
\end{equation}
In particular, if $H_1+H_2\ge(1+\varepsilon)L\log L$
for fixed $\varepsilon>0$, the set of $s\in[0,L]$ on which
$E_N(H_1+s,H_2+s)>L^{-\varepsilon/2}$ has measure
$O_{\delta,\varepsilon}(L^{1-\varepsilon/2})$.
\end{theorem}

Let $\mathcal H=H^1_0(0,1)$ with the usual $H^1$ norm.
For a real ordinate $\gamma$, write $v_\gamma$ for the
functional $b\mapsto\widehat b(\gamma)$ in the Hilbert dual
$\mathcal H^*$. Integration by parts gives
$\|v_\gamma\|_{\mathcal H^*}\ll(1+|\gamma|)^{-1}$.
For a translate with left endpoint $u$, set
\[
 A_N(u)=\frac{2}{\varphi(N)}
   \sum_{\chi\ {\rm odd}}
     \left\|\sum_{\gamma\in\mathcal Z_\chi}
                    e^{iu\gamma}v_\gamma\right\|_{\mathcal H^*}^2 .
\]
We have $L^{-1}\int_0^L A_N(H+s)\,\ud s\ll L$ uniformly
in real $H$. To prove this first truncate the ordinates.
Majorize the indicator of $0\le s/L\le1$ by a fixed
nonnegative Schwartz function whose Fourier transform is
supported in $[-1,1]$, as in the construction for
\cref{eq:shortzerocount}. On expanding the dual norm squared
and integrating over $s$, only pairs with
$|\gamma-\gamma'|\le1/L$ remain. The normalized integral
of each exponential has bounded absolute value.
In the bins $I_m=[m/L,(m+1)/L)$, write $M_{\chi,m}$ for
the number of zero labels. Pairs occur only in the same
or neighboring bins. The dual-norm bounds and
$2ab\le a^2+b^2$ bound the integral by a constant times
the character mean of
\[
 \sum_{m\in\mathbb Z}\frac{M_{\chi,m}^2}{1+(m/L)^2}.
\]
By \cref{eq:shortzerocount}, this mean is
$O(\sum_m(1+\log(2+|m|/L)/L)^2/(1+(m/L)^2))=O(L)$.
The truncated functional converges in $\mathcal H^*$ for
each character and fixed $N$, by \cref{eq:zeroshigh}
applied to the translated test. The bound passes to the
limit by Fatou's lemma.

On support $[u,u+1]$, the unweighted $H^1$ norm is
$O(e^{-u/L}\NN(F))$ when $u\ge0$. The covariance term of
\cref{eq:zeroform}, by Cauchy--Schwarz over the characters,
has normalized norm on the two windows bounded by
$O(e^{-(u+v)/L}\sqrt{A_N(u)A_N(v)})$.
The archimedean form has norm $O(L e^{-(u+v)/L})$ by
\cref{eq:archdensity}, Plancherel and the log growth of
the digamma function. Apply Cauchy--Schwarz over
$s\in[0,L]$ to the first bound and use the two estimates
for $A_N$. The error term of \cref{eq:zeroform} is uniform
on these windows, giving \cref{eq:shiftmean}. The final
assertion follows from the mean bound by Markov's inequality.

\subsection{A uniform upper product cutoff}

Fix a smooth function $\Theta\colon\mathbb R\to[0,1]$ with
$\Theta(t)=0$ for $t\le0$ and $\Theta(t)=1$ for $t\ge1$.
For $H\ge5L$, let $T_N^{\rm tail}(H)$ be the form of $T_N$
with both the arithmetic and continuous pair weights multiplied by
\begin{equation}
 \Theta\!\left(\frac{\log(xy)-H}{L}\right).
 \label{eq:uppertailmask}
\end{equation}
Denote the three terms in \cref{eq:split} with this weight by
$\mathcal A_H,\mathcal D_H,\mathcal M_H$. In the diagonal term
the weight is evaluated at $x=y$. Their norms below use $\NN$
in the two arguments.

Under $\grh(N)$ the combined deterministic form on unit log
packets with integer centers $j,k$ satisfies
\begin{equation}
 |(\mathcal D+\mathcal M)(F_j,G_k)|
       \ll \SSS(N)L e^{-|j-k|/4}\|F_j\|_{H^1}\|G_k\|_{H^1}.
 \label{eq:deterministicblock}
\end{equation}
To see this, take $j\le k$. The proof of \cref{eq:Mpacket}
with unweighted $H^1$ gives a normalized model block bound
$O((1+(k-L)_+)e^{-(k-j)/2})$. The diagonal has bound $O(1+k_+)$
on overlapping packets. These prove \cref{eq:deterministicblock}
for $k\le3L$. If $k>3L$ and $k-j>(k-L)/64$, the diagonal
vanishes for large $N$ and the model bound is $O(e^{-(k-j)/4})$.
In the remaining range $k>3L$, $k-j\le(k-L)/64$, we have $j\ge2L$.
Apply \cref{eq:modelresidue} with $\nu=1/100$.
Its error in the unweighted norms is bounded by a constant times
\[
 \exp\bigl(L/100+2(k-j)-(j-L)/8\bigr)\le e^{-(k-L)/16}.
\]
This is $O(e^{-(k-j)/4})$, as is the diagonal error
$O(e^{-j/2}(1+j+L)^3)$ on overlap.
The main terms together give $\mathcal W_N$ up to the
$O(L/\varphi(N))$ conductor correction on overlap.
For neighboring packets the symbol bounds this form by
$O(L)$ in $H^1$. On packets separated by at least one
its kernel is $O(e^{-3|r-s|/2})$.
This proves \cref{eq:deterministicblock}.

\begin{theorem}[An upper product cutoff at every modulus]
\label{thm:upperproduct}
Assume $\grh(N)$. Uniformly for $H\ge5L$,
\begin{align}
 \|\mathcal A_H\| &\ll_\Theta L^2 e^{-H/L},
       \label{eq:uppercharacter}\\
 \frac{\|\mathcal D_H+\mathcal M_H\|}{\SSS(N)}
       &\ll_\Theta L e^{-H/L}.
       \label{eq:uppermodel}
\end{align}
Consequently, for every $\kappa\ge1$ the tapered deletion of
products above $N^{2\log\log N+\kappa}$ at
$H=L(2\log L+\kappa)$ costs
$O_\Theta(e^{-\kappa})$ in normalized norm.
Already at $H=L(\log L+\kappa)$, the tail form differs from
its character term by $O_\Theta(e^{-\kappa})$
on the normalized scale, for sufficiently large $N$.
\end{theorem}

For the character term use a bounded-overlap smooth partition
into log intervals $[(j-1)L,(j+1)L]$, with derivatives
$O(L^{-1})$. The pieces with $j<0$ contain no arithmetic
terms. For $j\ge2$, \cref{eq:Utail} bounds the character
operator on the $j$th piece by $O(Le^{-j})$. For $j=0,1$
the same bound, with a different absolute constant, follows
by splitting off $r\le3$ and estimating its finitely many
prime powers by point evaluation.

On a packet pair write $t=(r+s)/L-j-k$, where $r=\log x$
and $s=\log y$. Multiply $\Theta(t+j+k-H/L)$ by a
fixed smooth bump equal to one on $[-2,2]$.
Its Fourier transform has uniformly bounded
$L^1$ norm with weight $(1+|\xi|)^2$, by four
integrations by parts. Inverting the transform separates
the weight into products of phase multipliers
$e^{i\xi r/L}$ and $e^{i\xi s/L}$. The product rule costs
$O((1+|\xi|)^2)$ in the two test norms. Integrating the
packet bounds gives $O_\Theta(L^2 e^{-j-k})$, and
the weighted form is zero when $j+k+2\le H/L$.
The row and column sums of these coefficients are
$O_\Theta(L^2 e^{-H/L})$ by a geometric sum.
The squared norms of the localized tests sum to a bounded
multiple of their starting squared norm, so the Schur
bound gives \cref{eq:uppercharacter}.

For the diagonal and model, use unit log packets with integer
centers $j,k$. Fourier inversion on $t=r+s-j-k$ bounds
the masked block by a constant times the bound in
\cref{eq:deterministicblock}. Only $j+k>H-2$ can occur.
For nonnegative $j,k$, conversion to the energy norm
costs $O(e^{-(j+k)/L})$. If $j<0<k$, it costs
$O(e^{-j/8}e^{-k/L})$, since $W_N(r)\gg e^{r/4}$
for $r\le0$. The additional factor is at most $e^{|j-k|/8}$.
The same reasoning applies for $k<0<j$.
The resulting coefficients on the normalized scale are
$O(L e^{-(j+k)/L}e^{-|j-k|/8})$, with row and column sums
$O(L e^{-H/L})$ on the retained pairs.
This proves \cref{eq:uppermodel}. The estimates extend from compact
smooth tests to the operator norms by density, and
$\SSS(N)\gg1$ gives the stated deletion error.

\subsection{A tighter cutoff when one input is large}

Let $\Theta$ be the same cutoff. For $\kappa\ge1$ put
\[
 p_0=L\bigl(\tfrac32\log L+\kappa\bigr),\qquad
 q_0=L(\log L+\kappa),\qquad
 \vartheta(r)=\Theta((r-q_0)/L).
\]
In the physical form multiply every pair weight by $1-m(r,s)$,
where $r=\log x$, $s=\log y$ and
\begin{equation}
 m(r,s)=\Theta((r+s-p_0)/L)
       \bigl(\vartheta(r)+\vartheta(s)-\vartheta(r)\vartheta(s)\bigr).
 \label{eq:asymuppertailmask}
\end{equation}
Write $E_{\rm ret}$ for the norm of this retained form divided
by $\SSS(N)$, and $E=\|T_N\|/\SSS(N)$.

\begin{theorem}[An upper cutoff for pairs with one large input]
\label{thm:asymupperproduct}
Assume $\grh(N)$. For sufficiently large even $N$, uniformly
in $\kappa\ge1$,
\begin{equation}
 |E-E_{\rm ret}|\ll_\Theta
      e^{-\kappa}\bigl(1+\sqrt{\min(E,E_{\rm ret})}\bigr).
 \label{eq:asymupperproduct}
\end{equation}
In particular, for any $\kappa=\kappa(N)\to\infty$, hypothesis S
is equivalent to $E_{\rm ret}\to0$. The retained form is supported
in the union of
\[
 \max(x,y)\le N^{\log L+\kappa+1},
 \qquad
 xy\le N^{(3/2)\log L+\kappa+1}.
\]
\end{theorem}

By the Schur bound and bounded overlap, \cref{eq:deterministicblock}
gives
\begin{equation}
 |(\mathcal D+\mathcal M)(F,G)|
       \ll \SSS(N)L\|F\|_{H^1}\|G\|_{H^1}.
 \label{eq:deterministicunweighted}
\end{equation}

For the deleted deterministic form the proof of
\cref{eq:uppermodel} applies at $H=p_0$. In the Fourier inversion
on a unit packet pair the additional factor
$\vartheta(r)+\vartheta(s)-\vartheta(r)\vartheta(s)$ is a sum
of three separated multipliers with uniformly bounded norms.
It preserves the bounds for each packet pair and the
restriction $j+k>p_0-2$ used there. The normalized deterministic
deletion is thus $O_\Theta(Le^{-p_0/L})$.
The finite norm of the deleted character form follows at once
from \cref{eq:Utail} and the packet proof of
\cref{eq:uppercharacter}, with the same three multipliers.

Let $e_*=\min(E,E_{\rm ret})$. On tests $B$ supported in
$0\le r\le q_0$ the two signed forms agree. By
\cref{eq:split,eq:deterministicunweighted}, positivity gives
\begin{equation}
 \frac{\|U(B)\|_\chi}{\sqrt{\SSS(N)}}
 \ll \sqrt{e_*}\,\NN(B)+\sqrt L\,\|B\|_{H^1}.
 \label{eq:asymupperpositivity}
\end{equation}
For the character deletion, we may multiply each argument by
the cutoff $\theta$ of \cref{eq:fullzeroell}, since it equals one
at every arithmetic atom. This costs a bounded factor in the
energy and puts both supports in $r\ge0$. Split the resulting
tests into lower pieces supported in $r\le q_0$ and upper pieces
supported in $r\ge q_0-L$, using the multiplier
$\Theta((r-q_0+L)/L)$ for the upper pieces. The lower--lower
deletion vanishes. For upper--upper, use the length-$L$ packet
bounds $O(Le^{-j})$ for $U$ and the Fourier inversion from
\cref{eq:uppercharacter}. The first nonzero packets have
$j,k\ge q_0/L-2$. The coefficients $O_\Theta(L^2e^{-j-k})$
have row and column sums $O_\Theta(L^2e^{-2q_0/L})$, so this
part is $O_\Theta(e^{-2\kappa})$.

For the mixed part take the lower test $F$ and the upper
test $G$, each of bounded energy norm. The mask on these
supports is $\Theta((r+s-p_0)/L)\vartheta(s)$. Localize
$\vartheta G$ into length-$L$ packets $G_k$ supported in
$[(k-1)L,(k+1)L]$, with $k\ge q_0/L-1$ and
$\sum_k\NN(G_k)^2\ll_\Theta1$.
In the variable $v=s/L-k$, insert a fixed smooth bump
equal to one for $|v|\le1$ and zero for $|v|\ge2$. Fourier
inversion now expresses $\Theta((r+s-p_0)/L)$ on the packet
as an integral of $a_\xi(r)e^{i\xi v}$. The multipliers satisfy
\[
 |a_\xi(r)|+L|a_\xi'(r)|\ll_\Theta (1+|\xi|)^{-4},
 \qquad a_\xi(r)=0\quad (r\le p_0-(k+2)L).
\]
These estimates follow by integration by parts in $v$,
also after differentiating in $r$.
Use \cref{eq:asymupperpositivity} for $a_\xi F$ and
\cref{eq:Utail} for $e^{-i\xi v}G_k$. The latter multiplier
costs $O(1+|\xi|)$ in energy. After integration over $\xi$,
the normalized mixed block is bounded by
\[
 O_\Theta\!\left(L e^{-k}\NN(G_k)
       \bigl(\sqrt{e_*}+\sqrt L\,b_k\bigr)\right),
 \qquad
 b_k^2=\int_{r>p_0-(k+2)L} (|F|^2+|F'|^2)\,\ud r.
\]
The part with $\sqrt{e_*}$ sums to
$O_\Theta(L e^{-q_0/L}\sqrt{e_*})$ by Cauchy--Schwarz.
For the other part, the sum of $e^{-2k}$ over the
indices with $r>p_0-(k+2)L$ is
$O(e^{-2p_0/L+2r/L})$. Hence
\[
 \sum_k e^{-2k}b_k^2
 \ll e^{-2p_0/L}\int_0^\infty e^{2r/L}
                 (|F|^2+|F'|^2)\,\ud r
 \ll e^{-2p_0/L}\NN(F)^2.
\]
Cauchy--Schwarz over the packets therefore bounds this part
by $O_\Theta(L^{3/2}e^{-p_0/L})$. The second mixed part follows
by self-adjointness. Adding the deterministic deletion and
taking the norm of the difference gives
\cref{eq:asymupperproduct}.

The lower product exclusion of \cref{thm:subcritical} can be
imposed on the retained form at the same cost
$O_\varepsilon(Q^{-1/2}+N^{-1/2+\varepsilon})$ for $Q\ge1$.
Indeed, the upper mask is zero on $xy\le N^2/Q$. On $x<1$
or $y<1$ there are only model terms. The $O(N^{-1})$ bound
below one in the proof of \cref{thm:subcritical} also applies
with the upper mask inserted. For fixed $x<1$, the difference
on $[bN-x,bN+x]$ is now applied to
$z(y)(1-m(\log x,\log y))$. Its derivative is bounded by
$O_\Theta(|z'(y)|+|z(y)|/y)$, whose squared integral with
weight $y^2$ is controlled by the same energy bound.

\section{Norm decay for almost all moduli on growing windows}
\label{sec:almostall}

For $R>0$ set
\[
 E_N(R)=\frac{1}{\SSS(N)}
     \sup_{\NN(F),\NN(G)\le1}|\B(f,g)|,
\]
where both tests have physical support in $(0,R]$. The right-hand
side of \cref{eq:form} defines this quantity without GRH. When
$\grh(N)$ holds it is the normalized norm of the corresponding
restriction of $T_N$.

\begin{theorem}[Mean-square decay of the restricted norm]
\label{thm:almostall}
For every fixed $A\ge0$ and $D>0$,
\begin{equation}
 \sum_{\substack{Y\le N\le2Y\\N\ {\rm even}}}
      E_N\bigl(N(\log N)^A\bigr)^2
     \ll_{A,D}Y(\log Y)^{-D}.
 \label{eq:almostlog}
\end{equation}
Under global Dirichlet GRH, for every fixed $0<\eta<1/3$
there is $c_\eta>0$ such that
\begin{equation}
 \sum_{\substack{Y\le N\le2Y\\N\ {\rm even}}}
      E_N(N^{1+\eta})^2
     \ll_\eta Y^{1-c_\eta}.
 \label{eq:almostpower}
\end{equation}
For example, one can take $c_{6/25}=1/2000$.
\end{theorem}

For $R,S>0$ write $E_N(R,S)$ for the same normalized supremum
with the first test supported in $x\le R$ and the second in
$y\le S$.

\begin{theorem}[Mixed supports at and above the critical product]
\label{thm:mixedpower}
Assume global Dirichlet GRH. For fixed $0<a\le b$ with
$b+a/2<2$ and fixed $C_1,C_2>0$, there is $c_{a,b}>0$ with
\begin{equation}
 \sum_{\substack{Y\le N\le2Y\\N\ \mathrm{even}}}
 E_N(C_1N^a,C_2N^b)^2 \ll_{a,b,C_1,C_2}Y^{1-c_{a,b}}.
 \label{eq:generalrectpowers}
\end{equation}
In particular, for $0<a<1$ one can take any fixed $b<2-a/2$
with $b\ge a$. In the narrower range $b=2-a+\kappa$, with
$\kappa=\min(a,1-a)/100$, one can take $c_{a,b}=\kappa/2$.
In this case
\begin{equation}
 \sum_{\substack{Y\le N\le2Y\\N\ {\rm even}}}
    E_N(C_1N^a,C_2N^{2-a+\kappa})^2
       \ll_{a,C_1,C_2}Y^{1-\kappa/2}.
 \label{eq:mixedpower}
\end{equation}
Thus, for each fixed $a$, the normalized mixed norm is at most
$Y^{-\kappa/8}$ outside $O_{a,C_1,C_2}(Y^{1-\kappa/4})$
even moduli in $[Y,2Y]$.
\end{theorem}

For fixed $B\ge0$, let $E_N^{\le}(B)$ be the normalized
norm of \cref{eq:form} retaining only pairs with
\begin{equation}
  xy\le R_N^*,\qquad R_N^*=N^2(\log N)^B ,
 \label{eq:productmask}
\end{equation}
in both the arithmetic and continuous terms. Here
$x,y$ are the physical variables. Write $E_{N,\ge M}^{\le}(B)$
for the restriction to tests supported in $x,y\ge M$.

\begin{theorem}[The product range above polylogarithmic inputs]
\label{thm:productwindow}
Assume global Dirichlet GRH. For every fixed $B\ge0$ and $D>0$
there is $A>0$ such that
\begin{equation}
 \sum_{\substack{Y\le N\le 2Y\\N\ \mathrm{even}}}
 \bigl(E_{N,\ge(\log N)^A}^{\le}(B)\bigr)^2
       \ll_{B,D}Y(\log Y)^{-D}.
 \label{eq:productwindow}
\end{equation}
\end{theorem}

\begin{theorem}[A power range above the critical product]
\label{thm:powerproduct}
Assume global Dirichlet GRH and fix $0<\delta<1$ and
$0<\gamma<\delta/2$. Let $E_N(\delta,\gamma)$ be the
normalized norm on tests supported in $x,y\ge N^\delta$
after multiplying the pair weights in both terms of
\cref{eq:form} by $\Phi(\log(xy/N^{2+\gamma}))$, with a fixed
smooth $\Phi\colon\mathbb R\to[0,1]$ equal to one on
$(-\infty,0]$ and zero on $[1,\infty)$. There is
$c_{\delta,\gamma}>0$ such that
\begin{equation}
 \sum_{\substack{Y\le N\le2Y\\ N\ \mathrm{even}}}
 E_N(\delta,\gamma)^2
     \ll_{\delta,\gamma,\Phi}Y^{1-c_{\delta,\gamma}} .
 \label{eq:powerproduct}
\end{equation}
This also bounds the mean square of the supremum of the
original normalized restriction norms over rectangles lying
in $x,y\ge N^\delta$ and $xy\le N^{2+\gamma}$.
The same power-saving mean-square conclusion holds for the
sharp product cutoff $\mathbf 1_{xy\le N^{2+\gamma}}$ in place
of the smooth multiplier, with a possibly different exponent.
\end{theorem}

\subsection{Mean squares of interval-restricted pair counts}

For intervals $I,J\subset[0,P]$, a sign $\sigma\in\{1,-1\}$,
and an integer $h\ge1$, put
\[
 R_\sigma(h;I,J)=
   \sum_{\substack{n\in I,\ m\in J\\m+\sigma n=h}}
                  \Lambda(n)\Lambda(m),\qquad
 L_\sigma(h;I,J)=\int_I\mathbf1_J(h-\sigma x)\,\ud x .
\]
The sums are over positive integers. Write
$\Delta_\sigma(h;I,J)=R_\sigma(h;I,J)-\SSS(h)L_\sigma(h;I,J)$,
extending $\SSS(h)$ by zero for odd $h$. Without GRH, uniformly
in $I,J,\sigma$, we have
\begin{equation}
 \sum_{1\le h\le2P}|\Delta_\sigma(h;I,J)|^2
      \ll_D P^3(\log P)^{-D}
       \qquad\text{for every $D>0$}.
 \label{eq:pairlog}
\end{equation}

Write $e(t)=e^{2\pi it}$,
$S_I(\alpha)=\sum_{n\in I}\Lambda(n)e(n\alpha)$ and
$v_I(\beta)=\int_I e(x\beta)\,\ud x$. Integrating
$S_J(\alpha)S_I(\sigma\alpha)e(-h\alpha)$ over a period
gives $R_\sigma(h;I,J)$. Fourier inversion on the real line
gives $L_\sigma$ by replacing $S_I,S_J$ with $v_I,v_J$.

Use disjoint major arcs $|\alpha-a/q|\le2Q/(qP)$ for reduced
$a/q$ and $q\le Q\le\sqrt P/4$. Dirichlet approximation on the
minor arcs gives $Q<q\le P/Q$ and $|\alpha-a/q|\le q^{-2}$.
Vaughan's exponential sum bound \cite{VaughanExp}, applied to
both initial sums defining $S_I$, gives
\[
 \sup_{\rm minor}|S_I(\alpha)|
     \ll (P Q^{-1/2}+P^{4/5})(\log P)^4.
\]
For endpoints of an initial sum below $P/\sqrt Q$ the Chebyshev
bound already suffices. For larger endpoints the approximating
denominator is at most the endpoint. By Parseval,
$\int_0^1|S_I(\alpha)|^2\,\ud\alpha\ll P\log P$.
Bessel's inequality for the product restricted to the minor arcs
then bounds their contribution to the mean square by
\begin{equation}
         O\bigl((P^3/Q+P^{13/5})(\log P)^9\bigr).
 \label{eq:almostminor}
\end{equation}

On the major arc at $a/q$ the main term for $S_I(a/q+\beta)$
is $\mu(q)v_I(\beta)/\varphi(q)$, and likewise for $J$.
The coefficient of the product of these
main terms after summing over $a$ is
$A_q(h)=\mu(q)^2c_q(h)/\varphi(q)^2$ with $c_q$ the Ramanujan
sum. The series $\sum_q A_q(h)$ is $\SSS(h)$ by its Euler
product for $h\ne0$. Completing each real integral past
$|\beta|=2Q/(qP)$ costs at most $O(Pq/Q)$ since
$|v_I(\beta)|\ll\min(P,|\beta|^{-1})$. Thus the total main-term
error from completing the integrals and the singular series is
bounded by
\begin{equation}
  C P\sum_{q\ge1}\min(1,q/Q)|A_q(h)|.
  \label{eq:almosttails}
\end{equation}
Here the untruncated sum of $|A_q(h)|$ is $O(h/\varphi(h))$
by the product over squarefree $q$, and hence $O(\log P)$ on
$1\le h\le2P$. Also, the squarefree Ramanujan sum satisfies
$|c_q(h)|\le\sum_{d\mid(q,h)}d$, so
$\sum_{1\le h\le2P}|c_q(h)|\le2P\tau(q)$, where $\tau$ counts
divisors. The bound $\tau(q)/\varphi(q)^2\ll_\varepsilon
q^{-2+\varepsilon}$ follows from the Euler factors. Consequently
the mean square sum of \cref{eq:almosttails} is
$O_\varepsilon(P^3(\log P)Q^{-1+\varepsilon})$.

For \cref{eq:pairlog}, set $Q=(\log P)^K$ with $K$ large
in terms of $D$. Siegel--Walfisz \cite{Walfisz} and partial
summation give a major-arc error for $S_I$ of
$O_{K,B}(P(\log P)^{-B})$ for any $B>0$, uniformly in both
endpoints of $I$. At endpoints below $P^{1/2}$ one may use the
Chebyshev bound, and above it $q\le (\log P)^K$ is in the
Siegel--Walfisz range. The absolute error from integrating the two
products over the major arcs is then $O_D(P(\log P)^{-D})$
by taking $B$ large, since their total length is $O(Q^2/P)$.
Together with \cref{eq:almostminor,eq:almosttails} this proves
\cref{eq:pairlog}.

\subsection{An integrated major-arc error under GRH}

Let $J$ be an interval in $[Z/e,eZ]$, with $2\le Z\le P$,
and let $q\le Z/(4e)$. As above, for a reduced $a/q$ write
$E_{J,a}(\beta)=S_J(a/q+\beta)-\mu(q)v_J(\beta)/\varphi(q)$.
For $0<B\le1/2$ we will use
\begin{equation}
 \int_{-B}^B\sum_{a\bmod q}^{*}|E_{J,a}(\beta)|^2\,\ud\beta
       \ll qZ B(\log(2P))^4 .
 \label{eq:arcintegrated}
\end{equation}
The same bound holds when the interval indicator is replaced
by a tent weight of height one. In that case $S_J$ and $v_J$
use the same weight.

To prove the estimate set $M_\chi(t)=\psi(t,\chi)-\delta_\chi t$,
where $\delta_\chi$ is one for the principal character and zero
otherwise. On a fixed interval of scales $t\asymp Z$, under GRH,
the truncated explicit formula writes $M_\chi(t)$ as
$-\sum_{|\gamma|\le T}t^{1/2+i\gamma}/(1/2+i\gamma)$
with an $O((\log(2P))^2)$ error for $T=cZ$ and a sufficiently
small fixed $c>0$. Use the formula for the inducing primitive
character and account for the prime powers at primes dividing
$q$ \cite[Lemmas~1--2]{KLM}. For $1\le H\le Z$, split the
sum at $U=Z/H$. Below $U$, express the difference between
the values at $t+H$ and $t$ by integrating
$u^{-1/2}\sum_{|\gamma|\le\min(U,T)}u^{i\gamma}$
over $[t,t+H]$. Cauchy--Schwarz and $u=e^r$ bound its
mean square in $t$ by $H^2$ times the integral of the
squared exponential sum over a log interval of bounded length.
For that integral, use a nonnegative majorant with bounded
Fourier support as in \cref{sec:zeros}. The unit-bin zero
count is $O(\log(2P))$ up to $T$, so the integral is
$O(U(\log(2P))^2)$. For the part above $U$ use the two
endpoint sums separately. Their mean square is at most
$O(Z^2(\log(2P))^2/U)$ by the same bin argument with
coefficients $(1/2+i\gamma)^{-1}$.
Together with the truncation error this gives
\[
 \int_{t\asymp Z}|M_\chi(t+H)-M_\chi(t)|^2\,\ud t
       \ll ZH(\log(2P))^4 .
\]
The integral is over any fixed range of comparable scales.

Now take $H=1/(2B)$ and let $\mu_\chi$ be the measure
$\sum_{n\in J}\Lambda(n)\chi(n)\delta_n-\delta_\chi\mathbf1_J\,\ud x$.
Plancherel after convolution with an interval of length $H$
bounds $\int_{-B}^B|\widehat\mu_\chi(\beta)|^2\,\ud\beta$
by a constant times
$H^{-2}\int_{\mathbb R}|\mu_\chi((t,t+H])|^2\,\ud t$,
since the Fourier transform of the interval has size
$\gg H$ on $[-B,B]$. If $H\le Z$, intervals lying strictly
inside $J$ are covered by the preceding mean square.
Intervals meeting an endpoint of $J$ cost at most
$O(ZH(\log(2P))^4)$ by the GRH pointwise bound for
$M_\chi$ and the length $O(H)$ of the endpoint ranges.
For $H\ge Z$ this pointwise argument bounds the integral
on its whole support, which has length $O(H)$.
Thus the Fourier integral is $O(ZB(\log(2P))^4)$.
Multiplicative Plancherel gives the sum of the squared
progression errors on units, and additive Plancherel
multiplies that sum by $q$. The nonunit prime powers
have total weight $O((\log(2P))^2)$ and are absorbed.
The additive transform of the unit main terms at reduced
$a/q$ is $\mu(q)v_J(\beta)/\varphi(q)$, giving
\cref{eq:arcintegrated}. A tent is the integral of its
superlevel interval indicators. Jensen's inequality gives
the tent version.

\subsection{Pair-count means at unequal scales}

Take $1\le X\le V$, put $P=eV$, and let $\mathcal I,\mathcal J$
be collections of at most $K$ disjoint subintervals of
$[X/e,eX]$ and $[V/e,eV]$, respectively. The interval lengths
are bounded by constant multiples of $W_x=X/K$ and
$W_y=P/K$ in the respective collections. Write
\[
 \mathcal R(X,V;K,Q)=
   \sum_{\substack{I\in\mathcal I,\ J\in\mathcal J\\
                   \sigma\in\{1,-1\}}}
      \sum_{1\le h\le2P}|\Delta_\sigma(h;I,J)|^2 .
\]
Assume $2\le Q\le\min(\sqrt P/4,X/(4e))$ and $W_x Q\ge P$.
Under global
Dirichlet GRH,
\begin{equation}
 \mathcal R(X,V;K,Q)\ll_\varepsilon P^\varepsilon
       \left(\frac{KX P^2}{Q}+KX P^{8/5}
            +KXP+K^2 X Q^3\right).
 \label{eq:pairpower}
\end{equation}
The displayed right side is uniform over all such collections.

For the minor arcs, sum the Bessel bound on both families
before estimating it. Bound $\sum_{J\in\mathcal J}|S_J|^2$
by its minor-arc supremum $O(K(P^2/Q+P^{8/5})(\log P)^8)$.
Parseval and disjointness give
$\int_0^1\sum_{I\in\mathcal I}|S_I|^2\ll X\log P$.
This proves the first two bounds on the right of
\cref{eq:pairpower} for the minor-arc contribution.

For completion of the singular integrals, the integral of
$|v_I(\beta)v_J(\beta)|$ over the real line is
$O(W_x\log(2P))$, by using $\min(W_x,|\beta|^{-1})$ and
$\min(W_y,|\beta|^{-1})$ for the respective factors.
Past $2Q/(qP)$ its integral is $O(Pq/Q)$.
Thus the completion error is at most
\[
 C W_x\log(2P)\sum_q \min(1,qP/(W_x Q))|A_q(h)|.
\]
The Ramanujan sum argument for \cref{eq:almosttails} gives
$O_\varepsilon(P^\varepsilon P^2 W_x/Q)$ in mean square per
interval pair, since $W_xQ/P\ge1$. Multiplying by $K^2$
gives the first term of \cref{eq:pairpower}.

For a reduced $a/q$ with $q\le X/(4e)$, put
$E_{I,a}(\beta)=S_I(a/q+\beta)-\mu(q)v_I(\beta)/\varphi(q)$.
On the smaller scale we have
\begin{equation}
  \sum_{a\bmod q}^{*}|E_{I,a}(\beta)|^2
       \ll qX(1+W_x|\beta|)^2(\log P)^4 .
  \label{eq:arcmeansquare}
\end{equation}
After subtracting the principal term, each initial portion
of $I$ has twisted Chebyshev error $O(\sqrt X(\log P)^2)$
under GRH \cite[Section~4]{LeeGRH}.
Partial summation introduces $1+W_x|\beta|$.
Multiplicative Plancherel bounds the sum of squared errors
in unit progressions by $O(X(1+W_x|\beta|)^2(\log P)^4)$.
The total weight of the nonunit prime powers is
$O((\log P)^2)$. Additive Plancherel over the numerators
multiplies the squared progression norm by $q$.
For reduced numerators the main term transforms to
$\mu(q)v_I(\beta)/\varphi(q)$ as required. This argument does
not require a restriction on $\beta$.

Apply Bessel to the product error on the disjoint major arcs.
For one main term and one error, use \cref{eq:arcintegrated}
against $|v_I(\beta)|^2$ or $|v_J(\beta)|^2$ on dyadic
frequency ranges starting at the reciprocal cell width.
This costs, up to logarithms, $qP W_x/\varphi(q)^2$ or
$qX W_y/\varphi(q)^2$ per denominator. In fact an integrated
error of size $O(qZB)$ on $|\beta|\le B$ with the bound
$|v(\beta)|\ll W/(1+W|\beta|)$ gives $O(qZW)$ by geometric
summation. Summing over $q$ and multiplying by $K^2$ gives
$O_\varepsilon(P^\varepsilon KXP)$.
For the product of two errors, use \cref{eq:arcmeansquare}
on the small scale and \cref{eq:arcintegrated} on the large
scale. On an arc of width $B=2Q/(qP)$ the integrated square
summed over numerators is, up to logarithms,
$O(q^2 XP B(1+W_x B)^2)$.
Since $W_x\ll P/K$, summing over $q\le Q$ bounds this
by $O(XQ^3\log(2Q))$. Multiplying by $K^2$ proves the last
term of \cref{eq:pairpower}.

\subsection{The unconditional logarithmic window}

For $Y\le N\le 2Y$ let $NH$ be a common upper bound for the
desired physical support, with $H\ge1$ depending on $Y$.
First use a smooth lower cutoff from $N/(2M)$ to $N/M$ with
$M\ge H$. Multiplication by this cutoff and its complement
is uniformly bounded in $\NN$ because their transition has
fixed log width. Terms with one lower-supported test lie
entirely in $xy\le N^2H/M$. The arithmetic and model bounds
in \cref{sec:subcritical} apply directly to \cref{eq:form}, with no
GRH needed for this step. The normalized error of discarding all
such terms is $O_\varepsilon((H/M)^{1/2}+Y^{-1/2+\varepsilon})$.

Partition $[0,P)$ into equal half-open cells of length
$1\le W\le Y/4$, where $P$ is the smallest multiple of $W$
strictly greater than $2YH$. After the lower cutoff, the
scaled tests $f,g$ in $u=x/N$ have supremum norm
$O(M^{1/2})$, $L^2$ norm $O(1)$ and derivative norm $O(M)$.
This follows from $f(u)=u^{-1/2}F(\log(Nu))$, the lower bound
$u\ge1/(2M)$ and \cref{eq:logenergy}. In the ranges used below
$N/(2M)>1$, so $W_N(\log(Nu))\gg1$ and the log $H^1$ norms
of the tests are bounded.

Replace the scaled tests by their cell averages. Poincar\'e's
inequality gives an $L^2(\ud u)$ error $O(WM/N)$. For the
arithmetic sum, write $a_I$ for the cell average of $f(x/N)$
on $I$. The squared pointwise error on the cell is bounded
by $W\int_I|\frac{\ud}{\ud x}f(x/N)|^2\,\ud x$. Each cell
contains at most $W+1$ integers, so the error as a sequence
over integer arguments has $\ell^2$ norm $O(WM/\sqrt N)$.
The step values have $\ell^2$ norm $O(\sqrt N)$ by contractivity
of averaging in $L^2$ and the same bound on integers per cell.
So do the original samples, since we take $WM\le Y$.
The unsigned arithmetic matrix has row and column sums
$O(H(\log P)^2)$, since the other integer lies in one of
two residue classes and occurs at most $O(H)$ times.
For the continuous model, each translation or reflection has
$L^2$ pairing norm at most one, and
$\sum_{b\le 2P/N}\SSS(bN)/\SSS(N)\ll H$ by \cref{eq:subseries}.
Together with $1/N$ in the arithmetic normalization, these
bounds give a replacement error $O(HMW Y^{-1}(\log P)^2)$.

The step tests can use the unrestricted counts $R_\sigma$.
For any one line $h=bN$, the unit condition removes at most
$O((\log P)^2)$ pairs, since each removed prime power is over a
prime dividing $N$ and the other integer is determined. This
also accounts for the equal-integer contributions on the sum
lines excluded in \cref{eq:form}, since they are nonunits.
The error after weights and normalization is
$O(Y^{-1} H M(\log P)^4)$.

Let $\mathcal I$ be the set of cells, $K=\lfloor2P/Y\rfloor$,
and use $1\le b\le K$ for the possible shifts. Cell averaging
is contractive in $L^2$, and $\|f\|_{L^2(\ud u)}=\|F\|_{L^2(\ud r)}$.
Thus the $\ell^2$ norm of the cell values of either unit test is
$O((N/W)^{1/2})$. For each cell pair and each orientation the
possible positive shifts lie in an interval of length at most
$2W<N$. Thus only one $b$ can contribute to that entry per
orientation. The Frobenius bound for the matrix of cell-pair
errors now gives a normalized step-form remainder at most
\[
 \frac{C}{W}\left(\sum_{\substack{I,J\in\mathcal I,\ 1\le b\le K\\
                                  \sigma\in\{1,-1\}}}
                |\Delta_\sigma(bN;I,J)|^2\right)^{1/2}.
\]
The other difference orientation interchanges $I$ and $J$.
In summing over $N$ and $b$, a fixed $h=bN\le2P$ occurs
at most $K$ times by injectivity for each $b$.
Writing $\mathcal R_W(P)$ for the sum of squared
$\Delta_\sigma(h;I,J)$ over the cell pairs, shifts $h$ and signs,
the step remainder has mean square at most
\begin{equation}
 O\bigl(W^{-2}K\,\mathcal R_W(P)\bigr).
 \label{eq:almostmatrix}
\end{equation}

For \cref{eq:almostlog}, take $H=(\log(2Y))^A$,
$M=(\log Y)^B$ and $W=Y(\log Y)^{-C}$. Choose $B$ and then $C$
large enough in terms of $A,D$ that the deterministic errors
are $O((\log Y)^{-D})$. Apply \cref{eq:pairlog} to each of the
$(P/W)^2$ cell pairs with a sufficiently large log saving
in \cref{eq:almostmatrix}.

\subsection{Separate interpolation meshes for the power window}

Assume global Dirichlet GRH. For $0<\eta<1/4$ we first prove
\cref{eq:almostpower} with an explicit choice of parameters.
use the lower cutoff above with $H=(2Y)^\eta$ and
$M=Y^{\eta+\rho}$ where $\rho>0$. Its normalized error is
$O(Y^{-\rho/2}+Y^{-2/5})$. Use the unit log packets of
\cref{sec:cutoff} on the central part. Over $N\in[Y,2Y]$
one fixed list of $O(\log Y)$ packet indices suffices.
For a pair with smaller and larger physical scales $X,V$,
we have
\begin{equation}
 X\gg Y^{1-\eta-\rho},\qquad V\ll Y^{1+\eta}.
 \label{eq:adaptivewindow}
\end{equation}
Pairs with $2eV<Y$ contribute zero, because both $x+y$ and
$|x-y|$ are then below $N$. For the others put $P=eV$,
so $P\gg Y$.

Subdivide each physical packet $[X/e,eX]$ and $[V/e,eV]$
into $K=\lceil Y^\omega\rceil$ equal cells, with $\omega>0$.
For the physical weight $w(x)=x^{-1/2}F(\log x)$ interpolate
its nodal values linearly in $x$, and set $F_{\rm lin}(r)$
to be $e^{r/2}$ times that interpolant at $e^r$.
It stays on the same packet and is bounded in $\NN$ by a
constant times $\NN(F)$. To see stability, on a cell the
derivative of the interpolant is the mean of $w'$.
Here $x\asymp X$ and the $\NN$ norm is equivalent to the
unweighted $H^1(\ud r)$ norm. In physical coordinates this
last norm is equivalent to
$(\|w\|_2^2+\|xw'\|_2^2)^{1/2}$. Cauchy--Schwarz for
the mean derivative proves the derivative bound, and
Poincar\'e's inequality for the residual proves the $L^2$
bound. The same argument works on the other packet.

The log residual $R=F-F_{\rm lin}$ vanishes at each node and
has $L^2$ and supremum bounds $O(K^{-1}\NN(F))$ and
$O(K^{-1/2}\NN(F))$, since each cell has log length $O(K^{-1})$.
Apply \cref{eq:zeroslow,eq:zeroshigh} with $T=K$ on this log
packet. The bound for the trivial zeros and the imprimitive
correction in \cref{eq:poslogexplicit} then gives
$\|U(R)\|\ll L K^{-1/2}\NN(F)$. The other character factor
has norm $O(L)$ on unit tests by \cref{eq:Utail}.
The diagonal costs $O(LK^{-1/2})$ by the supremum bound
on the residual. For the model, use the lattice estimate
of \cref{eq:Mpacket} with the $L^1$ bound on the physical
residual, as in the proof of \cref{eq:meshDM}. If the
log packet centers differ by $d$, the coefficient is bounded
by $O(L\min(e^{-d/2},K^{-1}e^{d/2}))$ since
$\log V=O(L)$. This is $O(LK^{-1/2})$.
Together these bounds give a normalized interpolation error
$O(L^2 K^{-1/2})$ on unit tests for every packet pair.

Each cell interpolant is a sum of an increasing and a
decreasing linear ramp times its endpoint values.
The $\ell^2$ norm of the nodal values of the physical
interpolant on the smaller packet is $O(W_x^{-1/2})$ on unit
tests, where $W_x=X/K$. Indeed,
$\int_0^1|(1-t)a+tb|^2\,\ud t\ge(|a|^2+|b|^2)/6$,
and the physical $L^2$ norm of the interpolant is bounded by
stability. On the other packet the bound is $O(W_y^{-1/2})$
with $W_y=P/K$. The shift range for each cell pair and sign
has length $O(P/K)$ and meets at most $O(1+P/(YK))$ positive
multiples $bN$. The two difference orientations interchange
the cell roles.

The unrestricted pair counts can again be used, since at
most $O((\log P)^2)$ pairs per line are removed by the unit
condition, which also excludes the equal-integer sum pairs.
With physical supremum weights $O(X^{-1/2})$ and
$O(V^{-1/2})$, the normalized error over the $O(P/Y)$ shifts
is $O(L^4Y^{-1}\sqrt{V/X})$.
The weighted $\Delta_\sigma$ for two cell ramps is an average
of the unweighted $\Delta_\sigma$ over two subintervals of
the cells. For example the increasing ramp is the integral,
over $t\in[0,1]$, of the indicator of the part of its cell
to the right of the fraction $t$. Jensen's inequality thus
preserves \cref{eq:pairpower} for the sum of squared weighted
errors of the four ramp choices. Each fixed pair of
integration parameters gives disjoint collections of
subintervals, to which that estimate applies directly.

Use the Frobenius bound with the two nodal $\ell^2$ norms
and Cauchy--Schwarz over the shifts contributing to an entry.
When the resulting squared error is summed over $N$, a
shift $h=bN$ occurs at most $\tau(h)\ll_\varepsilon P^\varepsilon$
times. Hence its mean square on the modulus interval is
bounded by the right side of \cref{eq:pairpower} times
$O_\varepsilon(P^\varepsilon(1+P/(YK))/(W_xW_y))$.
We may use $1+P/(YK)\ll P/Y$ since $P\gg Y$. Up to
arbitrarily small powers of $P$, the resulting terms are
\begin{equation}
 \frac{K^3 P^2}{YQ},\qquad \frac{K^3P^{8/5}}{Y},
 \qquad \frac{K^3P}{Y},\qquad \frac{K^4Q^3}{Y}.
 \label{eq:adaptiveterms}
\end{equation}

For $0<\eta<1/4$ choose $\theta$ with
\[
 \frac{2\eta}{1+\eta}<\theta<\tfrac12
\]
and take $Q=\lfloor P^\theta\rfloor$.
Using $P\ll Y^{1+\eta}$, the exponents in
\cref{eq:adaptiveterms} are bounded above by
\begin{equation}
 3\omega+(2-\theta)(1+\eta)-1,\quad
 3\omega+\tfrac85(1+\eta)-1,\quad
 3\omega+\eta,\quad 4\omega+3\theta(1+\eta)-1.
 \label{eq:adaptivexponents}
\end{equation}
All are strictly below one for sufficiently small $\omega>0$.
The bound $Q\le\sqrt P/4$ holds as $\theta<1/2$, and
$Q\le X/(4e)$ by \cref{eq:adaptivewindow} since
$\theta(1+\eta)<1-\eta-\rho$ for small $\rho$.
The condition $W_x Q\ge P$ follows from
$\theta(1+\eta)>2\eta+\rho+\omega$, again taking $\rho$
and $\omega$ sufficiently small.

Finally, summing over the $O((\log Y)^2)$ packet pairs and
using Cauchy--Schwarz costs only a logarithmic factor in
the mean square. The interpolation and unit-condition
errors have power savings by their uniform bounds.
The lower-cutoff error contributes $O(Y^{1-\rho})$ when
squared and summed over the moduli. This proves
\cref{eq:almostpower}. At $\eta=6/25$ we may take
$\omega=\rho=1/1000$ and $\theta=21/50$. The largest
exponent in \cref{eq:adaptivexponents} is $0.987$.
The deterministic errors have squared mean
$O(Y^{1-1/1000}(\log Y)^{O(1)})$, giving
$c_{6/25}=1/2000$.

\subsection{Signed counts for mixed supports through the critical product}

For the quantitative special case \eqref{eq:mixedpower},
remove the first support below
$N^{a-2\kappa}$ with a smooth cutoff of fixed log width
whose transition ends there. In the discarded terms we have
$xy\ll N^{2-\kappa}$, so \cref{thm:subcritical} costs
$O(Y^{-\kappa/2}+Y^{-2/5})$ on the normalized scale.
The part of the second support below $N/2$ gives no term
for large $N$, since the first variable is $o(N)$.
Use unit log packets and the separate linear interpolants
of the preceding subsection, now with
$K=\lceil Y^\kappa\rceil$. For each surviving packet pair
we use scales $X,V$ and $P=eV$ with
\begin{equation}
 Y^{a-2\kappa}\ll X\ll Y^a,\qquad
 Y\ll P\ll Y^{2-a+\kappa}.
 \label{eq:mixedscales}
\end{equation}
The interpolation error is again $O(L^2 K^{-1/2})$ on
unit tests. Both scales are positive powers of $Y$, and
the proof uses only the log length and this scale bound.
The unit restriction can be removed at cost
$O(L^4Y^{-1}\sqrt{P/X})$, as before.
The continuous model on the packet pair has normalized
norm $O(L\sqrt{X/P})$ by \cref{eq:Mpacket}, which has a
power saving since $a<1$.

Write the two linear interpolants in tent bases
$u_i,v_j$ on the two grids. The endpoint values vanish,
and each internal hat has height one and support on two
cells. There are $O(K)$ tents in each family, of widths
$O(W_x)$ and $O(W_y)$ for $W_x=X/K$ and $W_y=P/K$.
Use $S_{u_i},v_{u_i}$ and $S_{v_j},v_{v_j}$ for the
exponential sums and integrals with these weights. Since
$m\gg n$ on the two packets, the signed arithmetic count
at $h>0$ is the coefficient $D_{ij}(h)$ of
\[
       \bigl(S_{u_i}(-\alpha)-S_{u_i}(\alpha)\bigr)S_{v_j}(\alpha).
\]
The two signs correspond to $m-n=h$ and $m+n=h$.
Use major arcs at $a'/q$ for $q\le Q=\lfloor Y^{a/10}\rfloor$
with radius $2Q/(qX)$. They are disjoint since
$Q\le\sqrt X/4$ for large $Y$. Denote the Fourier coefficient
of the product of the major-arc main terms by $M_{ij}(h)$.

On the minor arcs, choose a Dirichlet approximation $a/q+\beta$
with $q\le X/Q$ and $|\beta|\le 2Q/(qX)$. This puts $q>Q$.
For $Q\ge 4e$ we have $q\le X/(4e)$ and can apply
\cref{eq:arcmeansquare} at this approximant. Its bound on
one error is $O(\sqrt{qX}(1+X|\beta|)(\log P)^2)$.
Since $Q<q\le X/Q$, this is $O(XQ^{-1/2}(\log P)^2)$.
The main term at this approximant is at most $O(X/\sqrt q)$
by $\varphi(q)\gg\sqrt q$. Thus the bound
$O(XQ^{-1/2}(\log P)^2)$ holds for every interval in the
smaller packet and for its tents by interval superposition.
Bounded $Q$ can be treated by Chebyshev's estimate.
Parseval for the larger variable summed over its tents costs
$O(P\log P)$, since their overlaps have bounded multiplicity.
The pointwise bound \cref{eq:arcmeansquare} also extends
to tents by interval superposition. For a single error on
the major arcs, integrate \cref{eq:arcintegrated} against
the square of the main term. Splitting $|\beta|$ at the
reciprocal tent width and then into dyadic ranges costs at
most $q X W_y/\varphi(q)^2$ or $q P W_x/\varphi(q)^2$ per
denominator and tent pair, up to logarithms. Since
$\sum_{q\le Q}q/\varphi(q)^2\ll(\log(2Q))^3$, this gives
$KXP$ times a fixed power of $\log(2P)$ after summing over
tent pairs. For two errors, combine
\cref{eq:arcmeansquare} for the small variable and
\cref{eq:arcintegrated} for the large one, using
$B=2Q/(qX)$. On summing over the numerators and integrating,
the cost for each denominator is, up to logs,
$O(qP Q(1+Q/(qK))^2)$ per tent pair.
Bessel over the minor and major arcs therefore gives, for
an absolute constant $C_0$,
\begin{equation}
 \sum_{i,j}\sum_{1\le h\le2P}|D_{ij}(h)-M_{ij}(h)|^2
 \ll (\log(2P))^{C_0}
   \left(\frac{KP X^2}{Q}+KXP+K^2P Q^3\right).
 \label{eq:mixedarcerror}
\end{equation}
The proof of \cref{eq:mixedarcerror} also applies to each
product $S_{u_i}(\alpha)S_{v_j}(\alpha)$,
$S_{u_i}(-\alpha)S_{v_j}(\alpha)$ and
$S_{u_i}(\alpha)S_{v_j}(-\alpha)$ separately, for
$X\le V$, $P=eV$ and $2\le Q\le\sqrt X/4$.
In each case $D_{ij}(h)$ and $M_{ij}(h)$ are the corresponding
arithmetic coefficient and the coefficient of the corresponding
product of major-arc main terms. The minor-arc estimate uses
the bound on the smaller variable, independent of the signs,
and the major-arc estimates above hold for either ordering.

In the major-arc main term, the sign produces the factor
$v_{u_i}(-\beta)-v_{u_i}(\beta)$, of size at most
$O(W_x\min(1,X|\beta|))$. The Fourier transform of the
larger tent, being a constant multiple of the transform
of a convolution of two intervals, is
$O(W_y(1+W_y|\beta|)^{-2})$. Since $W_y\gg X$ by
\cref{eq:mixedscales} and $\kappa<1-a$, the integral of
the product of these bounds is $O(W_x X\log(2P)/W_y)$.
The reduced numerators give the coefficient $A_q(h)$
from the earlier major-arc calculation, and
$\sum_q |A_q(h)|\ll\log(2P)$ for $1\le h\le2P$.
Consequently,
\[
 \sup_{1\le h\le2P}|M_{ij}(h)|
           \ll (W_x X/W_y)(\log(2P))^2 .
\]
The nodal $\ell^1$ norms are $O(K/\sqrt X)$ and
$O(K/\sqrt P)$ by the earlier $\ell^2$ estimates.
Summing the major-arc main term over tent pairs and the
$O(P/Y)$ shifts $h=bN$ therefore costs at most
$O(K^2 X^{3/2}(\log(2P))^2/(Y\sqrt P))$ on unit tests.
This too has a power saving throughout \cref{eq:mixedscales}.

For each tent pair, only the shifts that intersect one
of the two lines through its supports can contribute to
$D_{ij}$. Restrict to those shifts before splitting off
$M_{ij}$. The Frobenius and divisor estimates of the
preceding subsection bound the mean square of the
remaining form by \cref{eq:mixedarcerror} times
$O_\varepsilon(P^\varepsilon(P/Y)/(W_xW_y))$.
Here we have used the bound $O(P/Y)$ for the number of
shifts per entry. Up to arbitrarily small powers of $P$,
the resulting terms are
\[
 \frac{K^3XP}{YQ},\qquad
 \frac{K^3P}{Y},\qquad \frac{K^4P Q^3}{XY}.
\]
By \cref{eq:mixedscales} their exponents in $Y$ are at most
$1+4\kappa-a/10$, $1-a+4\kappa$, and
$1-17a/10+7\kappa$, respectively. Each is below $1-\kappa$
by the choice of $\kappa$. The uniform deterministic
errors, including interpolation and the product cutoff,
have squared mean $O(Y^{1-\kappa}(\log Y)^{O(1)})$.
Summing over the $O((\log Y)^2)$ packet pairs costs only
logarithms, proving \cref{eq:mixedpower}.

\subsection{Summation over the product range}

Write $\lambda=\log Y$ and put $\epsilon=(\log N)^{-J}$
with a fixed $J>B+D+10$. Use an auxiliary weight
$\Phi(\epsilon^{-1}\log(xy/R_N^*))$, where $\Phi$ is a fixed
smooth function from $\mathbb R$ to $[0,1]$ equal to one
on $(-\infty,0]$ and zero on $[1,\infty)$.
Localize both tests with the bounded-overlap
partition into unit log packets. On a pair of packets centered at
$j,k$, the auxiliary weight is a function of
$t=r+s-j-k$. Multiply that function by a fixed smooth bump equal
to one on $[-2,2]$ and apply Fourier inversion. The resulting
Fourier transform has $L^1$ norm $O(\epsilon^{-4})$ with weight
$(1+|\xi|)^2$, by four integrations by parts for $|\xi|\ge1$.
On writing the integrand as a scalar times
$e^{i\xi r}e^{i\xi s}$, the two phase multipliers on the tests
cost $O((1+|\xi|)^2)$ in $\NN$ by the product rule. Thus the
norm of the auxiliary form on each packet pair is
$O(\epsilon^{-4})$ times
the norm of the ordinary form on the same packet pair.

There are $O(\lambda^2)$ relevant pairs. In fact, both physical
scales are at least a constant times $\lambda^A$, and the
product cap requires $XP\ll Y^2\lambda^B$ for the smaller
scale $X$ and the larger upper endpoint $P=eV$. To compare the
auxiliary weight with the sharp cutoff, restrict to the band
$R_N^*\le xy\le e^\epsilon R_N^*$. On a congruence line
$y=\pm x+c$, the product $xy$ is a quadratic in $x$.
For $P$ larger than a sufficiently large constant times $X$,
its derivative has magnitude $\asymp P$ on the packets,
so the band has length $O(\epsilon X)$ on the line.
When $P\asymp X$, the preimage of an interval of length
$O(\epsilon X^2)$ under a quadratic of leading coefficient
$\pm1$ is a union of at most two intervals with total length
$O(\sqrt{\epsilon}X)$. Thus each line has at most
$O(\sqrt{\epsilon}X+1)$ contributing integers in the band.
Using the physical supremum bounds, the von Mangoldt
weights $O(\lambda^2)$, and \cref{eq:subseries} for the
continuous terms, the normalized difference on a packet
pair is at most
$O(\lambda^{2+B/2}(\sqrt\epsilon+X^{-1}))$.
Even after summing the packet bounds this is $O(\lambda^{-D/2})$
for the choice of $J$ and sufficiently large $A$.

We now bound the ordinary packet norms that enter the
Fourier estimate. Pairs with
$2P<Y$ contribute zero. If $X>Y^{9/10}$ then both packet
supports lie below $N^{1+1/8}$ for large $Y$ and their norms
are bounded by the norm in \cref{eq:almostpower}. It remains
to treat
\begin{equation}
 \lambda^A\ll X\le Y^{9/10},\qquad
 Y/2\le P\ll Y^2\lambda^B/X .
 \label{eq:productscales}
\end{equation}
In this range the ordinary interpolation and unit-condition
errors from the previous subsections are
$O(\lambda^2 K^{-1/2}+\lambda^{4+B/2}/X)$ for
$K=\lceil\lambda^k\rceil$. The interpolation proof uses
$r\ge2$ and $\log P=O(\lambda)$, which hold here. The
model has size $O(\lambda\sqrt{X/P})$ by \cref{eq:Mpacket}.
Use the signed tent argument with $Q=\lfloor\lambda^q\rfloor$.
If $A>2q$, then $Q\le\sqrt X/4$ and the arc bounds apply.
Also $P/K\gg X$, so the same argument for the main terms
gives $O(K^2 X^{3/2}\lambda^2/(Y\sqrt P))$. Both this term
and the model have a power saving in $Y$ on
\cref{eq:productscales}.

To keep logarithmic losses in the modulus sum, let $\tau(n)$
count divisors and choose a fixed exponent $m>0$. There are
$O(Y\lambda^{1-m})$ moduli with $\tau(N)>\lambda^m$, by
$\sum_{n\le U}\tau(n)\le U(1+\log U)$. On those moduli
we use the absolute packet bound
$O(\lambda^2\sqrt{XP}/Y)$, which follows by counting $O(X)$
integers of weight $O(\lambda^2)$ on each of $O(P/Y)$
lines, using the physical supremum bounds and
\cref{eq:subseries} for the continuous lines.
For the other moduli discard the shifts $h=bN$ with
$\tau(b)>\lambda^m$ from the interpolated arithmetic
sum before subtracting the major-arc main term. The same
divisor estimate gives $O((P/Y)\lambda^{1-m})$ such
shifts for each $N$. The physical supremum bounds for the
two interpolants then give a cost of
$O(\lambda^{3-m}\sqrt{XP}/Y)$. Each remaining shift $h$
can occur for at most $\tau(h)\le\lambda^{2m}$ moduli,
since $\tau(bN)\le\tau(b)\tau(N)$.

For the retained shifts, apply the Frobenius bound and
\cref{eq:mixedarcerror} with the same shift-count factor
$O((P/Y)/(W_xW_y))$ as in the fixed-power case. The mean
square of the residual is bounded by $Y\lambda^{C_1+2m+B}$
times
\[
  \lambda^{3k-q}+\lambda^{3k-A}+\lambda^{4k+3q-2A},
\]
where $C_1$ is absolute. Here $XP\ll Y^2\lambda^B$ was
used in each term. Put $D_1=D+8J$. To obtain
$O(Y\lambda^{-D_1-6})$ per
packet pair, first choose $k>D_1+12$ and $m>B+D_1+12$
for the interpolation and divisor errors. Then choose $q$
large enough to control the first term,
and finally choose $A>2q$ large enough to control the
remaining terms and the unit-condition error. These choices
are independent of the packet indices. The Fourier multiplier
loss is $O(\lambda^{4J})$ before squaring. Applying the
triangle inequality in the modulus mean-square norm over the
$O(\lambda^2)$ pairs and including the sharp boundary error
now proves \cref{eq:productwindow}.

\subsection{Completing the main terms in the power range}

We bound the ordinary form for a log packet pair at scales
$X\le V$ with $P=eV\ge Y/2$. We assume that $X$ grows as
at least a fixed positive power of $Y$ and that $\log P\ll\log Y$.
Use $K=\lceil Y^\omega\rceil$ cells per packet and
$Q=\lfloor\sqrt X/8\rfloor$. Linear interpolation and
removing the unit restriction cost
$O((\log Y)^2K^{-1/2}+(\log Y)^4Y^{-1}\sqrt{P/X})$ on
the normalized scale, as in the preceding proofs.
Write the interpolants in the tent bases and apply
\cref{eq:mixedarcerror} to each of the three orientations.
Here we compare each $M_{ij}(h)$ with its full continuous
term, including the singular series.

Write $J_{ij}(h)$ for the full real integral of the
product of the two continuous Fourier transforms against
$e(-h\beta)$, with the signs for the orientation in question.
By Fourier inversion it is the integral of the two tents
along that line, so $|J_{ij}(h)|\ll W_x$. For each $q\le Q$
the integral over $|\beta|>2Q/(qX)$ has absolute value
$O(W_x Xq/(W_y Q))$. This follows from the bound $O(W_x)$
on the transform of the smaller tent and the quadratic
decay $O(W_y(1+W_y|\beta|)^{-2})$ of the larger one.

The coefficient of this integral is
$A_q(h)=\mu(q)^2c_q(h)/\varphi(q)^2$ with the Ramanujan
sum $c_q(h)$. For $h\ge1$ and every $\epsilon>0$ we have
\begin{equation}
 \sum_{q>Q}|A_q(h)|
  +\frac1Q\sum_{q\le Q}q|A_q(h)|
   \ll_\epsilon \frac{P^\epsilon}{Q}
       \qquad (h\le 2P).
 \label{eq:ramanujancompletion}
\end{equation}
Indeed, for squarefree $q$, $|c_q(h)|\le\gcd(q,h)$.
Use $q/\varphi(q)\ll_\nu q^\nu$ and
$\gcd(q,h)\le\sum_{d\mid(q,h)}d$. In the first sum, for
each $d\mid h$ the sum over multiples $q=d\ell>Q$ is
$O_\nu(Q^{-1+2\nu})$ for $0<\nu<1/2$, by treating
$d\le Q$ and $d>Q$ using the convergent power sum.
In the second sum the inner sum over $\ell$ is bounded
by $O(Q^{2\nu}(1+\log Q))$ before division by $Q$.
Now use $\tau(h)\ll_\nu P^\nu$ and choose $\nu$ sufficiently
small in terms of $\epsilon$.

For an even shift $h$, the full sum of $A_q(h)$ is $\SSS(h)$
by the Euler product. Thus the difference between the
truncated and completed main terms is
\[
 |M_{ij}(h)-\SSS(h)J_{ij}(h)|
     \ll_\epsilon W_x(1+X/W_y)P^\epsilon/Q .
\]
On using $X\le P$ and $W_x=X/K$, summing this per-line
bound over the $O(P/Y)$ shifts and applying the nodal
$\ell^1$ bounds gives a normalized completion error of
\begin{equation}
 O_\epsilon\!\left(\frac{K^2\sqrt{XP}}{YQ}P^\epsilon\right).
 \label{eq:powercompletionerror}
\end{equation}
The completed terms sum to the continuous model with the
original signs. For each tent pair it suffices to sum over
shifts for which the physical line intersects the two
supports, since the arithmetic and the completed model
vanish at the other shifts.

For the remaining error use the Frobenius estimate,
Cauchy--Schwarz over the $O(P/Y)$ shifts in each entry,
and $\tau(h)\ll_\epsilon P^\epsilon$ when summing over
the moduli with $N\mid h$. With \cref{eq:mixedarcerror}
this costs at most
\begin{equation}
 O_\epsilon\left(\frac{P^{1+\epsilon}K^4\sqrt X}{Y}\right)
 \label{eq:completedpacket}
\end{equation}
in mean square, using $Q\asymp\sqrt X$ and resizing
$\epsilon$ to absorb the logarithms and divisor bound.
If $P\sqrt X\ll Y^{2-\sigma}$ for fixed $\sigma>0$, then
\cref{eq:completedpacket} is $O(Y^{1-c})$ for some $c>0$
by taking $\omega,\epsilon$ small. Since $P\ll Y^{2-\sigma}$,
\cref{eq:powercompletionerror} and the deterministic
interpolation and unit-condition errors also have power
savings with the same choices.

To prove \eqref{eq:generalrectpowers}, smoothly remove
the parts of each test below $N^{2-b-\rho}$, for a small
fixed $0<\rho<2-b$. Any pair discarded this way has
$xy\ll N^{2-\rho}$, so \cref{thm:subcritical} gives a
power-saving error. On the remaining packet pairs,
$X\gg Y^{2-b-\rho}$, $X\ll Y^a$, and $P\ll Y^b$.
Pairs with $2P<Y$ have no term. Thus
$P\sqrt X\ll Y^{b+a/2}$, and the strict inequality
$b+a/2<2$ gives the desired saving. Summing over
$O((\log Y)^2)$ packet pairs costs only logarithms.
The symmetric case $a=b=1+\eta$ completes the proof of
\cref{eq:almostpower} for $\eta<1/3$.

For the smooth cutoff in \cref{thm:powerproduct}, use
Fourier separation on unit log packets as in the preceding
subsection. The fixed cutoff costs $O_\Phi(1)$ per pair.
In the remaining $O((\log Y)^2)$ pairs we have
\begin{equation}
 X\gg Y^\delta,\qquad P\ge Y/2,\qquad XP\ll Y^{2+\gamma}.
 \label{eq:powerproductscales}
\end{equation}
Hence $P\sqrt X\ll Y^{2+\gamma-\delta/2}$, which gives a
power saving in the packet bound if $\gamma<\delta/2$.
This proves \cref{eq:powerproduct} for the smooth cutoff.

\subsection{Maximal interval counts for the sharp power cutoff}

For the sharp cutoff use the same packet range
\cref{eq:powerproductscales}, now without Fourier separation.
On a line $y=\pm x+c$, the condition $xy\le N^{2+\gamma}$
cuts the smaller coordinate into at most two intervals.
On each such interval the test product
$w(x)\overline{z(\pm x+c)}$ has supremum and total
variation $O((XP)^{-1/2})$ on unit tests. This follows
from the supremum and $L^1$ derivative bounds for the
physical weights on unit log packets. Partial summation
therefore reduces the weighted error for each orientation
to $(XP)^{-1/2}$ times the maximum unweighted discrepancy
over all subintervals of the smaller packet, with the
larger interval $J=[V/e,eV]$ fixed. Removal of the unit
condition again costs $O((\log Y)^4Y^{-1}\sqrt{P/X})$
after summing over the shifts.

Take $M$ a power of two in
$[Y^{\gamma/2+\rho},2Y^{\gamma/2+\rho}]$ for a fixed $\rho>0$,
and use a grid with cell width $X/M$, translated by
$u\in[0,X)$. Replacing an arbitrary
subinterval by the union of its full cells leaves at most
two partial cells. On a shift $h$, these have absolute
arithmetic weight $O((X/M+1)(\log(2P))^2)$ and absolute
model weight $O(\SSS(h)X/M)$. Apply the point-evaluation
normalization and sum over the $O(P/Y)$ shifts $h=bN$.
Using \cref{eq:subseries}, the total rounding error is
at most
\begin{equation}
 O\!\left((\log Y)^2\frac{\sqrt{XP}}Y
              \left(\frac1M+\frac1X\right)\right).
 \label{eq:sharprounding}
\end{equation}

Use the nested dyadic grids based at the same translation $u$,
with block lengths $\ell=2^j X/M$, $0\le j\le\log_2 M$.
Let $\mathcal I_\ell(u)$ be the full blocks of length $\ell$
lying in the smaller packet. Every consecutive union of
fine cells in the packet is the disjoint union of
$O(1+\log M)$ full blocks, using at most two at each
smaller level and $O(1)$ at the largest level.
Let $\mathbb E_u$ denote the average over $0\le u<X$.
For any nonnegative function of the left endpoint, its
expected sum over $\mathcal I_\ell(u)$ is at most
$\ell^{-1}$ times its integral over the packet, since
$X$ is a multiple of $\ell$.

We use the two mean-square bounds for the translated blocks
\begin{equation}
 \begin{split}
 \mathbb E_u\sum_{I\in\mathcal I_\ell(u)}
    \sum_{a\bmod q}^{*}|E_{I,a}(\beta)|^2
       &\ll qX(1+\ell|\beta|)^2(\log(2P))^4,\\
 \mathbb E_u\sum_{I\in\mathcal I_\ell(u)}
    \int_{-B}^B\sum_{a\bmod q}^{*}|E_{I,a}(\beta)|^2\,\ud\beta
       &\ll qXB(\log(2P))^4 .
 \end{split}
 \label{eq:translatedblocks}
\end{equation}
Here $q\le X/(4e)$ and $0<B\le1/2$. To see this, use
the signed measure
$\mu_\chi=\sum_n\Lambda(n)\chi(n)\delta_n-\delta_\chi\,\ud t$
on the packet. The short-interval estimate in the proof
of \cref{eq:arcintegrated} is
\[
 \int_{s\asymp X}|\mu_\chi((s,s+h])|^2\,\ud s
      \ll Xh(\log(2P))^4,\qquad 0<h\le X,
\]
where the measure can first be defined on the full positive
axis. For $h<1$ this follows by integrating the weight of
the at most one integer in each interval and adding the
principal term. Partial summation on $(s,s+\ell]$ gives
an $L^2(\ud s)$ bound
$O(\sqrt{X\ell}(1+\ell|\beta|)(\log(2P))^2)$
for its modulated mass. Indeed, the endpoint term uses
$h=\ell$, and the integral term is bounded by
$O(|\beta|(\log(2P))^2\int_0^\ell \sqrt{Xh}\,\ud h)$.
Divide the squared bound by $\ell$ and use multiplicative
and additive Plancherel as in \cref{eq:arcmeansquare}.

For the integrated bound put $H=1/(2B)$. Convolve the
restriction of $\mu_\chi$ to $(s,s+\ell]$ by an interval
of length $H$, as in the proof of \cref{eq:arcintegrated}.
After putting $t=s+v$ for the left endpoint of the second
interval, its intersection with $(s,s+\ell]$ has length
$h(v)$ depending only on $v$ and left endpoint
$s+\max(0,v)$. As $v$ ranges from $-H$ to $\ell$,
$\int h(v)\,\ud v=\ell H$. The short-interval estimate,
integrated over $v$ and divided by $\ell$, therefore
bounds the expected sum of squared convolution norms by
$O(XH(\log(2P))^4)$. Plancherel and the lower bound
$\gg H$ for the interval transform on $[-B,B]$ give
$O(XB(\log(2P))^4)$ before the Plancherel factor $q$.
The total weight of the nonunit prime powers in the
packet is $O((\log(2P))^2)$. Their squared weights summed
over disjoint blocks are bounded by the square of that
total, so they contribute within both bounds.

For one orientation write $D_I(h)$ and $M_I(h)$ for the
arithmetic and major-arc coefficients of the intervals
$I,J$. The estimates \eqref{eq:translatedblocks} give,
at each level and for $Q=\lfloor\sqrt X/8\rfloor$,
\begin{equation}
 \mathbb E_u \sum_{I\in\mathcal I_\ell(u)}\sum_{1\le h\le2P}
       |D_I(h)-M_I(h)|^2
      \ll P X^{3/2}(\log(2P))^{C_2}.
 \label{eq:dyadicintervalarc}
\end{equation}
For the minor arcs use the rational approximant with
$Q<q\le X/Q$ and $|\beta|\le 2Q/(qX)$. The first bound
in \eqref{eq:translatedblocks} and $\varphi(q)\gg\sqrt q$
give $\mathbb E_u\sum_I |S_I(\alpha)|^2\ll
(X^2/Q)(\log(2P))^4$. Use Parseval on $J$.
For one error on the major arcs, combine the integrated
bounds with the decay of the continuous factors, using
$\sum_I |v_I(\beta)|^2\ll X\ell(1+\ell|\beta|)^{-2}$
when the error is on $J$. Summation over the denominators
costs $O(XP(\log(2P))^{C_2})$. For two errors, use the
first bound of \eqref{eq:translatedblocks} pointwise on
the smaller variable and the integrated bound on $J$.
This costs $O(PQ^3(\log(2P))^{C_2})$ as in
\cref{eq:mixedarcerror}, since $\ell\le X$.
Bessel gives \eqref{eq:dyadicintervalarc} because
$X^2/Q$ and $Q^3$ are both $O(X^{3/2})$.

Write $J_I(h)$ for the continuous line length. Its
absolute size is $O(X)$, and the tails of the product
of its two Fourier factors satisfy
$\int_{|\beta|>2Q/(qX)}|v_I(\pm\beta)v_J(\pm\beta)|\,\ud\beta
\ll Xq/Q$, by the interval estimate $O(1/|\beta|)$ for
each factor. Thus \cref{eq:ramanujancompletion} gives
$|M_I(h)-\SSS(h)J_I(h)|\ll_\epsilon XP^\epsilon/Q$
for positive even $h$. Summing this uniform error over
the dyadic decomposition of any consecutive set of cells
and then over the shifts gives a normalized cost
$O_\epsilon((1+\log M)\sqrt{XP}P^\epsilon/(YQ))$.

For each translation, Cauchy--Schwarz bounds the square of
the maximum residual by $1+\log M$ times the sum of
$|D_I(h)-M_I(h)|^2$ over the levels and their blocks.
The packet norm being bounded is independent of the grid,
so we can average this upper bound in $u$. A second
Cauchy--Schwarz over $O(P/Y)$ shifts per modulus, followed
by the divisor bound on the number of moduli with $N\mid h$,
gives a mean square bounded by
\[
 O_\epsilon\left(\frac{P^{1+\epsilon}\sqrt X}{Y}\right),
\]
using \cref{eq:dyadicintervalarc}, absorbing the logarithms
from the levels, and applying the normalization factor
$(XP)^{-1}$ for squared packet norms. Since $\gamma<\delta/2$,
the last bound is $O(Y^{1-c})$ for sufficiently small
$\epsilon>0$ by \cref{eq:powerproductscales}. The rounding error
\eqref{eq:sharprounding}, completion error and unit-condition
error also have uniform power savings in that range.
Summing over the three orientations and the
$O((\log Y)^2)$ packet pairs proves the sharp-cutoff
assertion of \cref{thm:powerproduct}.

\section{Removing proper prime powers at a fixed modulus}
\label{sec:ppcut}

Call $p^k$ a proper prime power if $p$ is prime and $k\ge2$.
Let $E_N=\|T_N\|/\SSS(N)$ under $\grh(N)$. For $A\ge2$ let
$E_N^{[A]}$ be the normalized norm of \cref{eq:form} with every
proper prime power $n\ge A$ removed from its arithmetic terms.
The continuous model is the same in this definition.

\begin{theorem}[Proper prime powers and the small-input terms]
\label{thm:ppcut}
Assume $\grh(N)$ and put $L=\log N$. Uniformly for $2\le A\le N/8$,
write $\delta_A=A^{-1/4}\sqrt{\log(2A)}+N^{-1/5}$. Then
\begin{equation}
 \max(E_N,E_N^{[A]})
 \le (1+O(\delta_A/\sqrt L))\min(E_N,E_N^{[A]})
        +O(\sqrt L\,\delta_A).
 \label{eq:ppnorm}
\end{equation}
In particular, if $A=L^2(\log L)^2 K$, where
$1\le K=K(N)\le L$ and $K\to\infty$, the
additive term is $O(K^{-1/4})$. This leaves
$O(L\sqrt K)$ proper prime powers in the arithmetic form.

For the same choice of $A$, \cref{eq:goal} is equivalent to
$P_N(A)\to0$ and $J_N(A)\to0$. Here $P_N(A)$ is the normalized
norm on $x,y\ge A$ with the arithmetic sum restricted to
ordinary primes. The quantity $J_N(A)$ is the normalized norm,
for tests $F$ supported in $x\le2A$ and arbitrary $G$, of
\begin{equation}
 \sum_{\substack{2\le n\le2A\\(n,N)=1}}
   \frac{\Lambda(n)}{\sqrt n}F(\log n)
    \overline{\mathcal C_{N,n}(G)},\qquad
 \mathcal C_{N,n}(G)=
 \sum_{\substack{p\ {\rm prime}\\p\ge N-2A}}
     \frac{\log p}{\sqrt p}G(\log p)
     \bigl(\mathbf1_{p\equiv n\ (N)}-\mathbf1_{p\equiv -n\ (N)}\bigr).
 \label{eq:pprows}
\end{equation}
The latter form has rank $O(A/\log(2A))$.
\end{theorem}

\subsection{A bound on the proper-prime-power character vector}

For a set $\mathcal P$ of proper prime powers define
$U_{\mathcal P}$ by restricting $\ell_\chi$ in \cref{sec:cutoff}
to $n\in\mathcal P$. All bounds here are on compactly supported
tests. For $A\le n<N/3$, the positive integers under consideration
are distinct modulo $N$ and no two have sum divisible by $N$.
Thus odd-character orthogonality gives
\[
 \|U_{\mathcal P}(F)\|_\chi^2
     =\sum_{\substack{n\in\mathcal P\\(n,N)=1}}
       \frac{\Lambda(n)^2}{n}|F(\log n)|^2
     \ll \frac{\log(2A)}{\sqrt A}\NN(F)^2
     \qquad(\mathcal P\subset [A,N/3)).
\]
For the bound, Sobolev point evaluation on $r\ge0$ gives
$\sup |F(r)|\ll\NN(F)$. On $[X,2X)$, $X\ge2$, squares
contribute $O(\log(2X)/\sqrt X)$ to the unweighted diagonal
sum $\sum_{n\in\mathcal P}\Lambda(n)^2/n$. Chebyshev's bound
$\sum_{p\le t}\log p\ll t$ gives this estimate, and it gives
$O((\log(2X))^2 X^{-2/3})$ for all $k\ge3$ together.
Summing over dyadic intervals yields the displayed bound.
The same calculation bounds the change in $\mathcal D$
upon deleting all proper prime powers $n\ge A$ by
$O(\log(2A)/\sqrt A)$ in form norm.

For a subset $\mathcal P\subset[N/3,\infty)$, localize to the
unit log packets $F_j$ from \cref{sec:cutoff}, with $X=e^j$
and $j\ge L-O(1)$. Fix $k=2$ or $3$. A unit congruence
$b^k=u\bmod N$ has $O_k(k^{\omega(N)})$ solutions modulo $N$,
where $\omega(N)$ counts distinct prime factors.
The bound follows by counting the kernel of the power map at
each prime power and applying the Chinese remainder theorem,
as in \cref{sec:subcritical}. It is $O_{k,\varepsilon}(N^\varepsilon)$
for every $\varepsilon>0$. Each unit class therefore contains at
most $O_\varepsilon(N^\varepsilon(1+X^{1/k}/N))$ values $p^k$
in the packet. Bounding the odd-character norm by twice the
all-character mean square and applying Cauchy--Schwarz in
each residue class gives
\[
 \|U_{\mathcal P\cap\{p^k\}}(F_j)\|_\chi^2
    \ll_\varepsilon N^\varepsilon (1+j)
        (1+X^{1/k}/N)X^{1/k-1}e^{-2j/L}\NN(F_j)^2 .
\]
Here the squared coefficient sum on the packet uses Chebyshev's
bound, and $\|F_j\|_\infty\ll e^{-j/L}\NN(F_j)$.
For $k=2,3$ we have
$(1+X^{1/k}/N)X^{1/k-1}\le X^{-1/2}+N^{-1}$.
Summing the vector norms using
Cauchy--Schwarz in $j$ gives $O_\varepsilon(N^{-1/4+\varepsilon}
\NN(F))$. Indeed, the sum of
$(1+j)e^{-2j/L}(e^{-j/2}+N^{-1})$ over $j\ge L-O(1)$ is
$O(L/\sqrt N+L^2/N)$.
For all $k\ge4$ the sum of $\Lambda(p^k)/\sqrt{p^k}$ on
one packet is $O((1+j)X^{-1/4})$ by Chebyshev's bound.
The triangle inequality followed by summation in $j$ gives the
same $O_\varepsilon(N^{-1/4+\varepsilon}\NN(F))$ estimate.
In particular,
\begin{equation}
 \|U_{\mathcal P}(F)\|_\chi \ll N^{-1/5}\NN(F)
     \quad(\mathcal P\subset[N/3,\infty)),\qquad
 \|U_{\mathcal P}(F)\|_\chi \ll \delta_A\NN(F)
     \quad(\mathcal P\subset[A,\infty)).
 \label{eq:ppchar}
\end{equation}
These vector and diagonal estimates use only residue counting.

\subsection{Norm comparison and the progression profiles}

Let $E$ be the normalized norm of either form in
\cref{eq:ppnorm}. For a unit test, the corresponding character
vector has norm $O(\sqrt{\SSS(N)(L+E)})$ by positivity and
\cref{eq:DMtotal}. For the form with powers removed, use its
reduced diagonal, which is bounded by the original one.
Finiteness of both norms follows from the bounded original form
and the bound on the deleted vector in \cref{eq:ppchar}.
Pass to the other form by adding or subtracting the deleted
vector in each argument. By \cref{eq:ppchar} and $\SSS(N)\gg1$,
the two mixed terms cost $O(\delta_A\sqrt{L+E})$ on the
normalized scale. The term with two deleted vectors and the
change in the diagonal cost $O(\delta_A^2)$. Take the smaller
of $E_N,E_N^{[A]}$ and use
$\sqrt{L+E}\le\sqrt L+E/(2\sqrt L)$
to get \cref{eq:ppnorm}. The same argument compares these forms
on a common support restriction. For the stated choice of $A$,
$\log A\asymp\log L$ and
$\sqrt L\,\delta_A\ll K^{-1/4}$.
Chebyshev's prime counting bound gives $O(\sqrt A/\log A)$
proper prime powers below $A$, by treating squares separately
and bounding the rest by $O(A^{1/3}\log A)$.
An absolute norm comparison is also available. Split at $r=2$
with a smooth cutoff, use \cref{eq:Utail} on the positive
part, and count the finitely many prime powers on the lower
part to obtain $\|U(F)\|\ll L\NN(F)$ on the full space.
The same decomposition into deleted and retained vectors
then gives $|E_N-E_N^{[A]}|\ll L\delta_A$.

Let $H_N(A)$ be the normalized original form norm with both
supports in $x,y\ge A$, and let $B_N(2A)$ be the normalized
original form norm with the first support in $x\le2A$.
The support version of \cref{eq:ppnorm} compares $H_N(A)$
to $P_N(A)$. For the small-support restriction, any
off-diagonal pair has $y\ge N-2A$ since the nonzero lines
have shift of absolute value at least $N$. In particular the
two input supports cannot interact within $x,y\le2A$.
The continuous model on the small-support restriction costs
$O(A/N)$ by the bound for $\min(x,y)\le2A$ in
\cref{sec:subcritical}. For the arithmetic sum, use
\cref{eq:ppchar} to delete the proper prime powers $y\ge N-2A$.
The small-input character norm is $O(\sqrt{\log(2A)}\,\NN(F))$,
since odd-character orthogonality leaves only its arithmetic
diagonal, whose bound follows from the diagonal estimate in
\cref{sec:cutoff} on log packets with $e^j\ll A$.
Thus
\begin{equation}
 |B_N(2A)-J_N(A)|\ll A/N+N^{-1/5}\sqrt{\log(2A)} .
 \label{eq:ppboundrows}
\end{equation}
The resulting arithmetic form is exactly \cref{eq:pprows}.
The rank bound counts the primes and proper prime powers
below $2A$.

Finally, a smooth log cutoff from $A$ to $2A$ splits each
test into a part supported in $x\le2A$ and a part supported in
$x\ge A$. Both multipliers are bounded in $\NN$, uniformly
in $A$ and $N$. Expand the pairing as the term with both
high parts plus a term with the first low part and a term
with the second low part. Self-adjointness bounds the latter
two terms by the small-support norm. Thus
\[
  \max(H_N(A),B_N(2A))\le E_N\ll H_N(A)+B_N(2A).
\]
Applying \cref{eq:ppnorm,eq:ppboundrows} proves the
equivalent criterion.

\subsection{Truncation of the progression profiles}

For $Q\ge e$ and $\kappa\ge1$, put
\[
 Y_-=\frac{N^2}{2A Q^2},\qquad
 Y_+=N^{\,\log L+\frac12\log\log(2A)+\kappa}.
\]
Let $J_N^{\rm win}$ be $J_N(A)$ with the second
support restricted to $[Y_-,Y_+]$. For the same range of $A$
as in \cref{thm:ppcut}, we have
\begin{align}
 J_N^{\rm win}&\le J_N(A), \notag\\
 J_N(A)&\le
  (1+O(1/\log Q+1/L))J_N^{\rm win}
       +O\left(Q^{-1/2}+e^{-\kappa}
                  +N^{-1/5}\sqrt{\log(2A)}\right).
 \label{eq:ppwindow}
\end{align}
Indeed, split $G$ with a lower log transition on $[Y_-,QY_-]$
in physical coordinates and an upper log transition from
$H=\log Y_+-L$ to $\log Y_+$ in log coordinates.
The part in the retained interval has $\NN$ norm at most
$1+O(1/\log Q+1/L)$ times the original norm by the
multiplier estimate in \cref{sec:cutoff}. The discarded
parts have bounded multiplier norms.

For the lower discarded part, $xy\le (2A)(QY_-)=N^2/Q$.
The absolute arithmetic estimate in \cref{sec:subcritical}
applies to the prime-only second coordinate, giving a cost
$O(Q^{-1/2}+N^{-2/5})$ on the normalized scale.
For the upper discarded part, $y\ge e^H\ge N$ for all
sufficiently large $N$, uniformly in the stated parameters.
Its character vector with primes only has norm
$O((Le^{-H/L}+N^{-1/5})\NN(G))$ by
\cref{eq:Utail,eq:ppchar}. Pairing with the small-input vector
costs $O(\sqrt{\log(2A)}(Le^{-H/L}+N^{-1/5}))$
on the normalized scale, which yields \cref{eq:ppwindow}.
For instance, $A=L^2(\log L)^4$, $Q=(\log L)^2$ and
$\kappa=\log\log L$ give a vanishing error. The resulting
window runs from $N^{2-o(1)}$ to
$N^{\log\log N+O(\log\log\log N)}$.

\section{Norm decay of the progression profiles near the square of the modulus}
\label{sec:criticalprofiles}

Let $J_N(A;R)$ be the norm in \cref{eq:pprows} with the
second test supported in $y\le R$. We use the same
normalization by $\SSS(N)$. The restricted form can be
defined without GRH. Throughout this section $R_N=N^2(\log N)^B$
with $B\ge0$ fixed.

\begin{theorem}[Small-input terms through the critical scale]
\label{thm:criticalprofiles}
For every fixed $A\ge2$ and $D>0$,
\begin{equation}
 \#\{N\in[Y,2Y]: N\ {\rm even},\
                   J_N(A;R_N)>(\log Y)^{-D}\}=o(Y).
 \label{eq:fixedprofiledecay}
\end{equation}
For growing $A$, let $\mathcal F(Y)$ be the even moduli in
$[Y,2Y]$ with a prime factor in each of the intervals
\[
 (Y^{1/25},Y^{1/25+\nu}],\qquad
 (Y^{17/100},Y^{17/100+\nu}],\qquad \nu=1/2000.
\]
Then $\#\mathcal F(Y)\asymp Y$. For each fixed $C,D>0$,
uniformly for $2\le A_N\le(\log N)^C$,
\begin{equation}
 \sum_{N\in\mathcal F(Y)} J_N(A_N;R_N)^2
       \ll_{B,C,D}Y(\log Y)^{-D}.
 \label{eq:factoredprofiledecay}
\end{equation}
Both estimates also hold for the original normalized arithmetic
and model form with supports in $x\le2A_N$ and $y\le R_N$,
with $A_N=A$ in \cref{eq:fixedprofiledecay}.
\end{theorem}

\subsection{From progression counts to the profile norm}

Write $\lambda=\log Y$, let $A_*=A$ or $A_*=(\log(2Y))^C$
respectively, and take a sufficiently large fixed $M$. Set
$T_-=Y^2\lambda^{-M}$ and $T_+=(2Y)^2(\log(2Y))^B$.
A smooth log cutoff in $[T_-,2T_-]$ reduces the proof to
tests on $T_-\le y\le T_+$. The discarded part has
$xy\le4A_*T_-$. The unsigned arithmetic estimate in
\cref{sec:subcritical} therefore costs
$O((A_*\lambda^{-M})^{1/2}+Y^{-2/5})$ on the
normalized scale, by taking $Q=Y^2/(4A_*T_-)$.
This bound applies with the second coordinate restricted
to primes. The cutoffs have bounded $\NN$ multiplier norms.

For $(a,N)=1$, put
\[
 D_{N,a}(t)=\pi(t;N,a)-\frac{\pi(t)}{\varphi(N)},
 \qquad
 d_N(t)=\max_{\substack{1\le |a|\le2A_*\\(a,N)=1}}
                  \frac{|D_{N,a}(t)|}{\sqrt t}.
\]
On the retained support the lower constraint $p\ge N-2A_N$
is redundant. The common main term cancels between the
two residues. Partial summation with
$w(r)=r e^{-r/2}G(r)$ and Cauchy--Schwarz therefore give
\[
 |\mathcal C_{N,n}(G)|
   \ll \lambda\NN(G)
          \left(\int_{\log T_-}^{\log T_+} d_N(e^r)^2\,
                              \ud r\right)^{1/2}.
\]
Indeed, $|w'(r)|\ll\lambda e^{-r/2}(|G(r)|+|G'(r)|)$
on this interval, and the unweighted $H^1$ norm there is
bounded by $\NN(G)$. Boundary terms vanish for supported
smooth tests. Summing the first-argument weights in
\cref{eq:pprows} costs $O(\sqrt{A_*}\,\NN(F))$ by Sobolev
point evaluation and Chebyshev's bound.

Suppose a family $\mathcal G(Y)\subset[Y,2Y]$ satisfies
\begin{equation}
 \sum_{\substack{N\in\mathcal G(Y)\\(a,N)=1}}
       |D_{N,a}(t)|\ll_U t\lambda^{-U}
 \label{eq:profiledistribution}
\end{equation}
for every fixed $U>0$, uniformly in $T_-\le t\le T_+$
and $1\le |a|\le 2A_*$. Constants may depend on the fixed
parameters of the family. The trivial bound
$|D_{N,a}(t)|\ll \lambda t/Y$ gives
\[
 \sum_{N\in\mathcal G(Y)}d_N(t)^2
       \ll_U A_* \lambda^{1-U}\frac{t}{Y}
       \ll_{B,U}Y A_*\lambda^{1+B-U}.
\]
Here we bound the squared maximum by a sum over the
admissible residues. Integrating over a log interval of
length $O_{B,M}(\log\lambda)$ and using $\SSS(N)\gg1$ proves
the mean-square bound $O(Y\lambda^{-D})$ on $\mathcal G(Y)$.
Choose $M$ large enough in terms of $C,D$ for the lower
cutoff and then $U$ large enough to absorb the remaining
powers of $\lambda$.

\subsection{Families with the required progression estimates}

For fixed $A$, take a small absolute $c>0$ for which the
sieve bound \cref{eq:uniformsieve} holds. Take
$0<\delta<c/8$ sufficiently small, and let $\mathcal G_\delta(Y)$
be the even moduli in $[Y,2Y]$ divisible by some prime
$p\in(Y^{7\delta},Y^c]$. Maynard's fixed-residue theorem
\cite[Corollary~1.2]{MaynardI} gives \cref{eq:profiledistribution}
on this family. In fact $2Y<t^{1/2+\delta}$, and $p$ is a
divisor in the required interval
\[
 t^{3\delta}<p<
       \min\bigl(t^{1/10-12\delta/5},\,t^{1/2-20\delta}\bigr)
\]
when the second parameter of that corollary is $\delta$.
These inequalities hold for all $T_-\le t\le T_+$, since
$\log t/\log Y=2+o(1)$. The residue range is fixed, so
the finitely many constants in Maynard's bound can be
taken uniformly.

To count the complementary moduli, use
\cref{eq:uniformsieve} with one forbidden residue $0$
at each prime in $(Y^{7\delta},Y^c]$. Mertens' product
formula shows that the number is $O_c(\delta Y)+o_\delta(Y)$.
On $\mathcal G_\delta(Y)$ the mean-square bound with a
large enough exponent makes the count exceeding
$\lambda^{-D}$ equal to $o_\delta(Y)$. Taking the upper
limit as $Y\to\infty$ and then letting $\delta$ decrease to
zero proves \cref{eq:fixedprofiledecay}.

For \cref{eq:factoredprofiledecay}, use Maynard's estimate
with three factors \cite[Theorem~1.2]{MaynardIII}. Take
$\delta=1/10000$. Group the possible two prime factors into
dyadic ranges with upper endpoints $Q_2,Q_3$. Their exponents
relative to $t$ lie, up to $o(1)$, in
$[0.020,0.02025]$ and $[0.085,0.08525]$ respectively.
Consequently the conditions
\[
 t^{40\delta}<Q_2<t^{1/20-7\delta},\qquad
 t^{1/10+12\delta}/Q_2<Q_3<t^{1/10-4\delta}/Q_2^{3/5}
\]
hold with room in each inequality. Set
$Q_1=t^{1/2+\delta}/(Q_2Q_3)$. A factorization
$N=q_1p_2p_3$ in these ranges has $q_1\le Q_1$, since
$8Y<t^{1/2+\delta}$. For each fixed residue $a$ coprime
to $N$, choose the class $b=a\bmod q_1p_2$ in Maynard's
outer supremum. The inner supremum includes $a$ for
every $p_3$ in the bin. This bounds the sum in
\cref{eq:profiledistribution} uniformly in $a$ by the sum
of Maynard's bounds over $O(\lambda^2)$ bins. The power of
$\lambda$ lost in the bin count is absorbed by choosing
the saving in Maynard's bound large enough.

Finally, each pair of primes in the specified ranges
divides $Y/(2p_2p_3)+O(1)$ even integers in $[Y,2Y]$.
Summing over pairs gives a positive constant times $Y$,
up to $o(Y)$, by the prime reciprocal sum in each interval
and the bound $O(Y^{21/100+2\nu})$ on the number of pairs.
Each even modulus is counted only $O(1)$ times since
every such prime factor is at least $Y^{1/25}$. This proves
$\#\mathcal F(Y)\asymp Y$. The comparison
\cref{eq:ppboundrows} applies with the same second support
restriction, and its error is a power saving for
$A_N\le(\log N)^C$. This proves the assertions about the
original form.

\subsection{The same-label correction near the square of the modulus}

For each fixed $B\ge0$ and $D>0$, all but $o(Y)$ even
moduli in $[Y,2Y]$ satisfy
\begin{equation}
 |\mathcal P_N(F,G)|\le(\log Y)^{-D}\NN(F)\NN(G)
 \label{eq:samelabelmean}
\end{equation}
uniformly for tests on the high supports of \cref{thm:zeros}
whose log-coordinate differences are at most $\log R_N$ in
absolute value. Thus \cref{eq:samelabelcorrection} reduces
the form to distinct zero labels to this precision,
with a change of the implied constant, at these moduli
when $\grh(N)$ holds.

To prove the bound, use $h(v)=S_{F,G}(v)$ as a log test.
Its supremum and the supremum of its derivative are
$O(\NN(F)\NN(G))$ by Cauchy--Schwarz in the correlation.
It is supported in $|v|\le\log R_N$, so
$\NN(h)\ll_B\sqrt L\,\NN(F)\NN(G)$ using
$W_N(v)\le e^{2v/L}$. The part of $\mathcal P_N$ on proper
prime powers is the mean over odd characters of the
corresponding character vector restricted to $n\ge N-1$.
By \cref{eq:ppchar} it is bounded by
$O_B(N^{-1/5}\sqrt L\,\NN(F)\NN(G))$.

Split the prime sum at $T_-$ by the smooth cutoff used
above. The lower part costs $O(\lambda^{1-M/2}\|h\|_\infty)$
by summing over $bN\pm1\le2T_-$ as in
\cref{eq:samelabelcorrection}. On the upper part use the
partial summation estimate with residues $\pm1$ and the
bound on $\NN(h)$. The estimate for the integral of
$d_N(e^r)^2$ on $\mathcal G_\delta(Y)$, with the fixed
choice $A_*=2$, absorbs the factor $\sqrt L$ and gives an
arbitrary logarithmic saving in mean square. The same
exceptional set argument proves \cref{eq:samelabelmean}.

\subsection{The product range on the factored moduli}

\begin{theorem}[The product range with arbitrary inputs]
\label{thm:productfactored}
Assume global Dirichlet GRH and let $\mathcal F(Y)$ be the family
of positive lower density in \cref{thm:criticalprofiles}.
For every fixed $B\ge0$ and $D>0$,
\begin{equation}
 \sum_{N\in\mathcal F(Y)} E_N^{\le}(B)^2
     \ll_{B,D} Y(\log Y)^{-D}.
 \label{eq:productfactored}
\end{equation}
\end{theorem}

Consequently, for each fixed $B\ge0$ and $D>0$, the normalized
full operator norm differs by at most $(\log Y)^{-D}$ from the
normalized norm of the form retaining $xy>R_N^*$, except
at $O_{B,D}(Y(\log Y)^{-D})$ moduli in $\mathcal F(Y)$.
This follows from the triangle inequality and
\cref{eq:productfactored} with exponent $3D$.

Combining this estimate with \cref{thm:upperproduct} gives
a two-sided product restriction on the same moduli.
Take $H=L(2\log L+\kappa)$, $\kappa\ge1$, and multiply the
form by
\[
 \mathbf1_{xy>R_N^*}
       \left(1-\Theta\!\left(\frac{\log(xy)-H}{L}\right)\right).
\]
Its normalized norm differs from $\|T_N\|/\SSS(N)$ by
$O_{B,D,\Theta}((\log Y)^{-D}+e^{-\kappa}+N^{-c})$ outside
the same exceptional set. The lower and upper deleted terms
are disjoint for large $Y$, so their operator errors add.
For example $\kappa=\log\log L$ gives vanishing error with
a retained upper product $N^{2\log\log N+o(\log\log N)}$.

Use the exponent $A$ of \cref{thm:productwindow} and split
each test at $M=(\log N)^A$ by a smooth cutoff in log
coordinates, with the lower piece supported in $x\le2M$
and the upper one in $x\ge M$. The multiplier norms are
bounded. The term with two upper pieces is controlled by
\cref{eq:productwindow}. By symmetry it suffices to bound
the terms with the first support in $x\le2M$.

For the low-input model, let $I_\pm$ be the indicator of
$x(bN\pm x)\le R_N^*$. Write the signed larger weight as
\[
 I_+\bigl(z(bN+x)-z(bN-x)\bigr)+(I_+-I_-)z(bN-x).
\]
The first term has normalized norm $O(M/N)$ by the
derivative argument for $\min(x,y)\le2M$ in
\cref{sec:subcritical}. For each fixed $x$ the indicators
can differ at most once, since then
$|bN-R_N^*/x|\le x\le2M$. Also
$\SSS(bN)/\SSS(N)\ll (bN)^{1/4}$ by the product in
\cref{eq:series}, since each factor at a sufficiently large
prime is at most $p^{1/4}$. Sobolev point evaluation at
$y=bN-x$ bounds $|z(y)|$ by $O(y^{-1/2}\NN(G))$.
Thus this extra term costs at most
\[
 O\!\left((R_N^*)^{-1/4}\NN(G)
             \int_0^{2M}|w(x)|x^{1/4}\,\ud x\right)
     \ll (R_N^*)^{-1/4}M^{3/4}\NN(F)\NN(G)
\]
on the normalized scale. The integral bound follows by
weighted Cauchy--Schwarz in log coordinates as in
\cref{sec:subcritical}. This is a power saving.

In the low-input arithmetic form, every contributing
larger integer satisfies $y\ge N-2M$. Use the auxiliary
smooth weight from the proof of \cref{thm:productwindow},
with the same $\epsilon=(\log N)^{-J}$. In a fixed row $n$,
the two weights differ only for
$R_N^*/n<y\le e^\epsilon R_N^*/n$. In either residue class
there are $O(1+\epsilon R_N^*/(nN))$ integers in this
interval. Each contributes at most
$O((\log N)^2(R_N^*)^{-1/2})$ on unit tests, by point
evaluation. Summing over $2\le n\le2M$ gives an error
\[
 O\!\left((\log N)^2\left(
          \frac{M}{\sqrt{R_N^*}}
           +\frac{\epsilon\sqrt{R_N^*}}N\log(2M)\right)\right).
\]
This is $O((\log N)^{-D/2})$ with the chosen $J$.
For each integer $n$ in the first input, write
\[
 G_n(r)=G(r)\Phi\bigl(\epsilon^{-1}(\log n+r-\log R_N^*)\bigr).
\]
The signed progression term has test $G_n$, whose norm is
$O(\epsilon^{-1}\NN(G))$. By odd-character orthogonality and
\cref{eq:ppchar}, the part of this row from proper prime
powers $y\ge N-2M$ costs $O(N^{-1/5}\epsilon^{-1}\NN(G))$.
Summing the first weights costs $O(\sqrt M\,\NN(F))$ as in
the proof of \cref{thm:criticalprofiles}, so this deletion
has a power saving.

Write $\lambda=\log Y$. Use a smooth lower cutoff at
$T_-=Y^2\lambda^{-H}$ with transition on $[T_-,2T_-]$,
where $H$ is a sufficiently large fixed exponent.
The discarded part of the arithmetic form has
$xy\le4M T_-$, so the absolute arithmetic estimate of
\cref{sec:subcritical} bounds its normalized norm by
$O(\lambda^{(A-H)/2}+Y^{-2/5})$.
Every remaining row test is supported below
$T_+=(2Y)^2(\log(2Y))^{B+1}$ for sufficiently large $Y$,
since $n\ge2$ and the mask vanishes when
$ny\ge e^\epsilon R_N^*$. Apply the partial summation bound
of \cref{sec:criticalprofiles} to each row with this
modified test, using the multiplier bound uniformly in $n$.
With $A_*=(\log(2Y))^A$, the norm of the remaining form
is at most a constant times
\[
 \lambda^{1+J}\sqrt{A_*}
    \left(\int_{\log T_-}^{\log T_+}d_N(e^r)^2\,\ud r\right)^{1/2}.
\]
The estimate \eqref{eq:profiledistribution} on $\mathcal F(Y)$
holds for this $A_*$ and these endpoints, using the same
factor ranges in Maynard's theorem. Its arbitrary log saving
absorbs the factors in the displayed bound after squaring and
summing over $\mathcal F(Y)$. Taking $H>A+D+2$ makes the
lower cutoff negligible in the claimed mean square. This
proves \cref{eq:productfactored}.

\clearpage
\section{Conditional bounds from uniform prime-pair discrepancies}
\label{sec:pairconditional}

Uniform quantitative prime-pair estimates would give norm
decay at each modulus through a much larger support range.
The distinction between the shorter and longer variable in
such estimates matters for separated supports. Difference-line
estimates with square-root errors are also used in conditional
results on the smooth Hooley variance \cite[Theorem~4]{TZLiu}.
The signed form here uses both difference and reflection lines.

\begin{theorem}[A conditional support range]
\label{thm:pairconditional}
Fix $0\le\theta\le1/2$. Suppose that for every $\epsilon>0$,
uniformly for $2\le X\le V$, intervals
$I\subset[X/e,eX]$, $J\subset[V/e,eV]$, signs $\sigma=\pm1$,
and positive even $h\le2eV$, the discrepancy of
\cref{sec:almostall} satisfies
\begin{equation}
 |\Delta_\sigma(h;I,J)|
       \ll_\epsilon X^\theta V^{1/2-\theta+\epsilon}.
 \label{eq:uniformhyp}
\end{equation}
Then for each fixed
\[
       1\le\beta<\frac{2-2\theta}{3/2-2\theta}
\]
there is $c>0$ such that $E_N(N^\beta)\ll_{\beta,\theta} N^{-c}$
for every sufficiently large even $N$. Under $\grh(N)$ this
is a bound on the corresponding operator restriction.
The larger-scale square-root error ($\theta=0$) gives the
range $\beta<4/3$, and the shorter-scale square-root error
($\theta=1/2$) gives $\beta<2$.
\end{theorem}

Choose $\rho>0$ with $2-\beta-\rho>0$. Use a smooth lower
cutoff from $N^{2-\beta-\rho}/2$ to $N^{2-\beta-\rho}$.
Every pair discarded by it has $xy\le N^{2-\rho}$.
\Cref{thm:subcritical} bounds the normalized error by
$O(N^{-\rho/2}+N^{-2/5})$. Its physical-form estimate
applies without GRH. Localize the retained tests to unit
log packets at scales $X\le V$ as before. Then
$X\gg N^{2-\beta-\rho}$ and $V\ll N^\beta$.

On a single line $m+\sigma n=h$, the physical product
weight $w(n)\overline{z(h-\sigma n)}$ has supremum and total
variation $O((XV)^{-1/2})$ on unit tests. This follows by
bounding each physical weight and the $L^1$ norm of its
derivative by the reciprocal square root of its scale,
as in the proof of \cref{eq:Mpacket}. Partial summation
with \cref{eq:uniformhyp} therefore bounds the weighted
discrepancy on that line by
$O_\epsilon(X^{\theta-1/2}V^{-\theta+\epsilon})$.
The pairs removed by $(mn,N)=1$ have total von Mangoldt
weight $O((\log N)^4)$ per line before the physical weights
are inserted. This error is absorbed in \cref{eq:uniformhyp}
in the retained range because $X,V$ grow as positive powers
of $N$. The equal-integer sum pairs are included in that
removal. Lines with $h=0$ are excluded by construction
from the difference terms.

Nonzero packet pairs have $V\gg N$ and use $O(V/N)$ shifts
$h=bN$. Thus the normalized packet norm is
$O_\epsilon(N^{-1}X^{\theta-1/2}V^{1-\theta+\epsilon})$.
The bounds on $X,V$ give the exponent
\[
 -2+2\theta+(3/2-2\theta)\beta
            +(1/2-\theta)\rho+\beta\epsilon
\]
in $N$. It is strictly negative for small enough $\rho$
and $\epsilon$ in the stated range. The $O((\log N)^2)$
packet pairs and the lower-cutoff error preserve a power
saving. This proves \cref{thm:pairconditional}.

\clearpage
\section{Scope of the asymptotic problem}

The equivalence in \cref{thm:disc} identifies an arithmetic problem
already present on the balanced window. Proving that restriction
small uniformly in $N$ entails a Hardy--Littlewood asymptotic in
every fixed interior subinterval. Conversely, these subinterval
asymptotics suffice for the restriction, using only an upper bound
for the total mass near its ends. The separated-window certificates
exhibit parts of the full form outside the reach of this equivalence.

An independent preprint by Y. Liu~\cite{YLiu} reports an asymptotic
obstruction to hypothesis R under Dirichlet GRH from fluctuations
of two opposite shifts. Its lower bound is on the scale
$N^{-1/2}(\log N)^D$ for fixed $D$. The smallness condition
\cref{eq:goal} concerns the larger singular-series scale. The
finite-modulus witness in \cref{sec:witness} likewise addresses
fixed thresholds at one modulus. By \cref{thm:mesh}, proving
\cref{eq:goal} amounts to bounding the generalized eigenvalues
of a matrix of weighted correlation sums over the retained
interval. \Cref{thm:tracecompression} reduces the number of
spectral directions in that check, and \eqref{eq:tracemoment}
gives an equivalent test in terms of growing even trace moments.
\Cref{thm:signedtrace} fixes the leading constant in a
necessary mean-square law for the Dirichlet logarithmic
derivatives at distance $1/\log N$ from the critical line.
\Cref{thm:localderivative,thm:necessarypair} strengthen this
to bounded height intervals at every fixed multiple of that
distance and to the corresponding microscopic pair law.
\Cref{thm:restrictedzeropairs} establishes that pair law
under $\grh(N)$ for Fourier support in $(-1,1)$ and gives
mixed covariance when the two Fourier widths sum to at
most $2$. These statistics concern fixed smooth filters
and the leading zero covariance. The full-norm test still
requires uniform control at additive $o(\SSS(N))$ precision.
\Cref{thm:fullzerodata} removes long prime-power sums from
the finite matrix criterion. Its arithmetic entries retain
the two separate test transforms on the ordinates and the
Euler terms at primes dividing the modulus.
The pair deletion in \cref{thm:subcritical} places
the remaining arithmetic at and above products of order
$N^2/Q$. This includes both the separated support scales in
\cref{sec:rect} and the balanced count in \cref{sec:balanced}.
For each fixed upper cap $CN^2$, \cref{thm:criticaldiscrepancy}
identifies the norm problem in this range with a signed rectangle
discrepancy. In the most separated boxes the rectangle theorem
keeps the discrete smaller coordinate and uses cancellation of the
two continuous line contributions there.
\Cref{thm:criticalgrid} makes this an explicit test at fixed precision.
The number of signed arithmetic counts grows at most logarithmically
in $N$, and the number of entries that need the continuous model
is bounded independently of $N$.
By \cref{thm:ppcut}, the arithmetic on $x,y\ge A$ can be
restricted to primes, with the smaller input governed by
the progression profiles in \cref{eq:pprows}.
For those profiles, \cref{thm:criticalprofiles} controls the
part through $N^2(\log N)^B$ on the stated modulus families.
For the high-support block, \cref{thm:zeros} gives a fixed-modulus
covariance criterion with an additive precision requirement.
Known variance obstructions of the type studied by Fiorilli and
Martin and by Hayani occur at $N\asymp\log\log x$
\cite{FiorilliMartin,Hayani}. These support scales lie beyond
the upper cutoff of \cref{thm:cutoff}.
The estimates of \cref{thm:almostall,thm:mixedpower} give decay
on growing supported windows, including separated scales at
the critical product. \Cref{thm:powerproduct} gives a power range
above that product when both inputs are power-sized.
\Cref{thm:subcritical,thm:asymupperproduct} give a two-sided
product range at every modulus under $\grh(N)$.
On the factored family, \cref{thm:productfactored} also gives
mean-square decay through the critical region for arbitrary inputs.
A uniform estimate through the retained range, including the
critical region for general moduli, remains open.
The conditional range in \cref{thm:pairconditional} specifies
the strength of individual line estimates needed for uniform
decay below the square of the modulus.

\section*{Reproducibility}

The accompanying scripts were run with Python 3.11.7, NumPy 1.26.4,
SciPy 1.17.1 and mpmath 1.3.0. Run these commands from the directory
containing the scripts and the rational witness data.

{\small
\begin{verbatim}
python hypothesis_s_balanced_certificate.py 32762 1048576 4194286 4194304 16777216
python hypothesis_s_critical_certificate.py 4194286
python hypothesis_s_critical_certificate.py 4194304
python hypothesis_s_critical_certificate.py 16777216
python hypothesis_s_witness.py --weight-file hypothesis_s_witness_32762.json
\end{verbatim}
}

The derivations, computational code and manuscript were developed
with OpenAI Codex assistance.

The file \texttt{formalization/CriticalCollar.lean} verifies
the discriminant identity, vertex exclusion and separation
inequalities for the line geometry in section 9.3.
The sieve estimates and the operator comparison are proved
analytically in that section.
The file \texttt{formalization/CriticalPrefix.lean} verifies
the four-prefix identity and the dyadic index counts in
\cref{sec:criticalgrid}.

\clearpage
{\small
\begin{thebibliography}{99}
\bibitem{LiLiu}
J.~Li and J.~Liu,
\newblock \emph{A conditional proof of the Goldbach conjecture},
\newblock preprint, Mathematical Research Center, Shandong University,
13 November 2026.
\newblock \url{https://mrc.sdu.edu.cn/english/info/1018/1514.htm}.

\bibitem{Lichtman}
J.~D. Lichtman,
\newblock \emph{Primes in arithmetic progressions to large moduli,
and Goldbach beyond the square-root barrier},
\newblock arXiv:2309.08522, 2023.
\newblock \url{https://arxiv.org/abs/2309.08522}.

\bibitem{YLiu}
Y.~Liu,
\newblock \emph{Failure of the Li--Liu spectral hypothesis R under Dirichlet GRH},
\newblock Hexagon preprint 2611.00788v1, 21 November 2026.
\newblock \url{https://hexagonmath.org/2611.00788v1}.

\bibitem{BMOR}
M.~A. Bennett, G. Martin, K. O'Bryant and A. Rechnitzer,
\newblock \emph{Counting zeros of Dirichlet $L$-functions},
\newblock Math. Comp. \textbf{90} (2021), 1455--1482.
\newblock \url{https://arxiv.org/abs/2005.02989}.

\bibitem{MaiQian}
W. Mai and T. Qian,
\newblock \emph{Rational approximation in Hardy spaces on strips},
\newblock Complex Var. Elliptic Equ. \textbf{63} (2018), 1721--1738.
\newblock \url{https://www.fst.um.edu.mo/en/staff/documents/fsttq/MQ5.pdf}.

\bibitem{Helfgott}
H.~A. Helfgott,
\newblock \emph{The ternary Goldbach problem},
\newblock revised monograph draft, 2019, Part~IV.
\newblock \url{https://webusers.imj-prg.fr/~harald.helfgott/anglais/PartIV.pdf}.

\bibitem{Richert}
H.-E. Richert,
\newblock \emph{Lectures on sieve methods},
\newblock Tata Institute of Fundamental Research, Bombay, 1976.
\newblock \url{https://mathweb.tifr.res.in/sites/default/files/publications/ln/tifr55.pdf}.

\bibitem{TZLiu}
T.~Z. Liu,
\newblock \emph{Pair correlation of $L$-zeros and Hooley's conjecture},
\newblock Hexagon preprint 2611.00509v1, 16 November 2026.
\newblock \url{https://hexagonmath.org/2611.00509v1}.

\bibitem{DLMF}
NIST Digital Library of Mathematical Functions,
\newblock \emph{Chapter 5. Gamma function}.
\newblock \url{https://dlmf.nist.gov/5}.

\bibitem{HughesRudnick}
C.~P. Hughes and Z. Rudnick,
\newblock \emph{Linear statistics of low-lying zeros of $L$-functions},
\newblock Q. J. Math. \textbf{54} (2003), 309--333.
\newblock \url{https://arxiv.org/abs/math/0208230}.

\bibitem{CarneiroEtAl}
E. Carneiro, V. Chandee, F. Littmann and M.~B. Milinovich,
\newblock \emph{Hilbert spaces and the pair correlation of zeros of
the Riemann zeta-function},
\newblock J. Reine Angew. Math. \textbf{725} (2017), 143--182.
\newblock \url{https://arxiv.org/abs/1406.5462}.

\bibitem{HuaYang}
Z. Hua and X. Yang,
\newblock \emph{Simple and distinct zeros in a prime-modulus Dirichlet
family from near-microscopic to polylogarithmic heights},
\newblock arXiv:2608.16034v2, 2026.
\newblock \url{https://arxiv.org/abs/2608.16034v2}.

\bibitem{FiorilliMartin}
D. Fiorilli and G. Martin,
\newblock \emph{Disproving Hooley's conjecture},
\newblock J. Eur. Math. Soc. \textbf{25} (2023), 4791--4812.
\newblock \url{https://arxiv.org/abs/2008.05837}.

\bibitem{Hayani}
M. Hayani,
\newblock \emph{Disproving a weaker form of Hooley's conjecture},
\newblock Bull. Sci. Math. \textbf{199} (2025), 103582.
\newblock \url{https://arxiv.org/abs/2407.01045}.

\bibitem{VaughanExp}
R.~C. Vaughan,
\newblock \emph{Sommes trigonom\'etriques sur les nombres premiers},
\newblock C. R. Acad. Sci. Paris S\'er. A-B \textbf{285} (1977), A981--A983.

\bibitem{Walfisz}
A. Walfisz,
\newblock \emph{Zur additiven Zahlentheorie. II},
\newblock Math. Z. \textbf{40} (1936), 592--607.
\newblock \url{https://doi.org/10.1007/BF01218882}.

\bibitem{LeeGRH}
E.~S. Lee,
\newblock \emph{The prime number theorem for primes in arithmetic
progressions at large values},
\newblock Q. J. Math. \textbf{74} (2023), 1505--1533.
\newblock \url{https://arxiv.org/abs/2301.13457}.

\bibitem{KLM}
N. Kandhil, A. Languasco, and P. Moree,
\newblock \emph{Pair correlation of zeros of Dirichlet $L$-functions:
a possible path towards the conjectures of Chowla, Elliott--Halberstam
and Montgomery},
\newblock Math. Ann. \textbf{394} (2026), article~43.
\newblock \url{https://doi.org/10.1007/s00208-026-03383-y}.

\bibitem{MaynardI}
J. Maynard,
\newblock \emph{Primes in arithmetic progressions to large moduli I:
Fixed residue classes},
\newblock Mem. Amer. Math. Soc. \textbf{306} (2025), no.~1542.
\newblock \url{https://arxiv.org/abs/2006.06572}.

\bibitem{MaynardIII}
J. Maynard,
\newblock \emph{Primes in arithmetic progressions to large moduli III:
Uniform residue classes},
\newblock Mem. Amer. Math. Soc. \textbf{306} (2025), no.~1544.
\newblock \url{https://arxiv.org/abs/2006.08250}.
\end{thebibliography}}

\end{document}
